% Lesson 9 — Vocabulary: Words and expressions related to recreational activities — Anglais Tle A — Cours signé OnBuch+
% Programme officiel MINESEC. Compiler avec : tectonic EN09-vocabulary-words-and-expressions-related-to-recreationa.tex
\newcommand{\DOCMATIERE}{Anglais}
\newcommand{\DOCNIVEAU}{Tle A}
\newcommand{\DOCMODULE}{Chapter 1 : Family and Social Life}
\newcommand{\DOCLECON}{EN09}
\newcommand{\DOCTITRE}{Lesson 9 — Vocabulary: Words and expressions related to recreational activities}
% OnBuch+ — préambule « cours signés » ENRICHI (compilé avec tectonic / XeLaTeX)
% Système visuel « Pop » d'OnBuch+ : fond crème (jamais blanc), bordures encre
% épaisses, ombres « dures » décalées, titres Archivo Black en capitales.
% Version MATHÉMATIQUES : graphes (pgfplots/tikz), boîtes pédagogiques
% (définition, propriété, méthode, attention, le savais-tu, exercice résolu,
% exercice, TP, bilan), page de garde et signature OnBuch+ ancrée en bas.
%
% Macros attendues dans main.tex AVANT \input{preamble} :
%   \DOCMATIERE  \DOCNIVEAU  \DOCMODULE  \DOCLECON (numéro)  \DOCTITRE
\documentclass[11pt]{article}

\usepackage{fontspec}
\usepackage[english,french]{babel} % français = langue principale ; l'anglais via \eng{..} / textebox
\usepackage[a4paper,margin=2cm,top=3.4cm,bottom=2.8cm,headheight=1.4cm,footskip=1.3cm]{geometry}
\usepackage{amsmath,amssymb,amsfonts}
\usepackage{mathtools} % \xmapsto, \coloneqq... souvent utilisés par les modèles
\usepackage{cancel} % simplifications barrées (bilans de forces)
% \abs{z} (module, valeur absolue) et \norm{u} (norme), utilisés par les modèles
\providecommand{\abs}{}\let\abs\relax\DeclarePairedDelimiter{\abs}{\lvert}{\rvert}
\providecommand{\norm}{}\let\norm\relax\DeclarePairedDelimiter{\norm}{\lVert}{\rVert}
\usepackage[table]{xcolor} % [table] : fournit aussi \rowcolors (alternance de lignes)
\usepackage{pagecolor}
\usepackage{tikz}
\usetikzlibrary{arrows.meta,positioning,calc,shapes.geometric,shapes.misc,shapes.arrows,decorations.pathmorphing,decorations.pathreplacing,decorations.markings,patterns,shadows,mindmap,trees,fit,backgrounds,matrix,intersections,angles,quotes}
\usepackage{pgfplots}
\usepackage[version=4]{mhchem} % équations nucléaires \ce{^{235}_{92}U}
\usepackage[european,siunitx]{circuitikz} % schémas électriques
\pgfplotsset{compat=1.18}
\usepgfplotslibrary{fillbetween,groupplots} % aires entre courbes (calcul intégral)
\usepackage{enumitem}
\usepackage{fancyhdr}
% [most] charge toutes les bibliothèques tcolorbox (listings, minted, raster,
% documentation...) et gonfle inutilement la mémoire TeX sur un document long
% (~300 boîtes) : on ne charge que ce qui est réellement utilisé.
\usepackage[skins,breakable]{tcolorbox}
\usepackage{graphicx}
\usepackage{colortbl}
\usepackage{booktabs}
\usepackage{tabularx}
\usepackage{diagbox} % case d'en-tête diagonale des tableaux à double entrée
% Colonnes extensibles centrée / gauche / droite (C, L, R), souvent utilisées par les modèles
\newcolumntype{C}{>{\centering\arraybackslash}X}
\newcolumntype{L}{>{\raggedright\arraybackslash}X}
\newcolumntype{R}{>{\raggedleft\arraybackslash}X}
\usepackage{array}
\usepackage{multicol}
\usepackage{multirow}
\usepackage{siunitx}
% parse-numbers=false : accepte aussi \SI{2,5 \times 10^{-3}}{...}
\sisetup{locale=FR,output-decimal-marker={,},group-minimum-digits=4,parse-numbers=false}
\DeclareSIUnit\ohm{\ensuremath{\symup{\Omega}}} % U+2126 absent de la police
\DeclareSIUnit\division{div} % réglages d'oscilloscope (ms/div, V/div)
\DeclareSIUnit\tr{tr} % tours (vitesse de rotation, ex. \SI{1500}{\tr\per\minute})
\DeclareSIUnit\molar{M} % concentration molaire (ex. \SI{3}{\molar}, \SI{1}{\micro\molar})
\DeclareSIUnit\an{an} % année (le modèle confond parfois avec \year, un compteur TeX) — ex. \SI{5}{\kilo\meter\per\an}
\DeclareSIUnit\FCFA{FCFA} % franc CFA (ex. \SI{500}{\FCFA\per\kilo\gram})
\DeclareSIUnit\degreeBrix{\textdegree Brix} % teneur en sucre d'un sirop/jus (réfractométrie)
\DeclareSIUnit\month{mois} % le modèle utilise parfois \month, un compteur TeX, comme unité de durée
\usepackage{qrcode}
% \degree : commande fréquemment utilisée par les modèles pour un angle ou une
% température en dehors de \SI{}{} (non fournie par les paquets chargés).
\providecommand{\degree}{^{\circ}}
% \qed (fin de démonstration), utilisé par les modèles sans amsthm
\providecommand{\qed}{\hfill$\square$}
% \barre : double barre des valeurs interdites dans un tableau de variations (tkz-tab)
\providecommand{\barre}{\ensuremath{\|}}
% environnement proof (amsthm non chargé) utilisé par les modèles
\ifdefined\proof\else\newenvironment{proof}[1][Démonstration]{\par\noindent\textit{#1.}\ }{\qed\par}\fi
\usepackage{lastpage}
\usepackage{hyperref}
\hypersetup{hidelinks}

% ============================================================
% Palette « Pop » OnBuch+ — mêmes hex que la classe P (plus_ui.dart)
% ============================================================
\definecolor{poporange}{HTML}{F59321}
\definecolor{popdark}{HTML}{E07A0C}
\definecolor{popcream}{HTML}{FFF6E8}
\definecolor{popink}{HTML}{1C1714}
\definecolor{poppurple}{HTML}{7A5AE0}
\definecolor{popgreen}{HTML}{2E9E6A}
\definecolor{poppink}{HTML}{E0457B}
\definecolor{popblue}{HTML}{2F7DE1}
\definecolor{popgold}{HTML}{C98A12}
\definecolor{popmuted}{HTML}{8A7F76}
\definecolor{popline}{HTML}{EADFD2}
\definecolor{popgray}{HTML}{8A7F76}
\colorlet{popgrey}{popgray}
% Alias tolérants (le modèle nomme parfois les couleurs différemment)
\colorlet{popwhite}{white}
\colorlet{popblack}{popink}
\colorlet{popred}{poppink}
\colorlet{popyellow}{popgold}
% Teintes claires (fonds de tableaux/encadrés) — le modèle nomme parfois
% les couleurs avec un suffixe L (ex. popblueL) sans les avoir définies.
\colorlet{poporangeL}{poporange!20}
\colorlet{popdarkL}{popdark!20}
\colorlet{poppurpleL}{poppurple!20}
\colorlet{popgreenL}{popgreen!20}
\colorlet{poppinkL}{poppink!20}
\colorlet{popblueL}{popblue!20}
\colorlet{popgoldL}{popgold!20}
\colorlet{popredL}{popred!20}
\colorlet{popgrayL}{popgray!20}
% Teintes claires pour fonds de tableaux / zones
\colorlet{popblueL}{popblue!10!white}
\colorlet{poporangeL}{poporange!14!white}
\colorlet{popgreenL}{popgreen!12!white}
\colorlet{poppurpleL}{poppurple!12!white}
\colorlet{poppinkL}{poppink!12!white}
\colorlet{popgoldL}{popgold!14!white}

\pagecolor{popcream}

% ============================================================
% Typographie
% ============================================================
\defaultfontfeatures{Ligatures=TeX}
\setmainfont{PlusJakartaSans-Medium.ttf}[Path=../../../fonts/,BoldFont=PlusJakartaSans-Bold.ttf]
% Symboles, API et emojis absents de la police principale : repli sur DejaVu Sans (caractères actifs)
\newfontface\symfont{DejaVuSans.ttf}[Path=../../../fonts/]
\makeatletter
\newcommand\symactive[1]{\catcode#1=13 \begingroup\lccode`\~=#1 \lowercase{\endgroup\def~}{{\symfont\char#1}}}
\newcommand\symblank[1]{\catcode#1=13 \begingroup\lccode`\~=#1 \lowercase{\endgroup\def~}{}}
\newcommand\symhyph[1]{\catcode#1=13 \begingroup\lccode`\~=#1 \lowercase{\endgroup\def~}{\mbox{-}}}
\newcount\symc
\newcommand\symrange[2]{\symc=#1 \@whilenum\symc<#2 \do{\expandafter\symactive\expandafter{\the\symc}\advance\symc by 1 }}
\newcommand\symblankrange[2]{\symc=#1 \@whilenum\symc<#2 \do{\expandafter\symblank\expandafter{\the\symc}\advance\symc by 1 }}
\makeatother
\symrange{"0250}{"02FF}\symactive{"2713}\symactive{"2714}\symactive{"2717}\symactive{"2718}\symactive{"2610}\symactive{"2611}\symactive{"2612}
\symhyph{"2011}\symhyph{"2E1A}\symblankrange{"1F300}{"1FAFF}\symblank{"FE0F}
% Maths en sans-serif (Fira Math) avec chiffres et lettres droites de Plus
% Jakarta, pour que les formules se fondent dans le texte.
\usepackage[math-style=ISO,bold-style=ISO]{unicode-math}
\setmathfont{FiraMath-Regular.otf}[Path=../../../fonts/]
\setmathfont{PlusJakartaSans-Medium.ttf}[Path=../../../fonts/,range={up/num,up/{Latin,latin},\mathup}]
\usepackage{icomma} % virgule décimale sans espace en mode math
% Caractères mathématiques Unicode absents des polices de texte (Plus Jakarta,
% Archivo Black) : rendus par leur équivalent mathématique, en texte comme en
% formule — sinon ils s'affichent en carré vide (ex. « Arithmétique dans ℤ »).
\usepackage{newunicodechar}
\newunicodechar{ℝ}{\ensuremath{\mathbb{R}}}
\newunicodechar{ℤ}{\ensuremath{\mathbb{Z}}}
\newunicodechar{ℂ}{\ensuremath{\mathbb{C}}}
\newunicodechar{ℕ}{\ensuremath{\mathbb{N}}}
\newunicodechar{ℚ}{\ensuremath{\mathbb{Q}}}
\newunicodechar{𝔻}{\ensuremath{\mathbb{D}}}
\newunicodechar{α}{\ensuremath{\alpha}}
\newunicodechar{β}{\ensuremath{\beta}}
\newunicodechar{γ}{\ensuremath{\gamma}}
\newunicodechar{δ}{\ensuremath{\delta}}
\newunicodechar{ε}{\ensuremath{\varepsilon}}
\newunicodechar{θ}{\ensuremath{\theta}}
\newunicodechar{λ}{\ensuremath{\lambda}}
\newunicodechar{μ}{\ensuremath{\mu}}
\newunicodechar{σ}{\ensuremath{\sigma}}
\newunicodechar{τ}{\ensuremath{\tau}}
\newunicodechar{φ}{\ensuremath{\varphi}}
\newunicodechar{ψ}{\ensuremath{\psi}}
\newunicodechar{ω}{\ensuremath{\omega}}
\newunicodechar{Σ}{\ensuremath{\Sigma}}
% majuscules produites par \MakeUppercase dans les titres (α-aminés -> Α)
\newunicodechar{Α}{\ensuremath{\alpha}}
\newunicodechar{Β}{\ensuremath{\beta}}
\newunicodechar{Γ}{\ensuremath{\Gamma}}
\newunicodechar{Δ}{\ensuremath{\Delta}}
\newunicodechar{Ε}{\ensuremath{\varepsilon}}
\newunicodechar{Θ}{\ensuremath{\Theta}}
\newunicodechar{Λ}{\ensuremath{\Lambda}}
\newunicodechar{Μ}{\ensuremath{\mu}}
\newunicodechar{Τ}{\ensuremath{\tau}}
\newunicodechar{Φ}{\ensuremath{\Phi}}
\newunicodechar{Ψ}{\ensuremath{\Psi}}
\newunicodechar{Ω}{\ensuremath{\Omega}}
\newunicodechar{↦}{\ensuremath{\mapsto}}
\newunicodechar{∈}{\ensuremath{\in}}
\newunicodechar{∉}{\ensuremath{\notin}}
\newunicodechar{∧}{\ensuremath{\wedge}}
\newunicodechar{⇔}{\ensuremath{\Leftrightarrow}}
\newunicodechar{⟺}{\ensuremath{\Longleftrightarrow}}
\newunicodechar{⇒}{\ensuremath{\Rightarrow}}
\newunicodechar{⟹}{\ensuremath{\Longrightarrow}}
\newunicodechar{∘}{\ensuremath{\circ}}
\newunicodechar{∪}{\ensuremath{\cup}}
\newunicodechar{∩}{\ensuremath{\cap}}
\newunicodechar{≡}{\ensuremath{\equiv}}
\newunicodechar{⊂}{\ensuremath{\subset}}
\newunicodechar{⊥}{\ensuremath{\perp}}
\newunicodechar{∃}{\ensuremath{\exists}}
\newunicodechar{∀}{\ensuremath{\forall}}
\newunicodechar{∛}{\ensuremath{\sqrt[3]{\,}}}
\newunicodechar{∜}{\ensuremath{\sqrt[4]{\,}}}
\newunicodechar{⃗}{\ensuremath{{}^{\to}}}
\newunicodechar{ⁿ}{\ifmmode^{n}\else\textsuperscript{\ensuremath{n}}\fi}
\newunicodechar{ˣ}{\ifmmode^{x}\else\textsuperscript{\ensuremath{x}}\fi}
\newunicodechar{ᵇ}{\ifmmode^{b}\else\textsuperscript{\ensuremath{b}}\fi}
\newunicodechar{ʳ}{\ifmmode^{r}\else\textsuperscript{\ensuremath{r}}\fi}
\newunicodechar{ᵘ}{\ifmmode^{u}\else\textsuperscript{\ensuremath{u}}\fi}
\newunicodechar{ʸ}{\ifmmode^{y}\else\textsuperscript{\ensuremath{y}}\fi}
\newunicodechar{ᵃ}{\ifmmode^{a}\else\textsuperscript{\ensuremath{a}}\fi}
\newunicodechar{ᵉ}{\ifmmode^{e}\else\textsuperscript{\ensuremath{e}}\fi}
\newunicodechar{ᵏ}{\ifmmode^{k}\else\textsuperscript{\ensuremath{k}}\fi}
\newunicodechar{ᵖ}{\ifmmode^{p}\else\textsuperscript{\ensuremath{p}}\fi}
\newunicodechar{ᴺ}{\ifmmode^{N}\else\textsuperscript{\ensuremath{N}}\fi}
\newunicodechar{ᶜ}{\ifmmode^{c}\else\textsuperscript{\ensuremath{c}}\fi}
\newunicodechar{ʰ}{\ifmmode^{h}\else\textsuperscript{\ensuremath{h}}\fi}
\newunicodechar{ᵠ}{\ifmmode^{q}\else\textsuperscript{\ensuremath{q}}\fi}
\newunicodechar{ₙ}{\ifmmode_{n}\else\textsubscript{\ensuremath{n}}\fi}
\newunicodechar{ₐ}{\ifmmode_{a}\else\textsubscript{\ensuremath{a}}\fi}
\newunicodechar{ᵢ}{\ifmmode_{i}\else\textsubscript{\ensuremath{i}}\fi}
\newunicodechar{ᵣ}{\ifmmode_{r}\else\textsubscript{\ensuremath{r}}\fi}
\newunicodechar{ₖ}{\ifmmode_{k}\else\textsubscript{\ensuremath{k}}\fi}
\newunicodechar{ᵦ}{\ifmmode_{\beta}\else\textsubscript{\ensuremath{\beta}}\fi}
\newunicodechar{ₓ}{\ifmmode_{x}\else\textsubscript{\ensuremath{x}}\fi}
\newfontfamily\popdisplay{ArchivoBlack-Regular.ttf}[Path=../../../fonts/]
\newfontfamily\poplogo{SpaceGrotesk-Bold.ttf}[Path=../../../fonts/]
\newfontfamily\popsemi{PlusJakartaSans-SemiBold.ttf}[Path=../../../fonts/]
\newfontfamily\popxbold{PlusJakartaSans-ExtraBold.ttf}[Path=../../../fonts/]
\newfontfamily\popxboldls{PlusJakartaSans-ExtraBold.ttf}[Path=../../../fonts/,LetterSpace=3.0]

\setlist[enumerate]{leftmargin=1.7em,itemsep=5pt,topsep=4pt}
\setlist[itemize]{leftmargin=1.7em,itemsep=4pt,topsep=4pt,label={\color{poporange}\textbullet}}
\setlength{\parindent}{0pt}
\setlength{\parskip}{5pt}
\renewcommand{\arraystretch}{1.25}

% pgfplots : style Pop par défaut
\pgfplotsset{
  every axis/.append style={axis line style={popink,line width=1pt},tick style={popink},
    grid style={popline,line width=0.5pt},label style={font=\small},tick label style={font=\footnotesize}},
  popcurve/.style={poporange,line width=1.8pt,no markers,smooth},
}
% Même style disponible dans un \draw TikZ simple (hors pgfplots)
\tikzset{popcurve/.style={poporange,line width=1.8pt}}
\tikzset{popgrid/.style={popline,very thin}}
\colorlet{popcurve}{poporange} % le nom du style est parfois utilisé comme couleur

% ============================================================
% Pill / squiggle
% ============================================================
\newcommand{\poppill}[3]{%
  \tikz[baseline=-0.55ex]\node[draw=popink,line width=1.2pt,rounded corners=2.6pt,%
    fill=#2,inner xsep=6.5pt,inner ysep=3pt]{%
    \popxboldls\fontsize{7}{7}\selectfont\color{#3}\MakeUppercase{#1}};%
}
\newcommand{\popsquiggle}[1]{%
  \tikz[baseline=-0.4ex]{
    \draw[#1,line width=2.6pt,line cap=round]
      plot [smooth] coordinates {(0,0.09) (0.34,0.01) (0.68,0.21) (1.02,0.01) (1.36,0.10)};
  }%
}
% Mot-clé surligné (marqueur orange sous le texte)
\newcommand{\cle}[1]{{\popxbold\color{popdark}#1}}

% ============================================================
% En-tête / pied
% ============================================================
\pagestyle{fancy}
\fancyhf{}
\renewcommand{\headrulewidth}{0pt}
\renewcommand{\footrulewidth}{0pt}
\lhead{\raisebox{-0.15ex}{\poplogo\fontsize{13}{13}\selectfont\color{popink}OnBuch\color{poporange}+}}
\rhead{\poppill{\DOCMATIERE\ \textbullet\ \DOCNIVEAU\ \textbullet\ Leçon \DOCLECON}{white}{popink}}
\lfoot{\popxbold\fontsize{7.6}{7.6}\selectfont\color{popmuted}\MakeUppercase{Cours signé OnBuch+}\ \textbullet\ {\popsemi\DOCTITRE}}
\rfoot{\tikz[baseline=-0.6ex]\node[rounded corners=8.5pt,fill=popink,minimum height=17pt,minimum width=17pt,inner xsep=5pt,inner sep=0]{\popxbold\fontsize{8}{8}\selectfont\color{popcream}\thepage\,/\,\pageref*{LastPage}};}

% ============================================================
% Titres
% ============================================================
\newcommand{\chaptitre}[2]{%
  \begin{tcolorbox}[enhanced,colback=poporange,colframe=popink,boxrule=2.6pt,arc=11pt,
      drop shadow=popink,left=15pt,right=15pt,top=12pt,bottom=11pt,before skip=3pt,after skip=18pt]
    \poppill{#2}{popink}{popcream}\par\vspace{9pt}
    {\raggedright\hyphenpenalty=10000\exhyphenpenalty=10000\popdisplay\fontsize{22}{24}\selectfont\color{popcream}\MakeUppercase{#1}\par}
  \end{tcolorbox}
}
% Section numérotée : pastille orange avec le numéro + titre Archivo
\newcounter{popsec}
\newcommand{\coursec}[1]{%
  \stepcounter{popsec}\setcounter{popsub}{0}%
  \par\vspace{16pt}\noindent
  \begin{minipage}{\linewidth}
  \tikz[baseline=-0.7ex]\node[draw=popink,line width=1.6pt,fill=poporange,rounded corners=4pt,minimum size=22pt,inner sep=0]{\popdisplay\fontsize{12}{12}\selectfont\color{popcream}\thepopsec};%
  \hspace{8pt}{\popdisplay\fontsize{15}{17}\selectfont\color{popink}\MakeUppercase{#1}}\par
  \vspace{4pt}\noindent{\color{poporange}\rule{36pt}{3.6pt}}
  \end{minipage}\par\nopagebreak\vspace{8pt}
}
\newcounter{popsub}
\newcommand{\courssub}[1]{%
  \stepcounter{popsub}%
  \par\vspace{9pt}\noindent{\popxbold\fontsize{11.5}{13}\selectfont\color{popink}{\color{poporange}\thepopsec.\thepopsub}\hspace{6pt}#1}\par\nopagebreak\vspace{4pt}
}

% ============================================================
% Boîtes pédagogiques — même recette : fond blanc, bordure épaisse colorée,
% ombre dure encre, badge plein. Titre optionnel [#1].
% ============================================================
\tcbset{popbase/.style={breakable,enhanced,colback=white,boxrule=2.3pt,arc=9pt,
  shadow={3pt}{-3pt}{0pt}{popink},left=14pt,right=14pt,top=11pt,bottom=11pt,before skip=11pt,after skip=15pt,
  fontupper=\popsemi\fontsize{10.6}{14.5}\selectfont\color{popink},
  fontlower=\popsemi\fontsize{10.6}{14.5}\selectfont\color{popink}}}
% Coche dessinée (le glyphe ✓ n'existe pas dans les polices du document)
\newcommand{\popcheck}[1]{\tikz[baseline=-0.5ex]\draw[#1,line width=1.6pt,line cap=round,line join=round](-0.13,0.0)--(-0.04,-0.09)--(0.13,0.1);}
\providecommand{\coche}{\popcheck{popgreen}}
\providecommand{\resultat}[1]{\par\smallskip\noindent\fcolorbox{popgreen}{popgreenL}{\parbox{\dimexpr\linewidth-2\fboxsep-2\fboxrule\relax}{\textbf{Résultat.} #1}}\par\smallskip}
\newcommand{\popboxhead}[3]{\poppill{#1}{#2}{white}\ifx\relax#3\relax\else\hspace{6pt}{\popxbold\fontsize{10.6}{12}\selectfont\color{popink}#3}\fi\par\vspace{7pt}}

\newtcolorbox{aretenir}[1][]{popbase,colframe=popgreen,before upper={\popboxhead{À retenir}{popgreen}{#1}}}
% Texte en anglais (document de lecture, dialogue, extrait) : cadre
% d'étude. Un titre optionnel (source, auteur) s'affiche dans l'en-tête.
\newtcolorbox{textebox}[1][]{popbase,colframe=popblue,before upper={\popboxhead{Text}{popblue}{#1}\begin{otherlanguage}{english}},after upper={\end{otherlanguage}}}
% Marques ✓ / ✗ (forme correcte / fautive) souvent utilisées par les modèles
\providecommand{\cross}{\textcolor{poppink}{\ensuremath{\times}}}
\providecommand{\xmark}{\cross}\providecommand{\crossmark}{\cross}\providecommand{\cmark}{\ensuremath{\checkmark}}
% Ligne de réponse / justification à compléter (inventée par les modèles)
\providecommand{\justification}[1]{\par\noindent\textit{Justification~:} \dotfill\par}
% \textipa (package tipa non chargé) : la police ne couvre pas l'API -> simple passage
\providecommand{\textipa}[1]{#1}
% Soulignement (ulem non chargé) : \uline, \ul
\providecommand{\uline}[1]{\underline{#1}}\providecommand{\ul}[1]{\underline{#1}}
% \glossaire{\item ...} inventée par les modèles : liste de glossaire
\providecommand{\glossaire}[1]{\begin{itemize}#1\end{itemize}}
% \alert (beamer) inventée par les modèles : texte mis en évidence
\providecommand{\alert}[1]{\textbf{\textcolor{poppink}{#1}}}
% variantes ASCII d'ordinaux écrites en macro
\providecommand{\iere}{\textsuperscript{re}}\providecommand{\ieme}{\textsuperscript{e}}\providecommand{\ere}{\textsuperscript{re}}\providecommand{\eme}{\textsuperscript{e}}\providecommand{\er}{\textsuperscript{er}}\providecommand{\ie}{\textsuperscript{e}}
% Ordinaux français écrits en macro par les modèles : 1\ière, 3\ième
\newcommand{\ière}{\textsuperscript{re}}\newcommand{\ième}{\textsuperscript{e}}
% \exemple (inventée par les modèles) : étiquette « Exemple : »
\providecommand{\exemple}{\textbf{Exemple~:}\ }
% \square utilisable aussi hors mode math (cases à cocher)
\AtBeginDocument{\let\squareMath\square\renewcommand{\square}{\ensuremath{\squareMath}}}
% Mot ou phrase en anglais à l'intérieur d'un paragraphe français
\newcommand{\eng}[1]{\ifmmode\text{\textit{#1}}\else\begingroup\selectlanguage{english}\itshape #1\endgroup\fi}
\let\esp\eng % alias toléré
% flèches d'intonation inventées par les modèles
\providecommand{\textsearrow}{}\renewcommand{\textsearrow}{\ensuremath{\searrow}}\providecommand{\textnearrow}{}\renewcommand{\textnearrow}{\ensuremath{\nearrow}}
% \fr{..} : mot français (inventé par les modèles)
\providecommand{\fr}[1]{\textit{#1}}
\newtcolorbox{exemplebox}[1][]{popbase,colframe=poporange,before upper={\popboxhead{Exemple}{poporange}{#1}}}
\newtcolorbox{definition}[1][]{popbase,colframe=popblue,before upper={\popboxhead{Définition}{popblue}{#1}}}
\newtcolorbox{propriete}[1][]{popbase,colframe=poppurple,before upper={\popboxhead{Propriété}{poppurple}{#1}}}
\newtcolorbox{methode}[1][]{popbase,colframe=popgold,before upper={\popboxhead{Méthode}{popgold}{#1}}}
\newtcolorbox{attention}[1][]{popbase,colframe=poppink,before upper={\popboxhead{Attention}{poppink}{#1}}}
\newtcolorbox{experience}[1][]{popbase,colframe=popblue,colback=popblueL,before upper={\popboxhead{Expérience}{popblue}{#1}}}
% « Le savais-tu ? » : carte violette pleine (contexte, culture, Cameroun)
\newtcolorbox{savaistu}[1][]{popbase,colback=poppurple,colframe=popink,
  fontupper=\popsemi\fontsize{10.4}{14.2}\selectfont\color{white},
  before upper={\poppill{Le savais-tu\,?}{white}{poppurple}\ifx\relax#1\relax\else\hspace{6pt}{\popxbold\color{white}#1}\fi\par\vspace{7pt}}}
% Exercice résolu : énoncé en haut, \tcblower puis la solution sur fond crème
\newcounter{popexo}\newcounter{popexr}
\newtcolorbox{exoresolu}[1][]{popbase,colframe=poporange,colbacklower=popcream,
  segmentation style={popink,line width=1.2pt,dashed},
  before upper={\stepcounter{popexr}\popboxhead{Exercice résolu \thepopexr}{poporange}{#1}},
  before lower={\poppill{Solution}{popink}{popcream}\par\vspace{6pt}}}
% Exercice (non résolu en place) — la correction est dans la partie Corrigés
\newtcolorbox{exercice}[1][]{popbase,colframe=popink,
  before upper={\stepcounter{popexo}\popboxhead{Exercice \thepopexo}{popink}{#1}}}
% Corrigé
\newtcolorbox{corrige}[1][]{popbase,colframe=popgreen,colback=popgreenL,
  before upper={\popboxhead{Corrigé}{popgreen}{#1}}}
% Objectifs / prérequis / bilan
\newtcolorbox{objectifs}{popbase,colframe=popink,colback=poporangeL,
  before upper={\popboxhead{Objectifs de la leçon}{popink}{}}}
\newtcolorbox{prerequis}{popbase,colframe=popink,colback=popblueL,
  before upper={\popboxhead{Prérequis}{popblue}{}}}
\newtcolorbox{bilan}{popbase,colframe=popink,colback=white,boxrule=2.8pt,
  before upper={\popboxhead{Fiche bilan}{poporange}{L'essentiel à connaître pour l'examen}}}
% Figure encadrée avec légende
\newtcolorbox{popfigure}[1][]{enhanced,breakable=false,colback=white,colframe=popline,boxrule=1.4pt,arc=8pt,
  left=10pt,right=10pt,top=10pt,bottom=8pt,before skip=10pt,after skip=12pt,halign=center,
  lower separated=false,halign lower=center,
  fontlower=\popsemi\fontsize{9}{11.5}\selectfont\color{popmuted},
  #1}
\newcommand{\legende}[1]{\par\vspace{6pt}{\popsemi\fontsize{9}{11.5}\selectfont\color{popmuted}#1}}

% ============================================================
% Signature OnBuch+
% ============================================================
\newcommand{\poplogoinline}[1]{{\poplogo\fontsize{#1}{#1}\selectfont\color{popink}OnBuch\color{poporange}+}}
\newcommand{\signatureonbuch}{%
  \begin{tcolorbox}[enhanced,colback=popink,colframe=popink,boxrule=2.2pt,arc=10pt,
      drop shadow=poporange,left=14pt,right=14pt,top=10pt,bottom=10pt,before skip=0pt,after skip=0pt]
    \tikz[baseline=-0.7ex]\node[circle,fill=popgreen,draw=popcream,line width=1.3pt,minimum size=22pt,inner sep=0]{\popcheck{white}};%
    \hspace{9pt}\begin{minipage}[c]{0.62\linewidth}
      {\popxbold\fontsize{12}{13}\selectfont\color{popcream}Signé\ \textbullet\ {\poplogo OnBuch\color{poporange}+}}\par\vspace{2pt}
      {\raggedright\popsemi\fontsize{8}{10.5}\selectfont\color{popline}\MakeUppercase{Équipe pédagogique OnBuch+}\\\MakeUppercase{Conforme au programme officiel MINESEC}\par}
    \end{minipage}\hfill
    {\poplogo\fontsize{22}{22}\selectfont\color{popcream}OnBuch\color{poporange}+}
  \end{tcolorbox}%
}

% ============================================================
% Page de garde
% #1 titre · #2 matière · #3 niveau · #4 module · #5 n° leçon · #6 durée ·
% #7 description / objectifs (texte ou itemize)
% ============================================================
\newcommand{\pagegarde}[7]{%
  \thispagestyle{empty}
  \newgeometry{a4paper,margin=1.7cm,top=1.6cm,bottom=1.4cm}%
  \begin{tikzpicture}[remember picture,overlay]
    \fill[poporange!18] ([xshift=-3.2cm,yshift=-3.6cm]current page.north east) circle (4.6cm);
    \fill[poppurple!14] ([xshift=2.4cm,yshift=7.6cm]current page.south west) circle (3.8cm);
  \end{tikzpicture}%
  {\poplogo\fontsize{30}{30}\selectfont\color{popink}OnBuch\color{poporange}+}\hfill
  \raisebox{0.6ex}{\poppill{Cours signé \textbullet\ Premium}{popink}{popcream}}\par
  \vspace{1pt}\popsquiggle{poppurple}\par
  \vspace{8pt}
  \begin{tcolorbox}[enhanced,colback=poporange,colframe=popink,boxrule=2.8pt,arc=13pt,
      drop shadow=popink,left=18pt,right=18pt,top=12pt,bottom=13pt,after skip=10pt]
    \poppill{#2 \textbullet\ #3}{popink}{popcream}\hspace{5pt}\poppill{#4}{white}{popink}\par\vspace{8pt}
    {\popdisplay\fontsize{14}{14}\selectfont\color{popink}LEÇON #5}\par\vspace{4pt}
    {\raggedright\hyphenpenalty=10000\exhyphenpenalty=10000\popdisplay\fontsize{25}{27}\selectfont\color{popcream}\MakeUppercase{#1}\par}
    \vspace{8pt}
    {\popxbold\fontsize{9.5}{9.5}\selectfont\color{popink}\MakeUppercase{Durée conseillée : #6}}
  \end{tcolorbox}
  \begin{tcolorbox}[enhanced,colback=white,colframe=popink,boxrule=2.2pt,arc=10pt,
      drop shadow=popink,left=16pt,right=16pt,top=10pt,bottom=10pt,after skip=8pt]
    \poppill{Dans cette leçon}{poppurple}{white}\par\vspace{5pt}
    \popsemi\fontsize{9.3}{12.6}\selectfont\color{popink}#7
  \end{tcolorbox}
  \noindent\begin{minipage}[t]{0.487\linewidth}
    \begin{tcolorbox}[enhanced,colback=popgreen,colframe=popink,boxrule=2.2pt,arc=10pt,
        drop shadow=popink,left=11pt,right=11pt,top=10pt,bottom=10pt,height=3.1cm,valign=center]
      \tcbox[colback=white,colframe=popink,boxrule=1.3pt,arc=5pt,boxsep=0pt,nobeforeafter,
        left=3pt,right=3pt,top=3pt,bottom=3pt,valign=center]{\qrcode[height=1.9cm]{https://play.google.com/store/apps/details?id=com.luvvix.onbuch&hl=fr}}%
      \hspace{8pt}\begin{minipage}[c]{0.5\linewidth}
        {\popdisplay\fontsize{11}{12.5}\selectfont\color{white}\MakeUppercase{Télécharge OnBuch}}\par\vspace{4pt}
        {\raggedright\popsemi\fontsize{8.4}{11}\selectfont\color{white}Gratuit \textbullet\ Google Play\\Tuteur IA, annales, concours, résultats\par}
      \end{minipage}
    \end{tcolorbox}
  \end{minipage}\hfill
  \begin{minipage}[t]{0.487\linewidth}
    \begin{tcolorbox}[enhanced,colback=poppurple,colframe=popink,boxrule=2.2pt,arc=10pt,
        drop shadow=popink,left=11pt,right=11pt,top=10pt,bottom=10pt,height=3.1cm,valign=center]
      \tikz[baseline=-0.7ex]\node[circle,fill=white,draw=popink,line width=1.3pt,minimum size=1.6cm,inner sep=0]{\popdisplay\fontsize{24}{24}\selectfont\color{poppurple}+};%
      \hspace{8pt}\begin{minipage}[c]{0.6\linewidth}
        {\popdisplay\fontsize{11}{12.5}\selectfont\color{white}\MakeUppercase{Passe à OnBuch+}}\par\vspace{4pt}
        {\raggedright\popsemi\fontsize{8.4}{11}\selectfont\color{white}Tous les cours signés, Tuteur IA illimité, fiches PDF — onglet OnBuch+ de l'app\par}
      \end{minipage}
    \end{tcolorbox}
  \end{minipage}
  \vspace{8pt}
  \popcontenu
  \vfill
  \signatureonbuch
  \restoregeometry
  \newpage
}

% Rangée « Ce cours contient » (page de garde)
\newcommand{\poptile}[3]{% #1 couleur #2 gros texte #3 légende
  \begin{tcolorbox}[enhanced,colback=#1,colframe=popink,boxrule=2pt,arc=9pt,shadow={2.5pt}{-2.5pt}{0pt}{popink},
      left=6pt,right=6pt,top=9pt,bottom=9pt,halign=center,valign=center,height=1.75cm,width=\linewidth,nobeforeafter]
    {\popdisplay\fontsize{12.5}{13}\selectfont\color{white}\MakeUppercase{#2}}\par\vspace{3pt}
    {\centering\popsemi\fontsize{7.8}{9.5}\selectfont\color{white}#3\par}
  \end{tcolorbox}}
\newcommand{\popcontenu}{%
  \par\noindent{\popxboldls\fontsize{7.5}{7.5}\selectfont\color{popmuted}CE COURS SIGNÉ CONTIENT}\par\vspace{7pt}
  \noindent\begin{minipage}[t]{0.235\linewidth}\poptile{popblue}{Cours}{complet, illustré, pas à pas}\end{minipage}\hfill
  \begin{minipage}[t]{0.235\linewidth}\poptile{poporange}{Méthodes}{et exercices résolus}\end{minipage}\hfill
  \begin{minipage}[t]{0.235\linewidth}\poptile{poppink}{Exercices}{type examen + corrigés}\end{minipage}\hfill
  \begin{minipage}[t]{0.235\linewidth}\poptile{popgreen}{Fiche bilan}{l'essentiel à retenir}\end{minipage}\par}

% Sommaire
\newtcolorbox{sommaire}{enhanced,colback=white,colframe=popink,boxrule=2.2pt,arc=10pt,
  drop shadow=popink,left=18pt,right=18pt,top=13pt,bottom=13pt,before skip=6pt,after skip=20pt,
  before upper={\poppill{Sommaire}{poppurple}{white}\par\vspace{9pt}\popsemi\fontsize{10.6}{14.5}\selectfont\color{popink}}}

% Signature de fin de cours — poussée tout en bas de la dernière page
\newcommand{\findecours}{%
  \par\vfill\nopagebreak
  \begin{center}\popsquiggle{poporange}\end{center}\vspace{4pt}
  \signatureonbuch
}
\AtBeginDocument{\def\llbracket{\mathopen{[\mkern-3.5mu[}}\def\rrbracket{\mathclose{]\mkern-3.5mu]}}} % intervalles d'entiers (glyphe ⟦ absent de la police)
% Glyphes absents de Fira Math : redéfinis à partir de glyphes disponibles
\AtBeginDocument{%
  \def\setminus{\mathbin{\backslash}}\let\smallsetminus\setminus
  \def\vdots{\mathord{\vbox{\baselineskip3.2pt\lineskiplimit0pt\kern4pt\hbox{.}\hbox{.}\hbox{.}}}}%
  \def\ddots{\mathinner{\mkern1mu\raise6.5pt\hbox{.}\mkern2mu\raise3.6pt\hbox{.}\mkern2mu\raise0.7pt\hbox{.}\mkern1mu}}%
  \def\triangle{\mathord{\scalebox{0.9}{$\symup{\Delta}$}}}%
  \def\nabla{\mathord{\raisebox{\depth}{\scalebox{1}[-1]{$\symup{\Delta}$}}}}%
  \def\checkmark{\popcheck{popdark}}%
  \def\star{\mathbin{\ast}}\def\bigstar{\mathord{\ast}}%
  \def\diamond{\mathbin{\scalebox{0.6}{$\Diamond$}}}%
  \def\bigcup{\mathop{\mathchoice{\raisebox{-0.3em}{\scalebox{1.7}{$\cup$}}}{\scalebox{1.25}{$\cup$}}{\cup}{\cup}}\displaylimits}%
  \def\bigcap{\mathop{\mathchoice{\raisebox{-0.3em}{\scalebox{1.7}{$\cap$}}}{\scalebox{1.25}{$\cap$}}{\cap}{\cap}}\displaylimits}%
  \def\lesseqqgtr{\mathrel{\lessgtr}}\def\bigoplus{\mathop{\oplus}\displaylimits}%
}
\providecommand{\ssi}{si et seulement si} % abréviation fréquente
\AtBeginDocument{\providecommand{\wideparen}{\overparen}} % arc de cercle

%% FIGFIX-BEGIN v4 — correctifs globaux des figures (docs/onbuch-generation/tools/figfix)
\usepackage{adjustbox}\usepackage{etoolbox}
% toute figure TikZ/pgfplots plus large que la ligne est réduite à la largeur disponible
\BeforeBeginEnvironment{tikzpicture}{\begin{adjustbox}{max width=\linewidth}}
\AfterEndEnvironment{tikzpicture}{\end{adjustbox}}
% légendes pgfplots lisibles par-dessus les courbes
\pgfplotsset{every axis/.append style={legend style={fill=white,fill opacity=0.92,text opacity=1,draw=black!15,inner sep=3pt}}}
% étiquettes d'arêtes (syntaxe edge["..."]) avec fond blanc
\tikzset{every edge quotes/.append style={fill=white,inner sep=1.2pt}}
%% FIGFIX-END
\begin{document}

\pagegarde{Lesson 9 — Vocabulary: Words and expressions related to recreational activities}{Anglais}{Tle A}{Chapter 1 : Family and Social Life}{EN09}{2 h}{Cette leçon de vocabulaire présente les mots et expressions liés aux loisirs et activités récréatives (sports, jeux, sorties, hobbies). Elle fournit le lexique par champs, les collocations, les faux amis et les familles de mots, avec une mise en pratique orale et écrite en vue du Bac.
\vspace{4pt}
\begin{multicols}{2}\small
\begin{itemize}
\item À la fin de cette leçon, je sais nommer et classer les principales activités récréatives en anglais
\item À la fin de cette leçon, je sais employer les bonnes collocations (play / do / go + activité)
\item À la fin de cette leçon, je sais utiliser les expressions idiomatiques liées aux loisirs (to kill time, to take up a hobby...)
\item À la fin de cette leçon, je sais éviter les faux amis et les calques du français (to assist, to make sport...)
\item À la fin de cette leçon, je sais former des mots de la même famille (relax, relaxation, relaxing, relaxed)
\item À la fin de cette leçon, je sais prononcer correctement le lexique clé (leisure, hobby, amateur...)
\end{itemize}\end{multicols}}

\begin{sommaire}
\begin{multicols}{2}
\begin{enumerate}
\item Champs lexicaux : les grandes familles de loisirs
\item Collocations : play, do, go et autres verbes associés
\item Expressions idiomatiques et langage courant des loisirs
\item Faux amis et pièges des francophones
\item Familles de mots, prononciation et orthographe
\item Emploi en contexte : dialogues et court texte
\item Activité d'intégration
\item Méthodes et exercices résolus
\item Exercices
\item Corrigés des exercices
\item Fiche bilan
\end{enumerate}
\end{multicols}
\end{sommaire}

\begin{objectifs}
\begin{itemize}
\item À la fin de cette leçon, je sais nommer et classer les principales activités récréatives en anglais
\item À la fin de cette leçon, je sais employer les bonnes collocations (play / do / go + activité)
\item À la fin de cette leçon, je sais utiliser les expressions idiomatiques liées aux loisirs (to kill time, to take up a hobby...)
\item À la fin de cette leçon, je sais éviter les faux amis et les calques du français (to assist, to make sport...)
\item À la fin de cette leçon, je sais former des mots de la même famille (relax, relaxation, relaxing, relaxed)
\item À la fin de cette leçon, je sais prononcer correctement le lexique clé (leisure, hobby, amateur...)
\item À la fin de cette leçon, je sais parler de mes loisirs et de ceux des autres dans un dialogue
\item À la fin de cette leçon, je sais rédiger un court paragraphe sur les activités récréatives de ma communauté
\end{itemize}
\end{objectifs}

\begin{prerequis}
\begin{itemize}
\item Les verbes de goût et leurs constructions (like, enjoy, prefer + -ing) vus en Première
\item Le présent simple pour exprimer les habitudes
\item Le vocabulaire de base des sports et jeux (football, cards, dancing)
\item Les expressions de fréquence (often, twice a week, at weekends)
\item La distinction de base entre verbe, nom et adjectif
\end{itemize}
\end{prerequis}

\newpage

\coursec{Champs lexicaux : les grandes familles de loisirs}

\courssub{Situation d'accroche et problématique}

\begin{textebox}[Saturday afternoon in Bamenda]
It's Saturday afternoon in Bamenda. Some students are playing football at the commercial avenue field, others are watching a Nigerian movie at home, and a group of friends has gone hiking up Mount Cameroon for the weekend. When the teacher asks on Monday: \eng{“What did you do for recreation?”}, many students hesitate: \eng{“I made sport”}, \eng{“I assisted to a match”}... These are typical mistakes! How do English speakers really talk about their free time and hobbies?
\end{textebox}

\textbf{Problématique.} Ces élèves de Bamenda ont bien passé leur week-end, mais ils ne savent pas le raconter en anglais correct. Ils calquent le français : \eng{I made sport} (calque de « j'ai fait du sport »), \eng{I assisted to a match} (calque de « j'ai assisté à un match »). Pour parler naturellement de ses loisirs, il faut d'abord maîtriser le \cle{champ lexical} des activités récréatives : les mots justes, rangés par familles, avec leurs verbes associés. C'est l'objectif de cette première section.

\courssub{Les mots-clés : recreation, leisure, free time, pastime, hobby}

Avant de classer les activités, observons les mots qui désignent le temps libre lui-même. Lisez ces phrases :

\begin{exemplebox}[Observer avant de définir]
\begin{itemize}
\item \eng{“What do you do in your spare time?”} « Que fais-tu pendant ton temps libre ? »
\item \eng{“My favourite pastime is fishing.”} « Mon passe-temps préféré est la pêche. »
\item \eng{“Football is a popular recreation in Cameroon.”} « Le football est un loisir populaire au Cameroun. »
\item \eng{“She reads at her leisure.”} « Elle lit à loisir, à son rythme. »
\item \eng{“His hobby is collecting stamps.”} « Son hobby est la philatélie. »
\end{itemize}
\end{exemplebox}

\begin{definition}[Recreation, leisure, free time, pastime, hobby]
\begin{itemize}
\item \cle{recreation} (nom, plutôt soutenu) : une \emph{activité} pratiquée pendant le temps libre pour se détendre ou se divertir. Prononciation : « rè-kri-é-cheune ».
\item \cle{leisure} (nom, surtout GB) : le \emph{temps libre} lui-même, non l'activité. Prononciation : « lè-jeur ». Expression figée : \eng{at leisure} = « à loisir, sans se presser » ; \eng{leisure activities} = « activités de loisir ».
\item \cle{free time} / \cle{spare time} (noms, courants) : le temps libre, synonymes parfaits. On dit \eng{in my free time} ou \eng{in my spare time} (préposition \eng{in}).
\item \cle{pastime} (nom comptable) : un passe-temps, une activité qui fait passer le temps agréablement.
\item \cle{hobby} (nom comptable, pluriel \eng{hobbies}) : un loisir \emph{régulier et personnel}, souvent pratiqué seul (collectionner, coudre, jardiner).
\end{itemize}
\end{definition}

\begin{attention}[Faux amis et calques]
\begin{itemize}
\item On ne dit jamais \eng{I made sport} : le verbe est \eng{do} ou \eng{play} → \eng{I do sport} / \eng{I play football} (voir section 2).
\item \eng{assist} ne signifie pas « assister à » : \eng{I attended a match} ou \eng{I watched a match}. \eng{Assist} veut dire « aider ».
\item \eng{leisure} ne se prononce pas « lé-zure » mais « lè-jeur » (GB) ; les Américains disent « li-jeur ».
\item \eng{hobby} prend un \textbf{-ies} au pluriel : \eng{hobbies}, comme \eng{family → families}.
\end{itemize}
\end{attention}

\courssub{Texte support : recreation in Cameroon}

\begin{textebox}[Free time, from town to village]
In Cameroon, recreation varies from town to village. In Yaoundé and Douala, many young people spend their free time in front of a screen: they go to the cinema, play video games at game centres, or watch football matches in small video clubs. At the weekend, families often go out together: some attend weddings and traditional ceremonies, others visit relatives in another quarter of the town.

In rural areas, life follows a different rhythm. In the evening, children sit around their grandparents and listen to folk tales; storytelling remains a favourite pastime in many villages. During holidays, boys organise football matches on dusty playgrounds, while girls practise traditional dances for the next feast. In the North-West and South-West regions, the songo'o, a clever board game played with seeds or stones, still brings young and old together. Near the coast, in Limbe or Kribi, families spend Sundays at the beach, swimming and eating grilled fish.

Whatever the place, recreation has the same function everywhere: it helps people relax, make friends and keep their culture alive. As a teacher in Buea says: \eng{“A student who plays, sings or reads after class comes back to school with a fresh mind.”}
\end{textebox}

\begin{definition}[Glossaire du texte]
\begin{itemize}
\item \eng{screen} (n.) : écran ; \eng{game centre} (n.) : salle de jeux vidéo
\item \eng{to attend} (v.) : assister à, être présent à ; \eng{relatives} (n. pl.) : parents, membres de la famille
\item \eng{folk tale} (n.) : conte populaire ; \eng{storytelling} (n.) : l'art de raconter des histoires
\item \eng{dusty} (adj.) : poussiéreux ; \eng{playground} (n.) : terrain de jeu
\item \eng{feast} (n.) : fête ; \eng{board game} (n.) : jeu de société
\item \eng{to relax} (v.) : se détendre ; \eng{a fresh mind} (n.) : un esprit reposé
\end{itemize}
\end{definition}

\textbf{Questions de compréhension (oral).}
\begin{enumerate}
\item \eng{What do young people in Yaoundé do in their free time?}
\item \eng{Which pastimes are mentioned in rural areas?}
\item \eng{What is the songo'o and where is it played?}
\item \eng{According to the teacher in Buea, why is recreation important?}
\end{enumerate}

\courssub{Les cinq champs lexicaux des loisirs}

\begin{definition}[Champ lexical]
Un \cle{champ lexical} est un ensemble de mots qui se rapportent à un même thème. Classer le vocabulaire par champs aide la mémoire : au lieu d'apprendre des listes désordonnées, on retient des familles de mots liés entre eux.
\end{definition}

\textbf{Champ 1 — Sports et activités physiques.}
\begin{itemize}
\item \emph{Team sports} (sports collectifs) : \eng{football} (football), \eng{volleyball} (volley-ball), \eng{basketball} (basket-ball), \eng{handball} (handball), \eng{rugby} (rugby).
\item \emph{Individual sports} (sports individuels) : \eng{athletics} (athlétisme), \eng{swimming} (natation), \eng{cycling} (cyclisme), \eng{boxing} (boxe), \eng{wrestling} (lutte).
\item \emph{Outdoor activities} (activités de plein air) : \eng{hiking} (randonnée), \eng{jogging} (footing), \eng{camping} (camping), \eng{fishing} (pêche), \eng{hunting} (chasse).
\end{itemize}

\textbf{Champ 2 — Jeux et divertissements.}
\begin{itemize}
\item \emph{Board games} (jeux de société) : \eng{draughts} (jeu de dames), \eng{chess} (échecs, toujours singulier), \eng{ludo} (jeu de ludo / petits chevaux), \eng{scrabble} (scrabble).
\item \emph{Card games} (jeux de cartes) : \eng{playing cards} (jouer aux cartes), \eng{whist} (whist).
\item \emph{Electronic and modern games} : \eng{video games} (jeux vidéo), \eng{puzzles} (puzzles, casse-tête).
\item \emph{Traditional games} (jeux traditionnels) : \eng{songo'o} (songo'o, jeu de semis), \eng{dancing competitions} (concours de danse), \eng{hide-and-seek} (cache-cache).
\end{itemize}

\textbf{Champ 3 — Loisirs culturels et artistiques.}
\begin{itemize}
\item \eng{reading} (lecture), \eng{listening to music} (écouter de la musique), \eng{watching films} (regarder des films), \eng{dancing} (danse), \eng{singing in a choir} (chanter dans une chorale), \eng{painting} (peinture), \eng{drama} (théâtre), \eng{storytelling} (conte), \eng{photography} (photographie).
\end{itemize}

\textbf{Champ 4 — Loisirs sociaux.}
\begin{itemize}
\item \eng{chatting with friends} (bavarder avec des amis), \eng{going to parties} (aller à des fêtes), \eng{attending ceremonies} (assister à des cérémonies), \eng{visiting relatives} (rendre visite à la famille), \eng{hanging out} (traîner, passer du temps dehors, familier).
\end{itemize}

\textbf{Champ 5 — Les lieux des loisirs.}
\begin{itemize}
\item \eng{stadium} (stade), \eng{playground} (terrain de jeu), \eng{leisure centre} (centre de loisirs, GB), \eng{cinema} (cinéma), \eng{park} (parc), \eng{beach} (plage), \eng{community hall} (salle communautaire / des fêtes), \eng{swimming pool} (piscine).
\end{itemize}

\begin{savaistu}[Leisure centre ou rec center ?]
Au Royaume-Uni, on dit \eng{leisure centre} (prononcé « lè-jeur sèn-teur ») pour désigner le complexe sportif et culturel municipal. Aux États-Unis, on préfère \eng{recreation center}, abrégé familièrement en \eng{rec center} (« rèk sèn-teur »). Attention aussi à l'orthographe : \eng{centre} en anglais britannique, \eng{center} en anglais américain. Au Cameroun, l'anglais britannique est la norme à l'école.
\end{savaistu}

\courssub{Carte mentale et tableau récapitulatif}

\begin{popfigure}
\begin{tikzpicture}[mindmap, grow cyclic,
every node/.style={concept, align=center, font=\small},
root concept/.append style={concept color=popblue, font=\small\bfseries},
level 1/.append style={level distance=3.6cm, sibling angle=72},
level 2/.append style={level distance=2.2cm, sibling angle=40, font=\scriptsize}]
\node[root concept]{RECREATIONAL ACTIVITIES}
  child[concept color=poporange] { node {Sports} 
      child { node {football} } child { node {swimming} } child { node {hiking} } }
  child[concept color=popgreen] { node {Games} 
      child { node {chess} } child { node {songo'o} } child { node {video games} } }
  child[concept color=poppurple] { node {Arts \& Culture} 
      child { node {reading} } child { node {dancing} } child { node {storytelling} } }
  child[concept color=poppink] { node {Social life} 
      child { node {parties} } child { node {chatting} } child { node {ceremonies} } }
  child[concept color=popgold] { node {Places} 
      child { node {stadium} } child { node {beach} } child { node {cinema} } };
\end{tikzpicture}
\legende{Carte mentale : les cinq grandes familles d'activités récréatives, avec trois exemples par branche.}
\end{popfigure}

\begin{center}
\begin{tabularx}{\linewidth}{|l|X|X|X|}
\hline
\rowcolor{popblueL} \textbf{Champ} & \textbf{English} & \textbf{Français} & \textbf{Exemple en contexte} \\
\hline
Sports & football, athletics, hiking & football, athlétisme, randonnée & \eng{We go hiking during holidays.} \\
\hline
Games & chess, draughts, ludo & échecs, dames, ludo & \eng{My grandfather plays draughts.} \\
\hline
Arts \& culture & reading, drama, singing in a choir & lecture, théâtre, chant choral & \eng{She sings in the church choir.} \\
\hline
Social life & parties, ceremonies, chatting & fêtes, cérémonies, bavardage & \eng{They attended a wedding on Saturday.} \\
\hline
Places & stadium, beach, park & stade, plage, parc & \eng{The match is at the stadium.} \\
\hline
\end{tabularx}
\end{center}

\begin{methode}[Mémoriser le vocabulaire des loisirs]
\begin{enumerate}
\item \textbf{Regroupe} les mots par champ lexical (sports, jeux, arts, vie sociale, lieux) : le cerveau retient mieux les familles que les listes mélangées.
\item \textbf{Apprends chaque mot avec son verbe associé} : \eng{go hiking}, \eng{play chess}, \eng{listen to music}, \eng{watch a film}. Un mot isolé s'oublie ; une collocation se réutilise.
\item \textbf{Fabrique une phrase personnelle} pour chaque mot nouveau : \eng{My favourite pastime is fishing.} Une phrase vraie pour toi est inoubliable.
\item \textbf{Dessine ta propre carte mentale} au brouillon et teste-toi en cachant les branches.
\end{enumerate}
\end{methode}

\courssub{Exercice résolu et entraînement}

\begin{exoresolu}[Classer dans les cinq champs]
Énoncé — Classez les activités suivantes dans le bon champ lexical (Sports / Games / Arts \& Culture / Social life / Places) : \eng{swimming, storytelling, stadium, chess, attending a wedding, jogging, video games, beach, singing in a choir, hanging out with friends}.
\tcblower
Solution détaillée —
\begin{itemize}
\item \textbf{Sports} : \eng{swimming} (sport individuel), \eng{jogging} (activité de plein air).
\item \textbf{Games} : \eng{chess} (jeu de société), \eng{video games} (jeux électroniques).
\item \textbf{Arts \& Culture} : \eng{storytelling} (conte), \eng{singing in a choir} (activité musicale).
\item \textbf{Social life} : \eng{attending a wedding} (cérémonie), \eng{hanging out with friends} (loisir social).
\item \textbf{Places} : \eng{stadium} (lieu de sport), \eng{beach} (lieu de détente).
\end{itemize}
Astuce : les \emph{lieux} ne sont pas des activités ; on les reconnaît parce qu'on y va (\eng{go to the stadium / the beach}).
\end{exoresolu}

\begin{exercice}[À toi de jouer]
\begin{enumerate}
\item Trouve l'intrus dans chaque série et justifie en anglais : (a) \eng{football, volleyball, choir, basketball} ; (b) \eng{chess, draughts, swimming, ludo} ; (c) \eng{stadium, cinema, beach, fishing}.
\item Complète avec un mot du champ lexical : \eng{At the weekend, my uncle goes \_\_\_\_\_\_ by the river; he loves catching fish.} / \eng{In the evening, the children listen to \_\_\_\_\_\_ told by their grandmother.}
\item Traduis en anglais : « Mon passe-temps préféré est la lecture. » et « Que fais-tu pendant ton temps libre ? »
\end{enumerate}
\end{exercice}

\begin{corrige}[Exercice « À toi de jouer »]
\begin{enumerate}
\item (a) \eng{choir} : it belongs to arts and culture, not team sports. (b) \eng{swimming} : it is a sport, not a board game. (c) \eng{fishing} : it is an activity, not a place.
\item \eng{fishing} ; \eng{stories / folk tales}.
\item \eng{My favourite pastime is reading.} ; \eng{What do you do in your spare time?} (ou \eng{free time}).
\end{enumerate}
\end{corrige}

\begin{aretenir}[L'essentiel de la section]
\begin{itemize}
\item \eng{Recreation} désigne l'activité, \eng{leisure} et \eng{free/spare time} désignent le temps libre, \eng{pastime} et \eng{hobby} désignent un loisir précis.
\item Le vocabulaire des loisirs se range en cinq champs : \textbf{sports}, \textbf{games}, \textbf{arts \& culture}, \textbf{social life}, \textbf{places}.
\item On apprend chaque mot \emph{avec son verbe} : \eng{play football}, \eng{go hiking}, \eng{listen to music}.
\item Jamais \eng{make sport} ni \eng{assist to} : on dit \eng{do sport} et \eng{attend / watch a match}.
\end{itemize}
\end{aretenir}

\coursec{Collocations : play, do, go et autres verbes associés}

Dans la section précédente, nous avons classé les loisirs en grandes familles : sports, jeux, arts, activités de plein air. Mais connaître le nom d'une activité ne suffit pas pour parler correctement : il faut savoir \cle{quel verbe} l'accompagne. En français, un seul verbe, \emph{faire}, couvre presque tout : \emph{faire du football, faire du karaté, faire de la natation, faire du piano}. L'anglais, lui, répartit ces activités entre trois verbes principaux : \eng{play}, \eng{do} et \eng{go}. C'est l'une des collocations les plus testées à l'examen — et l'une des plus fautives chez les francophones.

\courssub{Observer avant de généraliser}

Lisons d'abord ce court dialogue entre deux camarades de classe à Yaoundé.

\begin{textebox}[At the end of the week — a dialogue]
\textbf{Ngong:} What do you usually do at the weekend, Amina?\par
\textbf{Amina:} On Saturdays, I \textbf{play} volleyball with my school team, and in the evening I sometimes \textbf{go} dancing with my cousins. And you?\par
\textbf{Ngong:} I \textbf{play} chess with my grandfather — he taught me when I was seven. On Sundays I \textbf{go} swimming at the municipal pool.\par
\textbf{Amina:} Do you still \textbf{do} karate?\par
\textbf{Ngong:} Yes, twice a week. My brother prefers music: he \textbf{plays} the guitar in a band.\par
\textbf{Amina:} Lucky him! I tried to \textbf{take up} the piano last year, but I \textbf{gave it up} after two months. I had no \textbf{free time} left.\par
\textbf{Ngong:} You should \textbf{join} a music club at school. It is a good way to \textbf{spend} your \textbf{spare time}.
\end{textebox}

\begin{definition}[Glossaire]
\begin{itemize}
\item \eng{chess} (n.) : les échecs
\item \eng{municipal pool} (n.) : piscine municipale
\item \eng{twice a week} : deux fois par semaine
\item \eng{band} (n.) : un groupe de musique
\item \eng{to take up} (phrasal verb) : commencer (une activité)
\item \eng{to give up} (phrasal verb) : abandonner
\item \eng{to join} (v.) : adhérer à, s'inscrire à
\end{itemize}
\end{definition}

\textbf{Questions d'observation.} (1) Quelles activités suivent \eng{play} ? (2) Quelles activités suivent \eng{do} ? (3) Quelles activités suivent \eng{go} ? (4) Que remarquez-vous sur la forme des activités qui suivent \eng{go} ?

Réponses : \eng{play} + volleyball, chess, the guitar, the piano ; \eng{do} + karate ; \eng{go} + dancing, swimming — toutes deux en \cle{-ing}. Nous pouvons maintenant énoncer la règle.

\courssub{La règle des trois verbes : play, do, go}

\begin{propriete}[Règle générale : play / do / go]
\begin{enumerate}
\item \cle{PLAY} + les sports avec ballon, les jeux et les instruments de musique : \eng{play football, play basketball, play chess, play cards, play the guitar}.
\item \cle{DO} + les activités individuelles, les arts martiaux et les activités de forme physique : \eng{do athletics, do karate, do judo, do gymnastics, do yoga, do aerobics}.
\item \cle{GO} + les activités dont le nom est un verbe en \cle{-ing} : \eng{go swimming, go hiking, go fishing, go shopping, go dancing, go jogging, go cycling}.
\end{enumerate}
\end{propriete}

\begin{exemplebox}[Exemples traduits]
\begin{itemize}
\item \eng{We play volleyball on Saturdays.} « Nous jouons au volley le samedi. »
\item \eng{She goes jogging every morning.} « Elle fait du jogging chaque matin. »
\item \eng{He does karate.} « Il fait du karaté. »
\item \eng{They play chess during the break.} « Ils jouent aux échecs pendant la récréation. »
\item \eng{My father goes fishing on the Wouri river.} « Mon père va à la pêche sur le fleuve Wouri. »
\item \eng{We do gymnastics on Tuesdays.} « Nous faisons de la gymnastique le mardi. »
\end{itemize}
\end{exemplebox}

\begin{attention}[L'article : play the piano mais play football]
Les \cle{instruments de musique} prennent l'article \eng{the} : \eng{play the piano, play the guitar, play the drums}. Les \cle{sports et jeux} ne prennent \textbf{pas} d'article : \eng{play football} (jamais \eng{play the football}). De même, pas d'article après \eng{do} et \eng{go} : \eng{do yoga}, \eng{go swimming}.
\end{attention}

\begin{attention}[Le piège du calque : « faire du sport »]
Ne traduisez jamais \emph{faire du sport} par \eng{to make sport} : c'est un calque du français, toujours faux. On dit \eng{to do sport} (anglais britannique) ou \eng{to play sports} (anglais américain). De même : \emph{faire du vélo} = \eng{to go cycling} ou \eng{to ride a bike}, jamais \eng{to make bike}.
\end{attention}

\begin{center}
\begin{tabularx}{\linewidth}{|X|X|X|}
\hline
\rowcolor{popblueL} Français & English & Remarque \\
\hline
faire du football & play football & sport avec ballon, sans article \\
\hline
faire du karaté & do karate & art martial, sans article \\
\hline
faire de la natation & go swimming & activité en -ing \\
\hline
jouer du piano & play the piano & instrument : avec \eng{the} \\
\hline
jouer aux échecs & play chess & jeu : sans article \\
\hline
faire du sport & do sport / play sports (US) & jamais \eng{make sport} \\
\hline
\end{tabularx}
\end{center}

\begin{popfigure}
\begin{tikzpicture}
\node[draw=popblue, fill=popblueL, rounded corners, align=center, font=\bfseries] (play) at (0,0) {PLAY\\ \footnotesize ball games, games,\\ \footnotesize instruments};
\node[draw=poporange, fill=poporangeL, rounded corners, align=center, font=\bfseries] (do) at (6,0) {DO\\ \footnotesize individual sports,\\ \footnotesize martial arts};
\node[draw=popgreen, fill=popgreenL, rounded corners, align=center, font=\bfseries] (go) at (12,0) {GO\\ \footnotesize + activity in -ing};
\node[draw=popblue, rounded corners, align=center, fill=white] at (0,-1.6) {\footnotesize football, volleyball,\\ \footnotesize basketball, tennis};
\node[draw=popblue, rounded corners, align=center, fill=white] at (0,-3) {\footnotesize chess, cards,\\ \footnotesize video games};
\node[draw=popblue, rounded corners, align=center, fill=white] at (0,-4.4) {\footnotesize the piano, the guitar,\\ \footnotesize the drums};
\node[draw=poporange, rounded corners, align=center, fill=white] at (6,-1.6) {\footnotesize karate, judo,\\ \footnotesize taekwondo};
\node[draw=poporange, rounded corners, align=center, fill=white] at (6,-3) {\footnotesize gymnastics, yoga,\\ \footnotesize aerobics, athletics};
\node[draw=popgreen, rounded corners, align=center, fill=white] at (12,-1.6) {\footnotesize swimming, jogging,\\ \footnotesize hiking, cycling};
\node[draw=popgreen, rounded corners, align=center, fill=white] at (12,-3) {\footnotesize fishing, shopping,\\ \footnotesize dancing, camping};
\draw[-{Stealth}, popblue, thick] (play) -- (0,-1.1);
\draw[-{Stealth}, poporange, thick] (do) -- (6,-1.1);
\draw[-{Stealth}, popgreen, thick] (go) -- (12,-1.1);
\end{tikzpicture}
\legende{Carte mentale : la répartition des activités entre play, do et go.}
\end{popfigure}

\courssub{Autres collocations essentielles des loisirs}

Au-delà des trois verbes vedettes, le vocabulaire des loisirs repose sur des collocations très fréquentes qu'il faut mémoriser comme des blocs.

\begin{definition}[Collocations avec les activités]
\begin{itemize}
\item \eng{watch a match} : regarder un match — \eng{We watched the match between Cameroon and Senegal.} « Nous avons regardé le match entre le Cameroun et le Sénégal. »
\item \eng{attend a concert} : assister à un concert — \eng{Thousands attended the concert in Douala.} « Des milliers de personnes ont assisté au concert à Douala. »
\item \eng{take part in a competition} : participer à une compétition — \eng{Our school took part in the regional competition.} « Notre école a participé à la compétition régionale. »
\item \eng{join a club} : adhérer à un club — \eng{She joined the drama club last term.} « Elle a rejoint le club de théâtre le trimestre dernier. »
\item \eng{practise an activity} : pratiquer, s'entraîner à une activité — \eng{He practises the piano every evening.} « Il pratique le piano chaque soir. »
\end{itemize}
\end{definition}

\begin{attention}[Orthographe : practise ou practice ?]
En anglais britannique, le \cle{verbe} s'écrit \eng{practise} et le \cle{nom} \eng{practice} : \eng{I practise football; football practice is at 4 p.m.} En anglais américain, les deux s'écrivent \eng{practice}.
\end{attention}

\begin{definition}[Collocations avec hobby]
\begin{itemize}
\item \eng{take up a hobby} : commencer un passe-temps — \eng{I took up photography during the holidays.} « Je me suis mis à la photographie pendant les vacances. »
\item \eng{pursue a hobby} : poursuivre, cultiver un passe-temps — \eng{She pursues her hobby with passion.} « Elle poursuit son passe-temps avec passion. »
\item \eng{give up a hobby} : abandonner un passe-temps — \eng{He gave up stamp collecting years ago.} « Il a abandonné la philatélie il y a des années. »
\end{itemize}
\end{definition}

\begin{definition}[Collocations avec time]
\begin{itemize}
\item \eng{free time / spare time} : temps libre — \eng{What do you do in your free time?} « Que fais-tu pendant ton temps libre ? »
\item \eng{spend time} : passer du temps — \eng{I spend a lot of time reading.} « Je passe beaucoup de temps à lire. »
\item \eng{kill time} : tuer le temps — \eng{We played cards to kill time at the bus station.} « Nous avons joué aux cartes pour tuer le temps à la gare routière. »
\item \eng{pass the time} : passer le temps — \eng{She listens to music to pass the time.} « Elle écoute de la musique pour passer le temps. »
\end{itemize}
\end{definition}

\begin{methode}[Comment mémoriser les collocations]
\begin{enumerate}
\item N'apprenez \textbf{jamais} une activité seule : apprenez toujours le verbe avec elle. Écrivez \eng{go fishing} dans votre cahier, pas seulement \eng{fishing}.
\item Constituez trois colonnes (play / do / go) et ajoutez chaque nouvelle activité dans la bonne colonne.
\item Testez-vous régulièrement : cachez la colonne des verbes et retrouvez-les.
\item Faites une phrase complète avec chaque collocation : c'est la phrase, non le mot isolé, qui reste en mémoire.
\end{enumerate}
\end{methode}

\begin{savaistu}[Pourquoi « the » devant les instruments ?]
La tradition anglaise considère l'instrument comme un objet précis, identifié : \eng{the piano} désigne « le piano » comme catégorie d'instrument. Le sport, lui, est vu comme une activité abstraite, d'où l'absence d'article : \eng{play football}. Curiosité : en anglais américain, on entend parfois \eng{play piano} sans article, surtout à l'oral — mais à l'examen, exigez de vous-même la forme britannique \eng{play the piano}.
\end{savaistu}

\begin{exoresolu}[Classer les activités]
Énoncé. Classez les 15 activités suivantes sous \eng{play}, \eng{do} ou \eng{go} : basketball, yoga, swimming, chess, the drums, athletics, hiking, volleyball, karate, fishing, the violin, cards, gymnastics, dancing, jogging.
\tcblower
Solution détaillée.
\begin{itemize}
\item \textbf{PLAY} (ballons, jeux, instruments) : play basketball, play volleyball (sports avec ballon) ; play chess, play cards (jeux) ; play the drums, play the violin (instruments, avec \eng{the}).
\item \textbf{DO} (individuel, arts martiaux, forme) : do yoga, do athletics, do karate, do gymnastics.
\item \textbf{GO} (verbe en -ing) : go swimming, go hiking, go fishing, go dancing, go jogging.
\end{itemize}
Astuce de vérification : si l'activité se termine en \eng{-ing}, elle va presque toujours avec \eng{go} ; si c'est un art martial ou une gymnastique, avec \eng{do} ; si un ballon, un jeu ou un instrument est impliqué, avec \eng{play}.
\end{exoresolu}

\begin{exercice}[À vous de jouer]
\begin{enumerate}
\item Complétez avec \eng{play}, \eng{do} ou \eng{go} à la forme correcte : (a) Every Sunday, we \ldots\ hiking near Mount Cameroon. (b) My sister \ldots\ the guitar very well. (c) They \ldots\ judo at the sports centre. (d) I \ldots\ tennis with my friends after school. (e) She \ldots\ shopping in Douala every month.
\item Traduisez : « Pendant mon temps libre, je fais de la natation et je joue aux échecs. »
\end{enumerate}
\end{exercice}

\begin{corrige}[Exercice « À vous de jouer »]
1. (a) \eng{go hiking} ; (b) \eng{plays the guitar} ; (c) \eng{do judo} ; (d) \eng{play tennis} ; (e) \eng{goes shopping}.\par
2. \eng{In my free time, I go swimming and I play chess.} — Attention : pas d'article devant \eng{chess}, et \eng{go swimming} pour traduire « faire de la natation ».
\end{corrige}

\begin{aretenir}[L'essentiel de la section]
\begin{itemize}
\item \eng{play} + sports avec ballon, jeux, instruments (avec \eng{the} pour les instruments).
\item \eng{do} + sports individuels, arts martiaux, gymnastique, yoga.
\item \eng{go} + activités en \eng{-ing} (swimming, hiking, fishing\ldots).
\item Autres collocations clés : \eng{watch a match, attend a concert, take part in a competition, join a club, take up / give up a hobby, spend free time}.
\item Jamais \eng{make sport} : dites \eng{do sport} ou \eng{play sports} (US).
\end{itemize}
\end{aretenir}

\coursec{Expressions idiomatiques et langage courant des loisirs}

Dans la section précédente, nous avons appris à combiner correctement les verbes \eng{play}, \eng{do} et \eng{go} avec les noms d'activités : \eng{play football}, \eng{do yoga}, \eng{go swimming}. Mais pour parler de ses loisirs comme un vrai anglophone, il ne suffit pas de nommer les activités : il faut aussi savoir dire qu'on \emph{aime} une activité, qu'on s'y est \emph{mis}, qu'on s'en sert pour \emph{se détendre}. C'est exactement le rôle des expressions idiomatiques, que les Anglais et les Américains utilisent constamment dans la conversation quotidienne. Cette section vous donne les clés de ce langage vivant : d'abord une définition claire, puis un texte d'observation, ensuite les grandes familles d'expressions, et enfin des exercices pour les réemployer à l'oral comme à l'écrit — une compétence directement évaluée au Baccalauréat, dans les épreuves de compréhension et de composition.

\courssub{Qu'est-ce qu'une expression idiomatique ?}

\begin{definition}[Idiomatic expression]
Une \cle{expression idiomatique} (ou \emph{idiom}) est un groupe de mots figé dont le sens global est différent du sens littéral de chaque mot pris séparément. On ne peut pas la comprendre ni la traduire mot à mot : il faut l'apprendre comme un tout. Par exemple, \eng{to let off steam} signifie « se défouler, évacuer son stress », et non « laisser sortir de la vapeur » ; \eng{it's not my cup of tea} signifie « ce n'est pas mon truc », et non « ce n'est pas ma tasse de thé ».
\end{definition}

\begin{attention}[Le piège de la traduction littérale]
Ne traduisez JAMAIS une expression idiomatique mot à mot, ni de l'anglais vers le français, ni du français vers l'anglais. \eng{To kill time} ne signifie pas « tuer quelqu'un » mais « occuper son temps libre en attendant ». De même, « se défouler » ne se traduit pas par une invention mot à mot du type \eng{to un-blow oneself}, mais par \eng{to let off steam}. Retenez chaque idiome comme un bloc indivisible, avec sa structure et son registre de langue.
\end{attention}

\courssub{Observation : un texte pour découvrir les expressions}

Lisez d'abord ce court texte. Relevez les expressions qui ne se comprennent pas littéralement, puis vérifiez vos réponses avec le glossaire.

\begin{textebox}[Friday evening in Yaoundé — a student's diary]
After a long week at school, I need to unwind. On Friday evenings, I usually chill out with my friends at a snack bar near the lycée, or I listen to makossa at home. Football is not really my cup of tea, but I'm crazy about dancing. My best friend Clarisse is keen on photography; she took it up last year and now she never goes anywhere without her camera. My brother, on the other hand, is a real fan of video games. He says they help him to kill time while he waits for the bendskin in the morning. My father thinks he should take his time and find a more useful hobby, but my mother says everyone needs some me time after a hard week. As for me, when I feel stressed before an exam, I go jogging to let off steam. It helps me to recharge my batteries and to keep fit at the same time. Last Saturday, our team lost a match, but our captain was a good sport and congratulated the winners. That's what I call having a good time: relaxing, laughing, and not taking yourself too seriously!
\end{textebox}

\textbf{Glossaire}

\begin{center}
\begin{tabularx}{\linewidth}{|l|l|X|}
\hline
\rowcolor{popblueL} \textbf{English} & \textbf{Nature} & \textbf{Français} \\
\hline
to unwind & verbe & se détendre, décompresser \\
\hline
to chill out & verbe (familier) & se relaxer, traîner tranquillement \\
\hline
snack bar & nom & snack, buvette \\
\hline
to take up (a hobby) & verbe & se mettre à (une activité) \\
\hline
to kill time & expression & tuer le temps, s'occuper en attendant \\
\hline
me time & nom (familier) & du temps pour soi \\
\hline
to let off steam & expression & se défouler, évacuer le stress \\
\hline
to recharge one's batteries & expression & recharger ses batteries, récupérer \\
\hline
to keep fit & expression & rester en forme \\
\hline
a good sport & nom & quelqu'un de fair-play, un bon joueur \\
\hline
bendskin & nom (Cameroun) & moto-taxi \\
\hline
\end{tabularx}
\end{center}

\textbf{Questions de compréhension}
\begin{enumerate}
\item What does the writer usually do on Friday evenings?
\item Why does Clarisse always carry her camera?
\item What does the writer do to let off steam before an exam?
\item Who is “a good sport” in the text, and why?
\end{enumerate}

\begin{corrige}[Questions de compréhension]
1. \eng{He/She usually chills out with friends at a snack bar near the lycée, or listens to makossa at home.}

2. \eng{Because she is keen on photography: she took it up last year and never goes anywhere without her camera.}

3. \eng{He/She goes jogging to let off steam before an exam.}

4. \eng{The team captain is “a good sport”: although his team lost the match, he congratulated the winners.}
\end{corrige}

\courssub{Expressions avec « time »}

Le mot \cle{time} entre dans de nombreuses expressions liées aux loisirs. Apprenez-les comme des blocs.

\begin{propriete}[Les grandes expressions avec « time »]
\begin{itemize}
\item \eng{to have a good time} = passer un bon moment, bien s'amuser. \emph{Attention :} on dit \eng{to have a good time}, jamais \eng{to have good times} au sens de « s'amuser ». Équivalent plus familier : \eng{to have fun}.
\item \eng{to kill time} = tuer le temps, s'occuper en attendant quelque chose.
\item \eng{to take your time} = prendre son temps, ne pas se presser. Le possessif s'accorde avec le sujet : \eng{take my time}, \eng{take his time}.
\item \eng{me time} = du temps pour soi (registre familier, très courant à l'oral).
\item \eng{to spend time doing something} = passer du temps à faire quelque chose (structure : \eng{spend + time + V-ing}).
\end{itemize}
\end{propriete}

\begin{exemplebox}[Exemples traduits]
\eng{We had a good time at the party last night.} « Nous avons passé un très bon moment à la fête hier soir. »

\eng{I read magazines to kill time at the bus station.} « Je lis des magazines pour tuer le temps à la gare routière. »

\eng{Take your time, the film starts at eight.} « Prends ton temps, le film commence à huit heures. »

\eng{On Sundays, I need some me time.} « Le dimanche, j'ai besoin de temps pour moi. »

\eng{She spends her free time painting.} « Elle passe son temps libre à peindre. »
\end{exemplebox}

\courssub{Expressions d'engagement : dire qu'on aime une activité}

L'anglais dispose d'une échelle riche pour exprimer le goût. Toutes ces expressions se construisent avec \textbf{to be + expression + V-ing} (ou + nom).

\begin{center}
\begin{tabularx}{\linewidth}{|l|X|l|}
\hline
\rowcolor{popblueL} \textbf{Expression} & \textbf{Exemple anglais} & \textbf{Français} \\
\hline
to be fond of + V-ing & I'm fond of reading. & J'aime (beaucoup) lire. \\
\hline
to be keen on + V-ing & I'm keen on hiking. & J'aime beaucoup la randonnée. \\
\hline
to be a fan of + nom & She's a fan of makossa. & Elle est fan de makossa. \\
\hline
to be crazy about + V-ing & He's crazy about dancing. & Il est fou de danse. \\
\hline
to be into + nom/V-ing & They're into video games. & Ils sont branchés jeux vidéo. \\
\hline
to take up + hobby & He took up photography last year. & Il s'est mis à la photo l'année dernière. \\
\hline
\end{tabularx}
\end{center}

\begin{propriete}[Degrés d'enthousiasme]
On peut classer ces expressions de la plus faible à la plus forte : \eng{to be fond of} (affection douce) \textrightarrow{} \eng{to be keen on} (enthousiasme net, très britannique) \textrightarrow{} \eng{to be a fan of} (admiration) \textrightarrow{} \eng{to be crazy about} (passion, registre familier). Le verbe \eng{to take up} n'exprime pas un goût mais un \emph{commencement} : se mettre à une nouvelle activité. Son contraire est \eng{to give up} (arrêter, abandonner). Attention aux prétérits irréguliers : \eng{take \textrightarrow{} took \textrightarrow{} taken} et \eng{give \textrightarrow{} gave \textrightarrow{} given}.
\end{propriete}

\begin{exemplebox}[Exemples traduits]
\eng{I'm keen on hiking.} « J'aime beaucoup la randonnée. »

\eng{He took up photography last year.} « Il s'est mis à la photo l'année dernière. »

\eng{My sister is crazy about cooking shows.} « Ma sœur est folle des émissions de cuisine. »

\eng{Are you into reggae music?} « Tu es branché musique reggae ? »

\eng{She gave up tennis because of her studies.} « Elle a abandonné le tennis à cause de ses études. »
\end{exemplebox}

\courssub{Expressions de détente}

\begin{propriete}[Quatre verbes pour dire « se détendre »]
\begin{itemize}
\item \eng{to unwind} = décompresser après un effort (image : se « dérouler » comme un ressort). Registre neutre. Prétérit irrégulier : \eng{unwind \textrightarrow{} unwound \textrightarrow{} unwound}.
\item \eng{to chill out} = se relaxer, ne rien faire de fatigant. Registre familier, très fréquent chez les jeunes.
\item \eng{to let off steam} = se défouler, évacuer le stress accumulé (souvent par le sport ou une activité physique).
\item \eng{to recharge one's batteries} = récupérer son énergie. \emph{Attention :} \eng{one's} s'accorde avec le sujet : \eng{I recharge my batteries}, \eng{she recharges her batteries}.
\end{itemize}
\end{propriete}

\begin{exemplebox}[Exemples traduits]
\eng{After exams, students need to unwind.} « Après les examens, les élèves ont besoin de décompresser. »

\eng{Let's chill out at my place this afternoon.} « Viens te relaxer chez moi cet après-midi. »

\eng{He plays basketball to let off steam.} « Il joue au basket pour se défouler. »

\eng{A weekend in Kribi will recharge your batteries.} « Un week-end à Kribi te rechargera les batteries. »
\end{exemplebox}

\courssub{Expressions sportives}

\begin{itemize}
\item \eng{to keep fit} = rester en forme. \eng{She swims twice a week to keep fit.} « Elle nage deux fois par semaine pour rester en forme. »
\item \eng{to work out} = s'entraîner, faire de la musculation ou du sport en salle. \eng{He works out at the gym every morning.} « Il s'entraîne à la salle tous les matins. »
\item \eng{to be a good sport} = être fair-play, savoir perdre avec le sourire. \eng{Even when he loses, he's a good sport.} « Même quand il perd, il est fair-play. »
\end{itemize}

\courssub{Expressions de préférence négative}

\begin{propriete}[Dire poliment « je n'aime pas »]
\begin{itemize}
\item \eng{It's not my cup of tea} = ce n'est pas mon truc, ce n'est pas mon genre. Expression typiquement britannique, presque toujours employée à la forme négative.
\item \eng{I'm not into + nom/V-ing} = je ne suis pas branché..., ça ne me branche pas. Registre familier.
\item \eng{I'm not really a fan of...} = je ne suis pas vraiment fan de... (plus doux et plus poli que \eng{I hate}).
\end{itemize}
\end{propriete}

\begin{exemplebox}[Exemples traduits]
\eng{Reading is not my cup of tea.} « La lecture, ce n'est pas mon truc. »

\eng{I'm not into horror films.} « Les films d'horreur, ce n'est pas mon truc. »

\eng{Camping? Thanks, but it's not really my cup of tea.} « Le camping ? Merci, mais ce n'est pas vraiment mon genre. »
\end{exemplebox}

\begin{attention}[Erreurs fréquentes des francophones]
\begin{itemize}
\item Ne dites pas \eng{I am fond to play} : après \eng{fond of}, \eng{keen on}, \eng{crazy about}, \eng{into}, utilisez toujours la forme en \textbf{-ing} : \eng{I'm fond of playing chess.}
\item \eng{To kill time} ne veut pas dire « perdre son temps » au sens négatif : c'est simplement « s'occuper en attendant ». Pour « perdre son temps », dites \eng{to waste time}.
\item Ne traduisez pas « ce n'est pas mon truc » par une invention littérale : l'idiome anglais est \eng{it's not my cup of tea}.
\item \eng{To work out} signifie « s'entraîner », pas « travailler dehors ».
\item Ne confondez pas \eng{to take up} (commencer une activité) et \eng{to take} tout seul : \eng{I took up jogging} = « je me suis mis au jogging », et non « j'ai pris le jogging ».
\end{itemize}
\end{attention}

\courssub{Les expressions en action : mini-dialogues}

\begin{textebox}[Mini-dialogues]
\textbf{1.} — What do you do to unwind after school?\\
— I usually chill out with my friends or listen to music. And you?\\
— I let off steam playing football with the neighbourhood boys.

\textbf{2.} — Do you want to come fishing with us on Saturday?\\
— Thanks, but fishing is not my cup of tea. I'm more into indoor games like chess.

\textbf{3.} — Have you taken up any new hobby this year?\\
— Yes, I've taken up photography. I'm crazy about it! It really helps me recharge my batteries.

\textbf{4.} — Hurry up, the match is starting!\\
— Relax, take your time. We have twenty minutes to kill before kick-off.
\end{textebox}

\begin{popfigure}
\begin{center}
\begin{tikzpicture}[node distance=1.1cm]
\node[draw=popblue, fill=popblueL, rounded corners, thick, align=center, font=\bfseries] (q) {How do you relax?};
\node[draw=popgreen, fill=popgreenL, rounded corners, thick, align=center, below left=1.2cm and 2.6cm of q] (a) {I unwind\\ by reading.};
\node[draw=poporange, fill=poporangeL, rounded corners, thick, align=center, below=1.2cm of q] (b) {I let off steam\\ playing football.};
\node[draw=poppurple, fill=poppurpleL, rounded corners, thick, align=center, below right=1.2cm and 2.6cm of q] (c) {I chill out\\ with friends.};
\draw[-{Stealth}, thick, popgreen] (q) -- (a);
\draw[-{Stealth}, thick, poporange] (q) -- (b);
\draw[-{Stealth}, thick, poppurple] (q) -- (c);
\end{tikzpicture}
\end{center}
\legende{Organigramme : trois réponses types à la question \eng{How do you relax?}}
\end{popfigure}

\courssub{Méthode : réutiliser les expressions dans votre propre production}

\begin{methode}[Construire une phrase personnelle avec chaque idiome]
\begin{enumerate}
\item Choisissez une expression de la leçon, par exemple \eng{to be keen on}.
\item Identifiez sa structure : ici, \eng{to be + keen on + V-ing}.
\item Remplacez le sujet par \eng{I} et l'activité par une activité VRAIE pour vous : \eng{I'm keen on dancing.}
\item Ajoutez un détail (temps, lieu, personne) pour enrichir : \eng{I'm keen on dancing makossa with my cousins at family parties.}
\item Faites de même avec une expression négative : \eng{Watching television is not my cup of tea; I prefer reading novels.}
\item Mémorisez vos phrases personnelles : on retient beaucoup mieux un idiome quand il parle de soi. Au Baccalauréat, ces phrases toutes prêtes enrichiront votre composition sur le thème des loisirs.
\end{enumerate}
\end{methode}

\begin{exoresolu}[Complétez avec l'expression idiomatique correcte]
Énoncé. Complétez chaque phrase avec l'expression qui convient : \eng{kill time}, \eng{let off steam}, \eng{not my cup of tea}, \eng{took up}, \eng{recharge my batteries}.
\begin{enumerate}
\item I play the drums when I need to \ldots\ after a difficult day.
\item Horror films are \ldots\ ; I prefer comedies.
\item My cousin \ldots\ karate two months ago and he already loves it.
\item I always carry a novel to \ldots\ while waiting at the hospital.
\item A short holiday at the seaside will help me \ldots\ .
\end{enumerate}
\tcblower
Solution détaillée.
\begin{enumerate}
\item \eng{let off steam} : on se « défoule » après une journée difficile ; le tambour est une activité physique qui évacue le stress.
\item \eng{not my cup of tea} : la phrase oppose les films d'horreur aux comédies, donc il s'agit d'une préférence négative.
\item \eng{took up} : « il y a deux mois » indique le \emph{début} d'une nouvelle activité ; attention au prétérit irrégulier \eng{take \textrightarrow{} took}.
\item \eng{kill time} : lire un roman en attendant, c'est « tuer le temps ».
\item \eng{recharge my batteries} : les vacances permettent de récupérer son énergie ; on garde \eng{my} car le sujet est \eng{I}.
\end{enumerate}
\end{exoresolu}

\begin{exercice}[À vous de jouer]
\begin{enumerate}
\item Traduisez en anglais : a) « Je me suis mis au jogging l'année dernière. » b) « Les échecs, ce n'est pas mon truc. » c) « Elle fait du sport pour se défouler. » d) « Prends ton temps, nous ne sommes pas pressés. »
\item Écrivez trois phrases personnelles vraies avec \eng{I'm fond of...}, \eng{I'm crazy about...} et \eng{... is not my cup of tea}.
\item Rédigez un mini-dialogue de quatre répliques entre deux amis qui parlent de la façon dont ils se détendent le week-end, en utilisant au moins trois idiomes de la leçon.
\end{enumerate}
\end{exercice}

\begin{corrige}[Exercice : À vous de jouer]
1. a) \eng{I took up jogging last year.} b) \eng{Chess is not my cup of tea.} c) \eng{She does sport to let off steam.} (ou \eng{She plays sport to let off steam.}) d) \eng{Take your time, we are not in a hurry.}

2. Réponses libres. Modèles possibles : \eng{I'm fond of cooking with my mother.} / \eng{I'm crazy about Cameroonian football.} / \eng{Gardening is not my cup of tea.}

3. Modèle de dialogue : — \eng{What do you do to unwind at the weekend?} — \eng{I usually chill out with my friends in Douala. And you?} — \eng{I work out at the gym to keep fit, and then I have some me time at home.} — \eng{That sounds great! I should take up a sport too.}
\end{corrige}

\begin{aretenir}[L'essentiel de la section]
Une expression idiomatique s'apprend comme un bloc, sans traduction mot à mot. Pour les loisirs, retenez quatre familles : les expressions avec \eng{time} (\eng{have a good time}, \eng{kill time}, \eng{take your time}, \eng{me time}), les expressions d'engagement (\eng{be keen on}, \eng{be fond of}, \eng{be crazy about}, \eng{take up}), les expressions de détente (\eng{unwind}, \eng{chill out}, \eng{let off steam}, \eng{recharge one's batteries}) et les préférences négatives (\eng{not my cup of tea}, \eng{I'm not into...}). Après \eng{of}, \eng{on}, \eng{about} et \eng{into}, le verbe se met toujours en \textbf{-ing}. Enfin, les possessifs s'accordent avec le sujet : \eng{recharge my batteries}, \eng{take your time}.
\end{aretenir}

\coursec{Faux amis et pièges des francophones}

Dans la section précédente, nous avons découvert les expressions idiomatiques et le langage courant des loisirs (\eng{to kill time}, \eng{to have a blast}, \eng{to take up a hobby}...). Mais il reste un danger invisible : les mots anglais qui \emph{ressemblent} au français. Ces mots « amis » en apparence sont en réalité des traîtres. Un élève de Terminale qui dit \eng{I assisted to the match} croit dire « j'ai assisté au match » — il vient en fait de produire une phrase doublement incorrecte. Cette section vous apprend à repérer et à neutraliser ces pièges, un par un.

\courssub{Qu'est-ce qu'un faux ami ?}

\begin{definition}[Faux ami (\eng{false friend} ou \eng{false cognate})]
Un \cle{faux ami} est un mot anglais qui ressemble par la forme à un mot français (même origine latine le plus souvent) mais dont le \emph{sens} ou l'\emph{emploi} diffère. Exemple : \eng{actually} ne signifie pas « actuellement » mais « en réalité » ; « actuellement » se dit \eng{currently} ou \eng{at present}.
\end{definition}

Le danger est maximal dans le domaine des loisirs, car le français et l'anglais y partagent beaucoup de mots d'origine latine : \emph{assister}, \emph{distraire}, \emph{amateur}, \emph{éventuellement}... Chacun cache un piège. Observons d'abord les erreurs, puis énonçons les règles.

\courssub{To assist : le faux ami n° 1}

\begin{attention}[Erreur fréquente : \eng{I assisted to the match}]
La phrase \eng{I assisted to the match} est \textbf{doublement fausse} :
\begin{enumerate}
\item \eng{to assist} ne signifie pas « assister à » (être présent) mais « \cle{aider} » ;
\item même dans son sens correct, \eng{assist} se construit \emph{sans} préposition devant la personne aidée : \eng{to assist somebody}.
\end{enumerate}
Pour dire « j'ai assisté au match », on dit : \eng{I attended the match} ou \eng{I watched the match} (si l'on insiste sur le fait de regarder) ou encore \eng{I was at the match}.
\end{attention}

\begin{propriete}[Assist / attend / watch : trois verbes à ne pas confondre]
\begin{itemize}
\item \eng{to attend + noun} (sans préposition) = assister à, être présent à un événement : \eng{to attend a concert, a match, a meeting, a wedding}. C'est un verbe plutôt soutenu.
\item \eng{to watch + noun} = regarder (un spectacle, un match, un film) : plus courant à l'oral.
\item \eng{to assist + somebody (in doing something)} = aider quelqu'un : sens officiel ou formel.
\end{itemize}
\end{propriete}

\begin{exemplebox}[Assist, attend, watch en contexte]
\eng{Thousands of fans attended the concert in Yaoundé.} « Des milliers de fans ont assisté au concert à Yaoundé. »

\eng{We watched the match on television last night.} « Nous avons regardé le match à la télévision hier soir. »

\eng{A nurse assisted the injured player on the pitch.} « Une infirmière a aidé le joueur blessé sur le terrain. »

\eng{She assisted the coach in training the young players.} « Elle a aidé l'entraîneur à former les jeunes joueurs. »
\end{exemplebox}

\courssub{To make sport : le calque interdit}

En français, on « fait du sport ». Le verbe \eng{to make} ressemble à « faire », donc l'élève francophone produit naturellement \eng{I make sport every weekend}. C'est un \cle{calque} : une traduction mot à mot qui ne correspond à aucun usage anglais.

\begin{propriete}[Sport : les bons verbes]
\begin{itemize}
\item \eng{to do sport} (britannique) ou \eng{to play sports} (américain) = faire du sport en général.
\item \eng{to play + sport with ball or opponent} : \eng{play football, play tennis, play chess}.
\item \eng{to do + individual or martial activity} : \eng{do athletics, do judo, do gymnastics, do yoga}.
\item \eng{to go + verb in -ing} : \eng{go swimming, go running, go cycling, go fishing}.
\end{itemize}
\eng{To make} ne se combine jamais avec \eng{sport}. On peut dire \eng{to make progress}, \eng{to make an effort}, mais pas \eng{to make sport}.
\end{propriete}

\begin{exemplebox}[Phrases correctes]
\eng{I do sport twice a week to keep fit.} « Je fais du sport deux fois par semaine pour rester en forme. »

\eng{My brother plays football for the school team.} « Mon frère joue au football dans l'équipe de l'école. »

\eng{We go swimming at the municipal pool every Saturday.} « Nous allons nager à la piscine municipale chaque samedi. »
\end{exemplebox}

\courssub{« Mes loisirs » ne se traduit pas toujours par \eng{my leisure}}

Le mot \eng{leisure} existe en anglais, mais il est \emph{incomptable} et désigne le \emph{temps libre} en général, le concept abstrait. Il ne désigne presque jamais les activités concrètes que vous pratiquez.

\begin{propriete}[Leisure, hobbies, free-time activities]
\begin{itemize}
\item \eng{leisure} (incomptable) = le temps libre, le loisir en tant que notion : \eng{Leisure is important for health.} « Le loisir est important pour la santé. » On le trouve surtout dans les expressions \eng{leisure time}, \eng{leisure activities}, \eng{at your leisure} (« à votre convenance »).
\item « Mes loisirs » (mes activités préférées) = \eng{my hobbies} ou \eng{my free-time activities} ou \eng{what I do in my free time}.
\item « Un loisir » (une activité précise) = \eng{a hobby}, \eng{a pastime}, \eng{a leisure activity}.
\end{itemize}
\end{propriete}

\begin{exemplebox}[Comparaison]
\eng{My hobbies are reading and dancing.} « Mes loisirs sont la lecture et la danse. » (et non \eng{my leisures} — \eng{leisure} n'a pas de pluriel dans ce sens)

\eng{What do you do in your leisure time?} « Que fais-tu pendant ton temps libre ? »

\eng{Reading is a popular pastime in my family.} « La lecture est un loisir populaire dans ma famille. »
\end{exemplebox}

\courssub{Pastime, pass time, past : l'orthographe piège}

\begin{savaistu}[D'où vient \eng{pastime} ?]
Le mot \eng{pastime} vient de \eng{pass} + \eng{time} : littéralement, une façon de « passer le temps ». Il s'écrit \textbf{en un seul mot}, avec un seul « t » : \eng{pastime}, et non \eng{passtime} ni \eng{past time}. La prononciation est « passe-taïme ».
\end{savaistu}

\begin{attention}[Past / passed : ne pas confondre]
\begin{itemize}
\item \eng{passed} = participe passé et prétérit du verbe \eng{to pass} : \eng{We passed the time playing cards.} « Nous avons passé le temps à jouer aux cartes. »
\item \eng{past} = adjectif, préposition ou nom (« passé ») : \eng{in the past}, \eng{half past three}, \eng{the past week}.
\item On écrit donc \eng{to pass the time} (verbe) mais \eng{a pastime} (nom composé, un seul mot, un seul « t »).
\end{itemize}
\end{attention}

\courssub{Amateur : un mot partagé, deux usages différents}

Le mot \eng{amateur} existe en anglais (prononcé « a-me-teur »), mais attention à deux pièges :

\begin{propriete}[Amateur en anglais]
\begin{enumerate}
\item \eng{an amateur} = une personne qui pratique une activité sans être professionnelle : \eng{an amateur player} = « un joueur amateur ». Le mot est un adjectif ou un nom, jamais un loisir.
\item « Un hobby » ne se dit jamais \eng{an amateur}. On dit \eng{a hobby}. La phrase « la photographie est mon loisir » se traduit \eng{Photography is my hobby}, jamais \eng{photography is my amateur}.
\item Nuance culturelle : en anglais, \eng{amateur} peut avoir une nuance péjorative (« peu compétent ») dans \eng{an amateurish performance} = « une prestation d'amateur ».
\end{enumerate}
\end{propriete}

\begin{exemplebox}[Exemples]
\eng{He is an amateur photographer; he does not sell his pictures.} « C'est un photographe amateur ; il ne vend pas ses photos. »

\eng{The tournament is open to both professionals and amateurs.} « Le tournoi est ouvert aux professionnels comme aux amateurs. »
\end{exemplebox}

\courssub{To distract : « distraire » ne veut pas dire « divertir »}

\begin{attention}[To distract $\neq$ divertir]
En français, « distraire » a deux sens : \emph{divertir} (agréable) et \emph{détourner l'attention} (gênant). L'anglais \eng{to distract} n'a gardé que le sens négatif : \cle{détourner l'attention}, empêcher de se concentrer.
\begin{itemize}
\item « Ce film a diverti les enfants » = \eng{This film entertained the children} ou \eng{amused the children}.
\item \eng{The noise distracted me from my homework.} « Le bruit m'a distrait de mes devoirs » (sens négatif, correct ici).
\end{itemize}
Dire \eng{the film distracted the children} signifie que le film les a empêchés de se concentrer — l'inverse de ce que vous vouliez dire !
\end{attention}

\begin{exemplebox}[Entertain, amuse, distract]
\eng{This film entertained the children for two hours.} « Ce film a diverti les enfants pendant deux heures. »

\eng{The clown amused the whole audience.} « Le clown a amusé tout le public. »

\eng{Please do not distract the players during the match.} « S'il vous plaît, ne distrayez pas les joueurs pendant le match. »

\eng{Games are a good form of entertainment.} « Les jeux sont une bonne forme de divertissement. »
\end{exemplebox}

\courssub{Eventually : « éventuellement » ou « finalement » ?}

\begin{propriete}[Eventually = finalement ; possibly = éventuellement]
\eng{Eventually} signifie « \cle{finalement} », « à la fin », « après un certain temps ». Il n'exprime jamais l'hypothèse.
\begin{itemize}
\item \eng{After a long match, our team eventually won.} « Après un long match, notre équipe a finalement gagné. »
\item « Éventuellement » (peut-être, si nécessaire) = \eng{possibly}, \eng{if necessary}, \eng{if need be}, \eng{maybe}.
\item \eng{If it rains, we can possibly play indoors.} « S'il pleut, on peut éventuellement jouer en salle. »
\end{itemize}
\end{propriete}

\courssub{Tableau récapitulatif des pièges}

\begin{center}
\begin{tabularx}{\linewidth}{|X|X|X|}
\hline
\rowcolor{popblueL} Français (piège) & Correct English & Remarque \\
\hline
assister à un concert & \eng{to attend a concert} & \eng{to assist} = aider \\
\hline
faire du sport & \eng{to do sport} / \eng{to play sports} & jamais \eng{to make sport} \\
\hline
mes loisirs & \eng{my hobbies} / \eng{my free-time activities} & \eng{leisure} = le temps libre (incomptable) \\
\hline
un passe-temps & \eng{a pastime} & un seul mot, un seul « t » \\
\hline
un joueur amateur & \eng{an amateur player} & « un hobby » = \eng{a hobby} \\
\hline
divertir les enfants & \eng{to entertain} / \eng{amuse the children} & \eng{to distract} = détourner l'attention \\
\hline
éventuellement & \eng{possibly} / \eng{if necessary} & \eng{eventually} = finalement \\
\hline
\end{tabularx}
\end{center}

\begin{popfigure}
\begin{tikzpicture}[font=\small]
\node[draw=popink, fill=poppinkL, rounded corners, align=center, minimum width=5.2cm, minimum height=0.9cm] (t1) at (0,0) {\textbf{FRENCH TRAP}\\ \eng{I assisted to the match.}};
\node[draw=popgreen, fill=popgreenL, rounded corners, align=center, minimum width=5.2cm, minimum height=0.9cm] (c1) at (8,0) {\textbf{CORRECT ENGLISH}\\ \eng{I attended the match.}};
\draw[-{Stealth[length=3mm]}, very thick, poporange] (t1) -- (c1);
\node[draw=popink, fill=poppinkL, rounded corners, align=center, minimum width=5.2cm, minimum height=0.9cm] (t2) at (0,-1.6) {\eng{I make sport on Saturdays.}};
\node[draw=popgreen, fill=popgreenL, rounded corners, align=center, minimum width=5.2cm, minimum height=0.9cm] (c2) at (8,-1.6) {\eng{I do sport on Saturdays.}};
\draw[-{Stealth[length=3mm]}, very thick, poporange] (t2) -- (c2);
\node[draw=popink, fill=poppinkL, rounded corners, align=center, minimum width=5.2cm, minimum height=0.9cm] (t3) at (0,-3.2) {\eng{The show distracted us.}\\ (sens voulu : diverti)};
\node[draw=popgreen, fill=popgreenL, rounded corners, align=center, minimum width=5.2cm, minimum height=0.9cm] (c3) at (8,-3.2) {\eng{The show entertained us.}};
\draw[-{Stealth[length=3mm]}, very thick, poporange] (t3) -- (c3);
\node[draw=popink, fill=poppinkL, rounded corners, align=center, minimum width=5.2cm, minimum height=0.9cm] (t4) at (0,-4.8) {\eng{We will eventually play}\\ (sens voulu : éventuellement)};
\node[draw=popgreen, fill=popgreenL, rounded corners, align=center, minimum width=5.2cm, minimum height=0.9cm] (c4) at (8,-4.8) {\eng{We will possibly play.}\\ \eng{if necessary}};
\draw[-{Stealth[length=3mm]}, very thick, poporange] (t4) -- (c4);
\end{tikzpicture}
\legende{Tableau de correction : à gauche, la phrase piégée calquée du français ; à droite, la version anglaise correcte.}
\end{popfigure}

\begin{methode}[Comment éviter les calques : la méthode de la phrase modèle]
Avant d'employer un verbe « transparent » (qui ressemble au français) comme \eng{assist}, \eng{distract} ou \eng{eventually}, suivez ces quatre étapes :
\begin{enumerate}
\item \textbf{Méfiance automatique} : si le mot anglais ressemble au français, considérez-le comme suspect jusqu'à preuve du contraire.
\item \textbf{Vérification} : cherchez le mot dans une phrase modèle complète (dictionnaire bilingue ou manuel), pas seulement sa traduction isolée.
\item \textbf{Test de substitution} : remplacez le mot par son équivalent supposé dans la phrase anglaise. \eng{I assisted the match} $\rightarrow$ « j'ai aidé le match » ? Absurde $\rightarrow$ le verbe est faux.
\item \textbf{Mémorisation par paires} : apprenez le faux ami et son correctif ensemble : \eng{assist} = aider / \eng{attend} = assister à.
\end{enumerate}
\end{methode}

\begin{exoresolu}[Corrigez les phrases piégées]
Énoncé — Chaque phrase contient un faux ami ou un calque. Corrigez-la.
\begin{enumerate}
\item \eng{I assisted to the football match last Sunday.}
\item \eng{My sister makes sport every morning.}
\item \eng{The magician distracted the children at the party.} (sens voulu : diverti)
\item \eng{We will eventually go swimming if the weather is fine.} (sens voulu : éventuellement)
\item \eng{Photography is my favourite amateur.}
\end{enumerate}
\tcblower
Solution détaillée :
\begin{enumerate}
\item \eng{I attended the football match last Sunday} (ou \eng{watched}). \eng{Assist} signifie « aider » et ne prend pas la préposition \eng{to}.
\item \eng{My sister does sport every morning} (ou \eng{plays sports}). Calque de « faire du sport » : \eng{make} est interdit avec \eng{sport}.
\item \eng{The magician entertained the children at the party} (ou \eng{amused}). \eng{Distract} = détourner l'attention, sens négatif.
\item \eng{We will possibly go swimming if the weather is fine.} \eng{Eventually} = finalement ; l'hypothèse se rend par \eng{possibly} ou \eng{if necessary}.
\item \eng{Photography is my favourite hobby.} \eng{An amateur} désigne une personne non professionnelle, jamais un loisir.
\end{enumerate}
\end{exoresolu}

\begin{exercice}[À vous de jouer]
Traduisez en anglais en évitant les faux amis :
\begin{enumerate}
\item Des milliers de spectateurs ont assisté au concert à Douala.
\item Ce jeu vidéo a diverti toute la famille.
\item Je fais du sport pour passer le temps.
\item Nous irons éventuellement à la pêche demain, si nécessaire.
\item Mon oncle est un joueur d'échecs amateur ; c'est son passe-temps préféré.
\end{enumerate}
\end{exercice}

\begin{corrige}[Exercice « À vous de jouer »]
\begin{enumerate}
\item \eng{Thousands of spectators attended the concert in Douala.}
\item \eng{This video game entertained the whole family.}
\item \eng{I do sport to pass the time.}
\item \eng{We will possibly go fishing tomorrow, if necessary.}
\item \eng{My uncle is an amateur chess player; it is his favourite pastime.}
\end{enumerate}
\end{corrige}

\begin{aretenir}[L'essentiel de la section]
\begin{itemize}
\item \eng{To assist} = aider ; « assister à » = \eng{to attend} (ou \eng{to watch}).
\item « Faire du sport » = \eng{to do sport / to play sports}, jamais \eng{make sport}.
\item « Mes loisirs » = \eng{my hobbies} ; \eng{leisure} est incomptable et abstrait.
\item \eng{A pastime} s'écrit en un seul mot ; \eng{to pass the time} avec le verbe \eng{pass} (\eng{passed}, pas \eng{past}).
\item \eng{To distract} = détourner l'attention ; « divertir » = \eng{to entertain / to amuse}.
\item \eng{Eventually} = finalement ; « éventuellement » = \eng{possibly / if necessary}.
\item Réflexe d'or : tout mot anglais qui ressemble au français doit être vérifié dans une phrase complète avant emploi.
\end{itemize}
\end{aretenir}

\coursec{Familles de mots, prononciation et orthographe}

Après avoir débusqué les faux amis et les calques du français, nous allons maintenant apprendre à \emph{faire fructifier} notre vocabulaire. En anglais comme en français, un même mot-racine donne naissance à toute une famille : verbe, nom, adjectifs. Connaître ces familles, c'est multiplier par quatre ou cinq son stock de mots sans effort supplémentaire — et c'est aussi savoir prononcer et orthographier correctement ces mots, deux points où le Bac ne pardonne pas.

\courssub{Pourquoi travailler par familles de mots ?}

Observez d'abord ces quatre phrases, toutes construites autour de la même racine :

\eng{We relaxed by the pool after the match.} « Nous nous sommes détendus au bord de la piscine après le match. »

\eng{Swimming is a very relaxing activity.} « La natation est une activité très relaxante. »

\eng{I felt completely relaxed.} « Je me sentais complètement détendu. »

\eng{Yoga is my favourite form of relaxation.} « Le yoga est ma forme de détente préférée. »

Vous l'avez remarqué : un seul mot de base, \eng{relax}, produit un verbe, deux adjectifs et un nom. C'est exactement le mécanisme du français \emph{détendre / détendant / détendu / détente}. Retenir les familles, c'est donc apprendre l'anglais de manière économique et élégante.

\begin{popfigure}
\begin{tikzpicture}[grow=down, level distance=2.1cm, sibling distance=3.4cm,
  every node/.style={draw=popblue, fill=popblueL, rounded corners=2pt, align=center, font=\small, text=popdark},
  edge from parent/.style={draw=popblue, thick}]
\node[fill=popblue, text=white, font=\bfseries, align=center]{RELAX\\ (racine)}
  child { node[align=center]{to relax\\ \emph{(verbe)}} }
  child { node[align=center]{relaxation\\ \emph{(nom)}} }
  child { node[align=center]{relaxing\\ \emph{(adj. activité)}} }
  child { node[align=center]{relaxed\\ \emph{(adj. ressenti)}} };
\node[draw=none, fill=none, font=\footnotesize\itshape, text=popmuted] at (-5.1,-3.1) {We relaxed at home.};
\node[draw=none, fill=none, font=\footnotesize\itshape, text=popmuted] at (-1.7,-3.1) {a moment of relaxation};
\node[draw=none, fill=none, font=\footnotesize\itshape, text=popmuted] at (1.7,-3.1) {a relaxing evening};
\node[draw=none, fill=none, font=\footnotesize\itshape, text=popmuted] at (5.1,-3.1) {I feel relaxed.};
\end{tikzpicture}
\legende{Arbre de famille : la racine RELAX et ses quatre dérivés, avec un exemple par branche.}
\end{popfigure}

\courssub{La famille de RELAX}

\begin{definition}[La famille de RELAX]
\begin{itemize}
\item \cle{to relax} (verbe) : se détendre, se relaxer. \eng{I relax by listening to music.} « Je me détends en écoutant de la musique. »
\item \cle{relaxation} (nom) : la détente, la relaxation. \eng{Reading is a form of relaxation.} « La lecture est une forme de détente. »
\item \cle{relaxing} (adjectif) : relaxant, reposant — qualifie \textbf{l'activité ou la chose}. \eng{a relaxing holiday} « des vacances reposantes ».
\item \cle{relaxed} (adjectif) : détendu, relaxé — qualifie \textbf{la personne et son ressenti}. \eng{a relaxed smile} « un sourire détendu ».
\end{itemize}
\end{definition}

\courssub{La règle d'or : -ing ou -ed ?}

C'est l'une des questions les plus rentables du Bac. Observez la paire suivante :

\eng{Fishing is very relaxing.} « La pêche est très relaxante. »

\eng{I felt relaxed after the game.} « Je me sentais détendu après le match. »

\begin{propriete}[Adjectifs en -ing et en -ed]
\begin{itemize}
\item L'adjectif en \cle{-ing} décrit \textbf{la chose, l'activité ou la situation} qui provoque un sentiment : \eng{a relaxing evening} « une soirée relaxante », \eng{an amusing story} « une histoire amusante », \eng{a boring film} « un film ennuyeux ».
\item L'adjectif en \cle{-ed} décrit \textbf{le ressenti de la personne} : \eng{I feel relaxed} « je me sens détendu », \eng{we were amused} « nous étions amusés », \eng{she is bored} « elle s'ennuie ».
\item Mémo : \textbf{-ing = la cause, -ed = l'effet}. La chose est \emph{-ing} ; la personne est \emph{-ed}.
\end{itemize}
\end{propriete}

\begin{exemplebox}[La paire -ing / -ed en action]
\eng{The documentary was fascinating, so the students were fascinated.} « Le documentaire était fascinant, donc les élèves étaient fascinés. »

\eng{This game is tiring; I am tired.} « Ce jeu est fatigant ; je suis fatigué. »

\eng{What an exciting match! The supporters were so excited.} « Quel match passionnant ! Les supporters étaient tellement excités. »
\end{exemplebox}

\begin{attention}[Piège : « je suis intéressant » ou « je suis intéressé » ?]
Ne dites jamais \eng{I am interesting} pour « je suis intéressé » : cela signifie « je suis une personne intéressante » ! De même, \eng{I am boring} veut dire « je suis ennuyeux » (les autres s'ennuient avec moi), alors que \eng{I am bored} signifie « je m'ennuie ». Au Bac, cette confusion coûte des points à chaque session.
\end{attention}

\courssub{Les familles d'AMUSE et d'ENTERTAIN}

\begin{definition}[La famille d'AMUSE]
\begin{itemize}
\item \cle{to amuse} (verbe) : amuser, divertir. \eng{The clown amused the children.} « Le clown a amusé les enfants. »
\item \cle{amusement} (nom) : l'amusement. \eng{an amusement park} « un parc d'attractions ».
\item \cle{amusing} (adjectif, la chose) : amusant. \eng{an amusing joke} « une blague amusante ».
\item \cle{amused} (adjectif, la personne) : amusé. \eng{She was amused by the story.} « Elle était amusée par l'histoire. »
\end{itemize}
\end{definition}

\begin{definition}[La famille d'ENTERTAIN]
\begin{itemize}
\item \cle{to entertain} (verbe) : divertir, recevoir (des invités). \eng{The band entertained the crowd.} « Le groupe a diverti la foule. »
\item \cle{entertainment} (nom) : le divertissement, les spectacles. \eng{the entertainment industry} « l'industrie du divertissement ».
\item \cle{entertaining} (adjectif) : divertissant. \eng{a very entertaining show} « un spectacle très divertissant ».
\item \cle{an entertainer} (nom de métier) : un artiste, un animateur. \eng{He works as an entertainer in hotels.} « Il travaille comme animateur dans des hôtels. »
\end{itemize}
\end{definition}

Notez la formation du nom de métier avec le suffixe \cle{-er} (comme \eng{play} → \eng{player}, \eng{sing} → \eng{singer}, \eng{teach} → \eng{teacher}) : c'est le suffixe le plus productif de l'anglais pour désigner « celui qui fait l'action ».

\courssub{Les familles de RECREATE et de PLAY}

\begin{definition}[La famille de RECREATE]
\begin{itemize}
\item \cle{recreation} (nom) : les loisirs, la récréation. \eng{Recreation is important for health.} « Les loisirs sont importants pour la santé. »
\item \cle{recreational} (adjectif) : récréatif, de loisirs — c'est le mot-clé de notre leçon : \eng{recreational activities} « les activités récréatives / de loisirs ». \eng{The school organises recreational activities every Friday.} « L'école organise des activités récréatives chaque vendredi. »
\end{itemize}
\end{definition}

\begin{attention}[Attention à la prononciation de \eng{recreation}]
\eng{Recreation} (loisirs) se prononce « rè-kri-é-cheune » avec un \emph{e} bref au début. Il ne faut pas le confondre avec \eng{re-creation} (« re-création », le fait de créer à nouveau), qui se prononce « ri-kri-é-cheune ». À l'écrit, le contexte lève l'ambiguïté ; à l'oral, soyez précis.
\end{attention}

\begin{definition}[La famille de PLAY]
\begin{itemize}
\item \cle{to play} (verbe) : jouer. \eng{to play football / the guitar / a role} « jouer au football / de la guitare / un rôle ».
\item \cle{a player} (nom) : un joueur. \eng{Eto'o was a great player.} « Eto'o était un grand joueur. »
\item \cle{playful} (adjectif) : joueur, espiègle. \eng{a playful puppy} « un chiot espiègle ».
\item \cle{a playground} (nom) : une cour de récréation, un terrain de jeux. \eng{The children are in the playground.} « Les enfants sont dans la cour de récréation. »
\item \cle{a play} (nom) : une pièce de théâtre. \eng{We saw a play at the cultural centre.} « Nous avons vu une pièce au centre culturel. »
\end{itemize}
\end{definition}

\begin{center}
\begin{tabularx}{\linewidth}{|l|l|X|X|}
\hline
\rowcolor{popblueL} \textbf{Racine} & \textbf{Nature} & \textbf{English} & \textbf{Français} \\
\hline
RELAX & verbe & to relax & se détendre \\
\hline
RELAX & nom & relaxation & la détente \\
\hline
RELAX & adj. (chose) & relaxing & relaxant \\
\hline
RELAX & adj. (personne) & relaxed & détendu \\
\hline
AMUSE & verbe & to amuse & amuser \\
\hline
AMUSE & nom & amusement & l'amusement \\
\hline
AMUSE & adj. (chose / personne) & amusing / amused & amusant / amusé \\
\hline
ENTERTAIN & verbe & to entertain & divertir \\
\hline
ENTERTAIN & nom & entertainment & le divertissement \\
\hline
ENTERTAIN & adj. & entertaining & divertissant \\
\hline
ENTERTAIN & métier & an entertainer & un animateur, un artiste \\
\hline
RECREATE & nom & recreation & les loisirs \\
\hline
RECREATE & adj. & recreational & récréatif \\
\hline
PLAY & métier & a player & un joueur \\
\hline
PLAY & adj. & playful & espiègle \\
\hline
PLAY & lieu & a playground & une cour de récréation \\
\hline
PLAY & œuvre & a play & une pièce de théâtre \\
\hline
\end{tabularx}
\end{center}

\courssub{Prononciation : les mots-pièges du thème}

\begin{attention}[Prononciations à mémoriser]
\begin{itemize}
\item \cle{leisure} se prononce « \textbf{lè-jeur} » en anglais britannique (et « li-jeur » en américain) — jamais « lé-jure ». \eng{leisure activities} « les activités de loisirs ».
\item \cle{choir} se prononce « \textbf{kouaïeur} » : le \emph{ch} se dit \emph{k} et le \emph{oi} se dit « ouaï ». \eng{She sings in the church choir.} « Elle chante dans la chorale de l'église. »
\item \cle{hobby} se prononce « \textbf{ho-bi} » : le \emph{h} anglais est aspiré (un souffle léger, comme dans \emph{house}) et le \emph{o} est bref, proche du \emph{o} français de « note ». Ne dites ni « o-bi » sans souffle, ni « ro-bi ».
\item \cle{amateur} se prononce « \textbf{a-me-teu} » en trois syllabes, avec l'accent sur la première : \textbf{A}-me-teu. Ne copiez pas la prononciation française « a-ma-teur ».
\item \cle{recreational} se prononce « rè-kri-é-cheu-neul », accent sur « é ».
\end{itemize}
\end{attention}

\courssub{Orthographe : quatre réflexes à acquérir}

\begin{propriete}[Orthographe du vocabulaire des loisirs]
\begin{enumerate}
\item \cle{pastime} s'écrit \textbf{en un seul mot} : \eng{Reading is my favourite pastime.} « La lecture est mon passe-temps préféré. »
\item \cle{free time} s'écrit \textbf{en deux mots} : \eng{What do you do in your free time?} « Que fais-tu pendant ton temps libre ? »
\item \cle{hobby} fait son pluriel en \textbf{hobbies} : consonne + \emph{y} → \emph{ies} (comme \eng{story} → \eng{stories}). \eng{My hobbies are swimming and dancing.} « Mes passe-temps sont la natation et la danse. »
\item En anglais britannique, le verbe s'écrit \cle{to practise} (avec \textbf{S}) et le nom \cle{practice} (avec \textbf{C}) : \eng{I practise the piano every day; practice makes perfect.} « Je m'entraîne au piano chaque jour ; c'est en forgeant qu'on devient forgeron. »
\end{enumerate}
\end{propriete}

\begin{savaistu}[Practise ou practice ?]
En anglais britannique, la distinction est nette : \eng{to practise} (verbe, avec S) et \eng{practice} (nom, avec C) — exactement comme \eng{to advise} / \eng{advice}. En anglais américain, en revanche, tout s'écrit avec C : \eng{to practice} et \eng{practice}. Astuce de mémorisation pour le britannique : \textbf{S} comme \emph{verbe d'action} (\emph{s}e entraîner), \textbf{C} comme \emph{chose} (le nom). Au Bac, les deux normes sont acceptées, mais soyez cohérent dans toute votre copie.
\end{savaistu}

\courssub{Enrichir son style au Bac : nominaliser}

\begin{methode}[Remplacer un verbe simple par un nom dérivé]
Pour varier le style et montrer votre maîtrise du lexique au correcteur, transformez une phrase verbale banale en construction nominale :
\begin{enumerate}
\item Repérez le verbe simple de votre phrase : \eng{We relaxed after the exam.}
\item Identifiez le nom de la même famille : \eng{relaxation}.
\item Reconstruisez la phrase autour du nom : \eng{After the exam, it was a moment of pure relaxation.}
\item Comparez : la seconde version est plus riche et plus « Bac ». Autres exemples : \eng{They entertained us} → \eng{They provided us with great entertainment} ; \eng{The show amused the children} → \eng{The show was a source of amusement for the children}.
\end{enumerate}
\end{methode}

\begin{exoresolu}[Familles de mots : complétez avec le bon dérivé]
Complétez chaque phrase avec la forme correcte du mot entre parenthèses.

1. Chess is a very \ldots\ game. (RELAX)

2. The children were \ldots\ by the magician's tricks. (AMUSE)

3. Our school offers many \ldots\ activities, such as drama and football. (RECREATE)

4. The \ldots\ sang beautifully at the Christmas concert. (CHOIR — donnez le sens et la prononciation)

5. She felt completely \ldots\ after her yoga session. (RELAX)

\tcblower
\textbf{Solution détaillée.}

1. \eng{relaxing} : on décrit \textbf{la chose} (le jeu d'échecs), donc adjectif en \emph{-ing}. « Les échecs sont un jeu très relaxant. »

2. \eng{amused} : on décrit \textbf{le ressenti des enfants}, donc adjectif en \emph{-ed}. « Les enfants étaient amusés par les tours du magicien. »

3. \eng{recreational} : il faut l'adjectif dérivé du nom \eng{recreation} pour qualifier \eng{activities}. « Notre école propose de nombreuses activités récréatives, comme le théâtre et le football. »

4. \eng{choir} signifie « la chorale » et se prononce « kouaïeur ». « La chorale a magnifiquement chanté au concert de Noël. »

5. \eng{relaxed} : \eng{she felt} indique un \textbf{ressenti}, donc \emph{-ed}. « Elle se sentait complètement détendue après sa séance de yoga. »
\end{exoresolu}

\begin{exercice}[À vous de jouer]
1. Traduisez : « Ce film est ennuyeux, les spectateurs s'ennuient. » (deux adjectifs différents !)

2. Trouvez le nom de métier dérivé de \eng{entertain}, puis employez-le dans une phrase.

3. Corrigez les erreurs : \eng{My hobbys are reading and to practise sport in my freetime.}

4. Classez dans deux colonnes « la chose » et « la personne » : \eng{amusing, relaxed, entertaining, amused, relaxing, bored}.
\end{exercice}

\begin{corrige}[Exercice « À vous de jouer »]
1. \eng{This film is boring; the spectators are bored.} Le film (chose) est \emph{boring} ; les spectateurs (personnes) sont \emph{bored}.

2. \eng{an entertainer}. Exemple : \eng{The entertainer made the whole village laugh during the festival.} « L'artiste a fait rire tout le village pendant le festival. »

3. \eng{My hobbies are reading and practising sport in my free time.} Trois corrections : \eng{hobby} → \eng{hobbies} (y → ies) ; \eng{to practise} → \eng{practising} (après \eng{are}, parallélisme avec \eng{reading}) ; \eng{freetime} → \eng{free time} (deux mots).

4. La chose : \eng{amusing, entertaining, relaxing}. La personne : \eng{relaxed, amused, bored}.
\end{corrige}

\begin{aretenir}[L'essentiel de la section]
\begin{itemize}
\item Une racine = une famille : verbe, nom, adjectifs (\eng{relax, relaxation, relaxing, relaxed}).
\item \textbf{-ing} décrit l'activité (la cause), \textbf{-ed} décrit le ressenti (l'effet).
\item Prononciations-pièges : \eng{leisure} « lè-jeur », \eng{choir} « kouaïeur », \eng{hobby} « ho-bi » (h aspiré), \eng{amateur} « A-me-teu ».
\item Orthographe : \eng{pastime} en un mot, \eng{free time} en deux mots, \eng{hobby} → \eng{hobbies}, \eng{to practise} (verbe GB) / \eng{practice} (nom GB).
\item Au Bac, nominalisez pour enrichir votre style : \eng{we relaxed} → \eng{a moment of relaxation}.
\end{itemize}
\end{aretenir}

\coursec{Emploi en contexte : dialogues et court texte}

Après avoir exploré les familles de mots, la prononciation et l'orthographe du lexique des loisirs, nous passons maintenant à l'étape cruciale : \cle{l'emploi en contexte}. Connaître des listes de mots ne suffit pas ; il faut savoir les assembler dans des phrases correctes, naturelles et adaptées à la situation de communication. Cette section vous propose deux dialogues modèles, un texte de lecture authentique ancré dans le contexte camerounais, des activités de compréhension et de classement lexical, pour finir par une production écrite guidée calquée sur les exigences du Baccalauréat.

\courssub{Dialogue modèle 1 : Évoquer son week-end (Past Simple \& Vocabulaire loisirs)}

Observons comment deux élèves, \eng{Paul} et \eng{Mary}, échangent sur leurs activités du week-end. Remarquez l'usage systématique du \cle{Past Simple} (prétérit) pour des actions terminées dans le passé, et la richesse du vocabulaire de loisirs (sport, cérémonie, plein air).

\begin{textebox}[Dialogue : "How was your weekend?"]
\textbf{Paul:} Hey Mary, how was your weekend? Did you do anything special?

\textbf{Mary:} Hi Paul! It was great, thanks. I \textbf{went hiking} on Mount Fako with my cousins on Saturday. The view from the summit was breathtaking. We also \textbf{had a picnic} near the crater lake. It was exhausting but really refreshing.

\textbf{Paul:} Wow, that sounds amazing! I didn't go out much. I \textbf{attended a traditional wedding} in my village near Bafoussam. There was a lot of \textbf{dancing}, \textbf{drumming} and delicious food — \textit{ndolé} and plantains!

\textbf{Mary:} Nice! I love weddings. Did you \textbf{take part in the traditional dances}?

\textbf{Paul:} Of course! I even tried the \textbf{balafon} for a few minutes. On Sunday, I just \textbf{chilled out} at home and \textbf{binge-watched} a series.
\end{textebox}

\begin{propriete}[Structures clés pour raconter son week-end]
\begin{itemize}
    \item \textbf{Question type :} \eng{What did you do (last weekend / on Saturday)?} / \eng{How was your weekend?}
    \item \textbf{Réponse (action précise) :} \eng{I + Past Simple (verbe irrégulier ou régulier) + complément.}
    \begin{center}
    \begin{tabularx}{\linewidth}{|X|X|X|}\hline
    \rowcolor{popblueL} \textbf{Verbe (Base)} & \textbf{Past Simple} & \textbf{Exemple contextuel} \\ \hline
    \eng{go} & \eng{went} & \eng{I \textbf{went hiking} / \textbf{went to a wedding}.} \\ \hline
    \eng{have} & \eng{had} & \eng{We \textbf{had a picnic} / \textbf{had fun}.} \\ \hline
    \eng{attend} & \eng{attended} & \eng{I \textbf{attended a wedding} / \textbf{attended choir practice}.} \\ \hline
    \eng{take part} & \eng{took part} & \eng{Did you \textbf{take part in the dances}?} \\ \hline
    \eng{chill out} & \eng{chilled out} & \eng{I just \textbf{chilled out} at home.} \\ \hline
    \eng{binge-watch} & \eng{binge-watched} & \eng{We \textbf{binge-watched} a series.} \\ \hline
    \end{tabularx}
    \end{center}
    \item \textbf{Adverbes de temps fréquents :} \eng{last weekend, on Saturday, on Sunday, yesterday, two days ago}.
\end{itemize}
\end{propriete}

\begin{exemplebox}[Vocabulaire utile : Cérémonies et plein air (Contexte Cameroun)]
\begin{center}
\begin{tabularx}{\linewidth}{|X|X|X|}\hline
\rowcolor{popgreenL} \textbf{English} & \textbf{Français} & \textbf{Collocation / Contexte} \\ \hline
\eng{traditional wedding} & mariage traditionnel & \eng{attend a traditional wedding} \\ \hline
\eng{bride / groom} & mariée / marié & \eng{The bride wore a beautiful \textit{toghu}.} \\ \hline
\eng{dowry ceremony} & cérémonie de la dot & \eng{The dowry ceremony lasted all day.} \\ \hline
\eng{hiking / trekking} & randonnée (montagne) & \eng{go hiking on Mount Cameroon / Mount Fako} \\ \hline
\eng{summit / peak} & sommet & \eng{reach the summit} \\ \hline
\eng{breathtaking view} & vue à couper le souffle & \eng{enjoy a breathtaking view} \\ \hline
\eng{crater lake} & lac de cratère & \eng{Lake Dissoni / Lake Monoun} \\ \hline
\end{tabularx}
\end{center}
\end{exemplebox}

\begin{attention}[Piège classique : "Assister" $\neq$ \eng{Assist}]
Ne traduisez jamais « j'ai assisté à un mariage » par \eng{*I assisted a wedding}.
\begin{itemize}
    \item \eng{Assist} = aider (\eng{I assisted the old lady.} = J'ai aidé la personne âgée).
    \item \eng{Attend} = être présent à un événement (\eng{I attended a wedding.} = J'ai assisté à un mariage).
\end{itemize}
\cle{Attend} est le verbe correct pour les cours, réunions, mariages, concerts, messes.
\end{attention}

\courssub{Dialogue modèle 2 : Inviter, accepter et refuser poliment}

Inviter un ami à une activité de loisir mobilise des structures de politesse (\cle{Would you like to...}, \cle{Do you want to...}) et des réponses types. Observons \eng{John} et \eng{Grace}.

\begin{textebox}[Dialogue : "Would you like to join us?"]
\textbf{John:} Hi Grace! \textbf{Are you free this Saturday afternoon}? A few friends and I are going to the community centre to \textbf{play draughts} and \textbf{play cards}. \textbf{Would you like to join us}?

\textbf{Grace:} \textbf{That sounds like fun!} \textbf{I'd love to,} \textbf{but} I \textbf{have choir practice} at 3 p.m. \textbf{I'm afraid I can't make it.}

\textbf{John:} \textbf{No worries.} Maybe \textbf{another time}? We often play on Sundays too.

\textbf{Grace:} \textbf{Sure, why not?} \textbf{I'll let you know} if I'm free next Sunday. Thanks for asking!
\end{textebox}

\begin{propriete}[Formules types : Inviter $\rightarrow$ Accepter / Refuser]
\begin{center}
\begin{tabularx}{\linewidth}{|X|X|}\hline
\rowcolor{poporangeL} \textbf{Fonction} & \textbf{Expressions anglaises (registre courant)} \\ \hline
\textbf{Inviter} & \eng{Would you like to + BV ?} (Politesse standard) \\
& \eng{Do you want to + BV ?} (Plus familier/amical) \\
& \eng{Are you free on [jour] / this [afternoon]?} (Vérifier dispo) \\
& \eng{Why don't we + BV ?} / \eng{What about + V-ing ?} (Suggestion) \\ \hline
\textbf{Accepter} & \eng{I'd love to.} / \eng{I'd love to join you.} \\
& \eng{That sounds great / fun / awesome.} \\
& \eng{Sure! / Why not? / Count me in!} \\ \hline
\textbf{Refuser (poliment)} & \eng{I'd love to, \textbf{but}...} (Structure obligatoire) \\
& \eng{I'm afraid I can't.} / \eng{I'm sorry, I can't make it.} \\
& \eng{I have another commitment.} / \eng{I have to + BV.} (Raison) \\
& \eng{Maybe another time.} / \eng{Rain check?} (Report) \\ \hline
\end{tabularx}
\end{center}
\end{propriete}

\begin{exemplebox}[Exemples traduits (Modèles à mémoriser)]
\begin{itemize}
    \item \eng{Would you like to join us for a game of draughts?} = Voudrais-tu te joindre à nous pour une partie de dames ?
    \item \eng{I'd love to, but I have choir practice.} = J'aimerais bien, mais j'ai répétition de chorale.
    \item \eng{Are you free on Saturday evening?} = Es-tu libre samedi soir ?
    \item \eng{I'm afraid I can't make it.} = Je crains de ne pas pouvoir venir / Je ne pourrai pas venir.
    \item \eng{Count me in for the hike!} = Compte sur moi pour la rando ! (Très familier/enthousiaste)
\end{itemize}
\end{exemplebox}

\begin{attention}[Erreur majeure des francophones : "I very like..."]
\begin{itemize}
    \item \textbf{INCORRECT :} \eng{*I very like football.} / \eng{*I very enjoyed the party.}
    \item \textbf{CORRECT :} \eng{I \textbf{like football very much}.} / \eng{I \textbf{really like} football.} / \eng{I \textbf{enjoyed the party a lot}.}
    \item \textbf{Alternatives sophistiquées (Bac+) :} \eng{I'm \textbf{fond of} football.} / \eng{I'm \textbf{keen on} hiking.} / \eng{I \textbf{am passionate about} music.}
\end{itemize}
\cle{Règle :} \eng{Very} ne peut pas modifier un verbe directement. Il modifie un adjectif/adverbe (\eng{very good}). Pour intensifier un verbe : \eng{very much}, \eng{a lot}, \eng{really}, ou changer de verbe (\eng{love, enjoy, be fond of}).
\end{attention}

\courssub{Court texte de lecture : « Recreation in my village »}

Le texte suivant mobilise le lexique des cinq champs étudiés (Sport, Arts/Musique, Jeux traditionnels, Vie sociale/Famille, Plein air/Nature). Lisez-le attentivement, puis répondez aux questions.

\begin{textebox}[Recreation in my village (near Bafoussam)]
In my village near Bafoussam, recreation brings people together and strengthens our cultural identity. Every Sunday afternoon, the community field comes alive. Young men \textbf{play football} fiercely, cheered by supporters drinking \textit{foléré} juice, while older men sit under the mango trees to \textbf{play draughts} (songo'o) or \textbf{play cards}. 

Meanwhile, the women's \textbf{choir} rehearses hymns for the Sunday service, their voices blending with the sound of the \textbf{balafon} played by a griot. Children, free from school, prefer \textbf{traditional games} like \textbf{hide-and-seek} or \textbf{marbles} in the red dust. 

As the sun sets, families \textbf{gather} around a fire for \textbf{storytelling} sessions. Elders narrate folktales about the spider \textit{Kwaku Ananse} or the history of our clan. These \textbf{pastimes} are not just entertainment; they transmit values, teach history, and \textbf{strengthen social bonds} between generations. Weekends here are never boring; they are a celebration of community life.
\end{textebox}

\begin{definition}[Glossaire du texte (Mots difficiles)]
\begin{itemize}
    \item \eng{comes alive} (loc. v.) : s'anime, devient vivant.
    \item \eng{foléré} (n.) : jus de fleurs d'hibiscus (bissap), boisson locale populaire.
    \item \eng{draughts (UK) / checkers (US)} (n.) : jeu de dames. \eng{Songo'o} : nom local camerounais.
    \item \eng{rehearses} (v.) : répète (spectacle, chorale). \eng{Choir practice} = répétition de chorale.
    \item \eng{blending with} (part. prés.) : se mêlant à.
    \item \eng{griot} (n.) : conteur/musicien traditionnel ouest-africain.
    \item \eng{marbles} (n. pl.) : billes (jeu).
    \item \eng{folktales} (n. pl.) : contes populaires, légendes.
    \item \eng{pastimes} (n. pl.) : passe-temps, loisirs (plus formel que \eng{hobbies}).
    \item \eng{strengthen} (v.) : renforcer, consolider. \eng{Bond} (n.) : lien.
\end{itemize}
\end{definition}

\begin{exercice}[Compréhension du texte : « Recreation in my village »]
\textbf{A. True or False? Justify with a quote from the text.}
\begin{enumerate}
    \item Only young people play football on the community field.
    \item The women's choir sings with the accompaniment of a balafon.
    \item Children play modern video games in the dust.
    \item Storytelling happens in the morning before church.
    \item The narrator thinks recreational activities are useless.
\end{enumerate}

\textbf{B. Multiple Choice Questions (MCQ).}
\begin{enumerate}
    \item The expression "comes alive" (line 2) means:
    \begin{itemize}
        \item[a)] The field is dangerous.
        \item[b)] The field becomes active and busy.
        \item[c)] The field is repaired.
    \end{itemize}
    \item "Songo'o" is mentioned as a synonym for:
    \begin{itemize}
        \item[a)] Football.
        \item[b)] Draughts / Checkers.
        \item[c)] Marbles.
    \end{itemize}
    \item The folktales are told by:
    \begin{itemize}
        \item[a)] The children.
        \item[b)] The griot / Elders.
        \item[c)] The choir women.
    \end{itemize}
\end{enumerate}

\textbf{C. Short Answer Questions.}
\begin{enumerate}
    \item What do the older men do while the young men play football?
    \item Name two traditional games played by children.
    \item According to the last paragraph, what are the three functions of these pastimes?
\end{enumerate}
\end{exercice}

\begin{corrige}[Exercice : Compréhension « Recreation in my village »]
\textbf{A. True or False?}
\begin{enumerate}
    \item \textbf{False.} Quote: "\eng{Young men play football fiercely, cheered by supporters... while older men sit under the mango trees...}". Older men are present but play draughts/cards.
    \item \textbf{True.} Quote: "\eng{...their voices blending with the sound of the balafon played by a griot.}".
    \item \textbf{False.} Quote: "\eng{Children... prefer traditional games like hide-and-seek or marbles in the red dust.}".
    \item \textbf{False.} Quote: "\eng{As the sun sets, families gather around a fire for storytelling sessions.}". It happens in the evening.
    \item \textbf{False.} Quote: "\eng{These pastimes are not just entertainment; they transmit values, teach history, and strengthen social bonds...}". They are very useful.
\end{enumerate}

\textbf{B. MCQ.}
\begin{enumerate}
    \item \textbf{b)} The context (football, supporters, choir, games) shows activity.
    \item \textbf{b)} "\eng{...older men sit under the mango trees to play draughts (songo'o)...}".
    \item \textbf{b)} "\eng{Elders narrate folktales...}".
\end{enumerate}

\textbf{C. Short Answers.}
\begin{enumerate}
    \item They sit under mango trees to play draughts (songo'o) or cards.
    \item Hide-and-seek and marbles. (Also acceptable: traditional games like songo'o if considered for kids, but text separates "older men play draughts" and "children prefer traditional games like...").
    \item 1) Transmit values. 2) Teach history. 3) Strengthen social bonds between generations.
\end{enumerate}
\end{corrige}

\courssub{Relevé et classification du lexique dans le texte}

\begin{exercice}[Activité de classement : 10 mots du thème]
Relevez \textbf{10 mots ou expressions} du texte « Recreation in my village » appartenant au champ des loisirs. Classez-les dans le tableau ci-dessous selon leur famille lexicale (Champ). Un exemple est donné.
\begin{center}
\begin{tabularx}{\linewidth}{|X|X|X|}\hline
\rowcolor{poppurpleL} \textbf{Mot / Expression (English)} & \textbf{Traduction (Français)} & \textbf{Champ lexical (Catégorie)} \\ \hline
\eng{play football} & jouer au football & Sports / Jeux d'équipe \\ \hline
 &  &  \\ \hline
 &  &  \\ \hline
 &  &  \\ \hline
 &  &  \\ \hline
 &  &  \\ \hline
 &  &  \\ \hline
 &  &  \\ \hline
 &  &  \\ \hline
 &  &  \\ \hline
 &  &  \\ \hline
\end{tabularx}
\end{center}
\emph{Indication : Cherchez des verbes d'action (play, rehearse, gather, narrate, strengthen), des noms d'activités (choir, storytelling, pastimes, folktales, draughts, marbles, hide-and-seek) et des noms de lieux/acteurs (community field, griot, elders, supporters).}
\end{exercice}

\begin{corrige}[Exercice : Relevé et classification (Proposition de correction)]
\begin{center}
\begin{tabularx}{\linewidth}{|X|X|X|}\hline
\rowcolor{poppurpleL} \textbf{Mot / Expression (English)} & \textbf{Traduction (Français)} & \textbf{Champ lexical (Catégorie)} \\ \hline
\eng{play football} & jouer au football & Sports / Jeux d'équipe \\ \hline
\eng{play draughts (songo'o)} & jouer aux dames & Jeux de société / Traditionnels \\ \hline
\eng{play cards} & jouer aux cartes & Jeux de société \\ \hline
\eng{choir (rehearses)} & chorale (répète) & Arts / Musique \\ \hline
\eng{balafon} & balafon & Arts / Musique / Instruments \\ \hline
\eng{traditional games} & jeux traditionnels & Jeux traditionnels \\ \hline
\eng{hide-and-seek} & cache-cache & Jeux d'enfants / Traditionnels \\ \hline
\eng{marbles} & billes & Jeux d'enfants / Traditionnels \\ \hline
\eng{storytelling} & contes / narration & Vie sociale / Culture orale \\ \hline
\eng{pastimes} & passe-temps / loisirs & Vie sociale / Général \\ \hline
\eng{strengthen social bonds} & renforcer les liens sociaux & Vie sociale / Bienfaits \\ \hline
\eng{gather} & se rassembler / se réunir & Vie sociale / Action \\ \hline
\end{tabularx}
\end{center}
\emph{Note : D'autres réponses sont possibles (ex: \eng{folktales} $\rightarrow$ Culture orale, \eng{griot} $\rightarrow$ Arts/Musique, \eng{community field} $\rightarrow$ Lieu de sport). L'essentiel est la justification de la catégorie.}
\end{corrige}

\courssub{Production guidée : « My favourite recreational activity »}

Au Baccalauréat, l'expression écrite (Composition / Guided Writing) demande souvent de décrire une activité, un événement ou une habitude. Voici la méthode infaillible pour réussir un paragraphe de 8 à 10 lignes (environ 100-120 mots) sur ce thème.

\begin{methode}[Réussir le paragraphe « My favourite recreational activity » au Bac+]
\textbf{Sujet type :} "Write a paragraph about your favourite recreational activity. Say what it is, when and where you do it, who you do it with, and why you like it (benefits)."
\begin{enumerate}
    \item \textbf{Choisir une activité précise} que vous maîtrisez lexicalement (ex: \eng{playing football}, \eng{hiking}, \eng{singing in the choir}, \eng{playing draughts}, \eng{reading novels}, \eng{gardening}). Évitez les activités trop vagues (\eng{hanging out}).
    \item \textbf{Respecter le plan imposé (4 mouvements) :}
    \begin{itemize}
        \item \textbf{Introduction (Topic Sentence) :} Nommer l'activité + verbe d'opinion (\eng{My favourite recreational activity is... / I am keen on... / I really enjoy...}).
        \item \textbf{Description (When, Where, With whom) :} Fréquence (\eng{every weekend, twice a week}), Lieu (\eng{at the stadium, in the village, at the community centre}), Partenaires (\eng{with my friends, with my cousins, alone}).
        \item \textbf{Bénéfices (Why?) :} Physique (\eng{keep fit, stay healthy}), Mental (\eng{relax, relieve stress, unwind}), Social (\eng{strengthen bonds, meet new people, teamwork}), Culturel (\eng{preserve tradition, learn history}).
        \item \textbf{Conclusion (Closing sentence) :} Sentiment final / Projet (\eng{I will never stop... / It is essential to my well-being.}).
    \end{itemize}
    \item \textbf{Contraintes linguistiques (Score boosters) :}
    \begin{itemize}
        \item Utiliser \textbf{au moins 5 mots/expressions du lexique} de la leçon (champs + collocations).
        \item Utiliser \textbf{au moins 2 collocations correctes} (\eng{play football}, \eng{go hiking}, \eng{do yoga}, \eng{attend choir practice}, \eng{take part in}).
        \item Varier les \textbf{familles de mots} (ex: \eng{relax} (v.) $\rightarrow$ \eng{relaxing} (adj.) $\rightarrow$ \eng{relaxation} (n.) ; \eng{compete} $\rightarrow$ \eng{competition} $\rightarrow$ \eng{competitive}).
        \item Utiliser des \textbf{connecteurs logiques} : \eng{Firstly, Moreover, Besides, As a result, Finally, In conclusion}.
    \end{itemize}
\end{enumerate}
\end{methode}

\begin{popfigure}
\begin{tikzpicture}[
    node distance=1.2cm and 1.5cm,
    box/.style={draw=popblue, thick, rounded corners, fill=popblueL, text width=3.2cm, align=center, minimum height=1.8cm, font=\small},
    arrow/.style={-Stealth, thick, popblue},
    title/.style={font=\bfseries\small, text=popdark}
]
% Nodes
\node[box, title, align=center] (intro){Introduction\\ \textbf{Activity + Opinion}\\\eng{My favourite activity is...}\\\eng{I am fond of...}};
\node[box, title, right=of intro, align=center] (desc){Description\\ \textbf{When / Where / With whom}\\\eng{Every Saturday at the stadium}\\\eng{with my teammates}};
\node[box, title, right=of desc, align=center] (benef){Benefits\\ \textbf{Physical / Mental / Social}\\\eng{Keeps me fit, relieves stress,}\\\eng{strengthens bonds}};
\node[box, title, right=of benef, align=center] (conc){Conclusion\\ \textbf{Feeling / Future}\\\eng{It is vital for my balance.}\\\eng{I'll never quit.}};

% Arrows
\draw[arrow] (intro.east) -- (desc.west) node[midway, above, font=\tiny, popmuted] {Firstly / To begin with};
\draw[arrow] (desc.east) -- (benef.west) node[midway, above, font=\tiny, popmuted] {Moreover / Besides};
\draw[arrow] (benef.east) -- (conc.west) node[midway, above, font=\tiny, popmuted] {Finally / In conclusion};

% Global Title
\node[above=0.3cm of intro.north, anchor=south, font=\bfseries\normalsize, text=popdark] {Plan du paragraphe type (8-10 lignes)};
\end{tikzpicture}
\legende{Schéma de structure pour la production écrite guidée.}
\end{popfigure}

\begin{exoresolu}[Modèle de paragraphe : « My favourite recreational activity »]
\textbf{Consigne :} Rédigez un paragraphe de 8-10 lignes en suivant le plan et les contraintes de la méthode.

\textbf{Production proposée :}

My favourite recreational activity is \textbf{playing football}. I \textbf{am keen on} this sport because it combines physical effort and team spirit. \textbf{Firstly}, I \textbf{play football every Saturday afternoon} at the municipal stadium in my neighbourhood \textbf{with my childhood friends}. We have formed an amateur team called "The Lions". \textbf{Moreover}, during the holidays, we often \textbf{organise friendly matches} against neighbouring villages. \textbf{Besides} the fun, this activity brings me huge benefits. \textbf{Physically}, it \textbf{keeps me fit} and \textbf{builds my stamina}. \textbf{Mentally}, after a stressful week of classes, running on the pitch helps me \textbf{unwind} and \textbf{relieve stress}. \textbf{Socially}, it \textbf{strengthens the bonds} with my teammates; we support each other on and off the field. \textbf{In conclusion}, football is not just a \textbf{pastime} for me; it is a \textbf{passion} that structures my week and preserves my well-being. I will never stop playing as long as my legs carry me.

\tcblower

\textbf{Analyse du modèle (Correction) :}
\begin{itemize}
    \item \textbf{Plan respecté :} Intro (l.1), Description Fréquence/Lieu/Partenaires (l.2-3), Bénéfices Physique/Mental/Social (l.4-6), Conclusion (l.7).
    \item \textbf{Lexique thème (5+) :} \eng{playing football, recreational activity, stadium, team, matches, pastime, passion, well-being, unwind, relieve stress, strengthen bonds, keep fit}.
    \item \textbf{Collocations (2+) :} \eng{play football, keep fit, relieve stress, strengthen bonds, unwind}.
    \item \textbf{Familles de mots :} \eng{recreation (n.) / recreational (adj.)} ; \eng{passion (n.) / passionate (adj.)} ; \eng{stress (n.) / stressful (adj.)}.
    \item \textbf{Connecteurs :} \eng{Firstly, Moreover, Besides, In conclusion}.
    \item \textbf{Grammaire :} Present Simple (habitudes), Present Continuous (pas utilisé ici mais possible pour "Right now I am playing..."), \eng{helps me unwind} (BV sans \eng{to} après \eng{help}).
\end{itemize}
\end{exoresolu}

\begin{exercice}[À vous de jouer ! Production écrite]
Rédigez un paragraphe de 8 à 10 lignes sur \textbf{votre} activité de loisir préférée en respectant \textbf{strictement} le plan et les contraintes de la \textbf{Méthode} ci-dessus.
\textbf{Sujet :} "My favourite recreational activity."
\textbf{Contraintes :} 1) Plan en 4 étapes. 2) 5 mots du lexique minimum. 3) 2 collocations (play/do/go). 4) 1 famille de mots variée (ex: compete/competition). 5) Connecteurs logiques.
\end{exercice}

\begin{savaistu}[Le lexique des loisirs au Baccalauréat]
Au Cameroun, les sujets d'expression écrite du Baccalauréat (Séries A, C, D, TI) portent très fréquemment sur la \cle{vie sociale} et les \cle{relations humaines}. Les sujets types incluent : 
\begin{itemize}
    \item "Write a letter to your friend describing how you spent your last holiday." (Récit : Past Simple + Vocab loisirs).
    \item "What are the advantages and disadvantages of practising a sport?" (Essai argumentatif : \eng{keep fit, teamwork, injuries, time-consuming}).
    \item "Traditional games are disappearing. Do you agree?" (Opinion : \eng{songo'o, folktales, cultural identity, modern technology, video games}).
    \item "Describe a festival or a ceremony you attended." (Description : \eng{wedding, dowry, choir, balafon, storytelling, attire}).
\end{itemize}
Maîtriser ce lexique (champs, collocations, idiomes, familles de mots) ne sert pas seulement à réussir un exercice de vocabulaire : c'est la \cle{boîte à outils indispensable} pour traiter sereinement tous les sujets d'expression écrite sur les loisirs.
\end{savaistu}


\newpage

\coursec{Activité d'intégration}

\coursec{Lesson 9 — Vocabulary: Words and expressions related to recreational activities}

\courssub{Champs lexicaux : les grandes familles de loisirs}
\begin{textebox}[Texte support : « Leisure Time in Yaoundé »]
\eng{After a long week of classes at Lycée Bilingue, students look forward to the weekend. Some prefer \textbf{outdoor activities}: they \textbf{hang out} at the \textbf{call-box} near the school gate, \textbf{play football} on the dusty \textbf{pitch} behind the stadium, or \textbf{go jogging} along the streets of Bastos. Others enjoy \textbf{indoor activities}: \textbf{playing video games} like FIFA on their phones, \textbf{watching series} on Netflix, or \textbf{reading novels}. Music lovers \textbf{go dancing} at \textbf{bendskin} parties or \textbf{chill out} listening to \textbf{Coupé-Décalé} and Afrobeat. Whether it is \textbf{socialising} with friends or \textbf{unwinding} alone, everyone needs \textbf{leisure time} to \textbf{recharge their batteries}.}
\end{textebox}

\begin{definition}[Champs lexicaux principaux -- Vocabulaire de base]
\begin{center}
\begin{tabularx}{\linewidth}{|X|X|X|}
\hline
\rowcolor{popblueL} \textbf{Catégorie} & \textbf{Mots-clés (English)} & \textbf{Traduction / Contexte camerounais} \\
\hline
\textbf{Outdoor Activities} & \eng{play football, go jogging, go hiking, go fishing, play basketball, hang out, chill (out), go for a walk} & Activités de plein air. \eng{Pitch} = terrain (souvent en terre). \eng{Call-box} = ancienne cabine téléphonique, point de rencontre. \\
\hline
\textbf{Indoor Activities} & \eng{play video games, play cards, play draughts (dames), play ludo, watch movies/series, read, listen to music, do puzzles, do crosswords} & Activités à l'intérieur. \eng{Draughts} (UK) = \eng{Checkers} (US). \eng{Ludo} très populaire. \\
\hline
\textbf{Arts \& Culture} & \eng{go dancing, go clubbing, go to the cinema, play an instrument (guitar, drums), sing, draw, paint, act} & \eng{Go dancing / clubbing} (boîte, \eng{bendskin} party). \eng{Play the guitar} (instrument $\to$ \eng{play}). \\
\hline
\textbf{Sports \& Fitness} & \eng{do gymnastics, do yoga, do aerobics, do judo, do karate, do exercise, work out, go swimming, go cycling} & \eng{Work out} (verbe phrasal) = s'entraîner. \eng{Do sport} (générique). \\
\hline
\textbf{Socialising \& Relaxing} & \eng{socialise, unwind, relax, chill out, hang out (with), catch up (with), meet up (with), have fun, enjoy oneself} & Verbes d'état/d'action sociale. \eng{Catch up with} = rattraper le temps / discuter longuement. \\
\hline
\end{tabularx}
\end{center}
\end{definition}

\courssub{Collocations : \eng{play}, \eng{do}, \eng{go} et autres verbes associés}
\begin{propriete}[Règle d'or : \eng{PLAY} / \eng{DO} / \eng{GO} + activité]
\begin{center}
\begin{tabularx}{\linewidth}{|X|X|X|}
\hline
\rowcolor{popblueL} \textbf{Verbe} & \textbf{Activités types (Exemples)} & \textbf{Règle / Contexte} \\
\hline
\eng{PLAY} & \eng{football, basketball, handball, volleyball, tennis, cards, video games, draughts, ludo, chess, the guitar, the drums} & \textbf{Jeux de balle / compétition}, \textbf{jeux de société}, \textbf{instruments de musique} (\eng{play the piano}). \eng{Play football} (jamais *\eng{do football}). \\
\hline
\eng{DO} & \eng{gymnastics, yoga, aerobics, judo, karate, exercise, workout, crosswords, puzzles, sport (générique)} & \textbf{Activités individuelles}, \textbf{exercices corporels}, \textbf{jeux cérébraux}. \eng{Do sport} (terme générique), \eng{do some exercise}. \\
\hline
\eng{GO} & \eng{swimming, dancing, jogging, running, hiking, fishing, shopping, clubbing, bowling, skiing, cycling} & \textbf{Activités en \eng{-ing}} impliquant un \textbf{déplacement vers un lieu}. \eng{Go dancing} (boîte, \eng{bendskin}), \eng{go clubbing}, \eng{go shopping}. \\
\hline
\end{tabularx}
\end{center}
\emph{Autres verbes utiles :} \eng{watch} (TV, a match, a movie), \eng{listen to} (music, the radio), \eng{surf} (the Internet, the net), \eng{scroll} (social media), \eng{attend} (a concert, a party).
\end{propriete}

\begin{exemplebox}[Exemples traduits -- Collocations]
\begin{itemize}
    \item \eng{My brother \textbf{plays football} every Saturday.} « Mon frère joue au foot chaque samedi. »
    \item \eng{She \textbf{does yoga} to relax.} « Elle fait du yoga pour se détendre. »
    \item \eng{We often \textbf{go dancing} at the weekend.} « Nous allons souvent danser le week-end. »
    \item \eng{They \textbf{play video games} to \textbf{kill time}.} « Ils jouent aux jeux vidéo pour tuer le temps. »
    \item \eng{He \textbf{works out} at the gym three times a week.} « Il s'entraîne à la salle trois fois par semaine. »
\end{itemize}
\end{exemplebox}

\courssub{Expressions idiomatiques et langage courant des loisirs}
\begin{definition}[Idiomes \& Expressions figées (Niveau Terminale -- À réemployer à l'oral et à l'écrit)]
\begin{center}
\begin{tabularx}{\linewidth}{|l|X|l|}
\hline
\rowcolor{poporangeL} \textbf{Expression} & \textbf{Sens / Contexte} & \textbf{Exemple traduit} \\
\hline
\eng{Hang out (with)} & Passer du temps (sans but précis), traîner & \eng{We \textbf{hang out} at the call-box.} (On traîne au call-box.) \\
\hline
\eng{Chill out / Chill} & Se détendre, relaxer (argot jeune) & \eng{Just \textbf{chill out}, exams are over.} (Détends-toi, les examens sont finis.) \\
\hline
\eng{Be one's cup of tea} & Être ce qu'on aime, correspondre à ses goûts & \eng{Dancing is \textbf{my cup of tea}.} (La danse, c'est mon truc.) \\
\hline
\eng{Recharge one's batteries} & Se ressourcer, récupérer de l'énergie & \eng{Reading \textbf{recharges my batteries}.} (La lecture me ressource.) \\
\hline
\eng{Kill time} & Occuper le temps (souvent en attendant) & \eng{I play games to \textbf{kill time}.} (Je joue pour tuer le temps.) \\
\hline
\eng{A guilty pleasure} & Un plaisir qu'on assume peu (culpabilité légère) & \eng{Watching telenovelas is \textbf{a guilty pleasure}.} (Regarder des telenovelas est un plaisir coupable.) \\
\hline
\eng{Have fun / Enjoy oneself} & S'amuser (verbes standards) & \eng{We \textbf{had fun} at the party.} (On s'est bien amusés à la fête.) \\
\hline
\eng{Socialize / Unwind} & Verbes standards : voir du monde / se détendre & \eng{I \textbf{unwind} by listening to music.} (Je me détends en écoutant de la musique.) \\
\hline
\end{tabularx}
\end{center}
\end{definition}

\courssub{Faux amis et pièges classiques des francophones}
\begin{attention}[Erreurs fréquentes -- \textbf{À bannir absolument}]
\begin{itemize}
    \item \textbf{\eng{Recreation} $\neq$ Récréation (pause scolaire).} \eng{Recreation} = loisirs, activités de détente (registre formel). La pause de 10h = \eng{break} (UK) ou \eng{recess} (US). \textit{Correct} : \eng{Recreational activities}, \eng{a recreation centre}.
    \item \textbf{\eng{Amusement} $\neq$ Amusement (général).} \eng{Amusement} = divertissement (souvent \eng{amusement park}). Pour « s'amuser » : \eng{have fun}, \eng{enjoy oneself}. \textit{Piège} : \eng{*We made an amusement} $\to$ \eng{We had fun}.
    \item \textbf{\eng{Distraction} $\neq$ Distraction (loisir).} \eng{Distraction} = chose qui détourne l'attention (négatif : \eng{distracted driving}). Pour « passe-temps » : \eng{hobby}, \eng{pastime}, \eng{leisure activity}.
    \item \textbf{Collocations verbales :} \eng{*Make sport} (FAUX) $\to$ \eng{Do sport} / \eng{Play sport}. \eng{*Make a party} $\to$ \eng{Have a party} / \eng{Throw a party}. \eng{*Pass time} (verbe) $\to$ \eng{Spend time} / \eng{Kill time}.
    \item \textbf{Prononciation critique :} \eng{Leisure} [« lé-jou » UK / « li-jour » US], \textbf{jamais} *\eng{« lai-jour »}. \eng{Recreation} [« ré-kri-é-ï-chène »], accent tonique sur \eng{-é-}. \eng{Relax} [« ri-laks »], pas *\eng{« ré-lax »}.
\end{itemize}
\end{attention}

\courssub{Familles de mots, prononciation et orthographe}
\begin{propriete}[Word Formation -- Dérivation nominale et adjectivale (Utile pour la synthèse écrite)]
\begin{center}
\begin{tabular}{|l|l|l|l|}
\hline
\rowcolor{popgreenL} \textbf{Verbe} & \textbf{Nom (Action/État)} & \textbf{Nom (Agent)} & \textbf{Adjectifs} \\
\hline
\eng{Relax} & \eng{Relaxation} & \eng{---} & \eng{Relaxed (état) / Relaxing (effet)} \\
\hline
\eng{Entertain} & \eng{Entertainment} & \eng{Entertainer} & \eng{Entertaining / Entertained} \\
\hline
\eng{Compete} & \eng{Competition} & \eng{Competitor} & \eng{Competitive} \\
\hline
\eng{Participate} & \eng{Participation} & \eng{Participant} & \eng{Participatory} \\
\hline
\eng{Admire} & \eng{Admiration} & \eng{Admirer} & \eng{Admirable} \\
\hline
\eng{Organise} & \eng{Organisation} & \eng{Organiser} & \eng{Organised / Organisational} \\
\hline
\end{tabular}
\end{center}
\emph{Règle \eng{-ed} / \eng{-ing} :} \eng{Relaxed} (sentiment ressenti : \eng{I feel relaxed}) vs \eng{Relaxing} (cause/effet : \eng{The music is relaxing}). Même chose pour \eng{bored/boring}, \eng{interested/interesting}, \eng{tired/tiring}.
\end{propriete}

\begin{exemplebox}[Orthographe \& Prononciation -- Points de vigilance]
\begin{itemize}
    \item \eng{Leisure} : \eng{ei} se prononce /e/ (UK) ou /i:/ (US). Pas de \eng{z}.
    \item \eng{Recreation} : double \eng{c} (dur + mou), \eng{tion} final. Accent sur l'antépénultième.
    \item \eng{Amusement} : \eng{muse} + \eng{ment}. Pas *\eng{amusement} avec deux \eng{m}.
    \item \eng{Guilty} : \eng{ui} = /aɪ/. \eng{Pleasure} : \eng{s} = /ʒ/.
    \item \eng{Hang out} : \eng{hang} (irrégulier : hang-hung-hung) mais \eng{hang out} régulier souvent traité comme verbe à particule (hung out).
\end{itemize}
\end{exemplebox}

\courssub{Emploi en contexte : enquête et production d'affiche (Projet « Youth \& Leisure »)}
\begin{textebox}[Contexte : Projet « Youth \& Leisure » pour le Magazine du Lycée Bilingue de Yaoundé]
\eng{The English Club of Lycée Bilingue de Yaoundé is preparing a special edition of the school magazine entitled "Youth \& Leisure in Cameroon". As young reporters, your editorial team (groups of 4) has been assigned a mission: investigate how Terminale A students relax after school and during weekends. Your final product will be a wall poster titled "How Our Class Relaxes" displayed in the school library during the Cultural Week. The best poster will be featured in the magazine.}
\end{textebox}

\begin{textebox}[Modèle de dialogue pour l'enquête (Script d'écoute / Support oral)]
\eng{\textbf{Student A (Reporter):} Excuse me, we're conducting a survey for the school magazine. Do you have a minute? \\
\textbf{Student B (Respondent):} Sure, go ahead. \\
\textbf{Student A:} Great. First, \textbf{what do you usually do in your spare time}? \\
\textbf{Student B:} Well, I \textbf{hang out} with my friends at the \textbf{call-box} near the school gate. We just \textbf{chill out} and listen to Afrobeat. \\
\textbf{Student A:} Nice. \textbf{Do you prefer playing or watching sports}? \\
\textbf{Student B:} I prefer \textbf{playing football} on the dusty pitch behind the stadium. \textbf{It's my cup of tea}! Watching is boring. \\
\textbf{Student A:} What about indoor activities? \textbf{Do you ever play video games or board games}? \\
\textbf{Student B:} Sometimes I \textbf{play video games} on my phone, mostly FIFA. But honestly, \textbf{dancing at "bendskin" parties} \textbf{recharges my batteries} after exams. \\
\textbf{Student A:} Last question: \textbf{How often do you go out}? \\
\textbf{Student B:} \textbf{Twice a week}, usually Friday and Saturday nights. \\
\textbf{Student A:} Thank you very much for your time! \\
\textbf{Student B:} You're welcome. Good luck with the magazine!}
\end{textebox}

\courssub{Consignes de l'activité intégrée (Durée totale : ~1h30 -- 2h)}
\begin{enumerate}
    \item \textbf{Préparation du questionnaire (Groupe, 15 min)} \\
    Rédigez \textbf{5 questions variées} en anglais pour votre enquête. Vous devez impérativement inclure :
    \begin{itemize}
        \item Une question sur les activités de temps libre (\eng{spare time / free time}).
        \item Une question utilisant la collocation \eng{play / do / go + activité} (ex. \eng{Do you go swimming?}).
        \item Une question sur les préférences (\eng{prefer / enjoy / fancy + -ing}).
        \item Une question sur la fréquence (\eng{How often...?}).
        \item Une question ouverte incitant à utiliser une expression idiomatique (ex. \eng{What activity helps you recharge your batteries?}).
    \end{itemize}
    \textit{Conseil :} Utilisez le dialogue modèle comme guide pour la politesse (\eng{Excuse me, Do you have a minute?}) et la relance (\eng{What about...?}).

    \item \textbf{Réalisation de l'enquête — « Mingling » (Classe entière, 20 min)} \\
    Levez-vous. Chaque élève doit interroger \textbf{6 camarades d'autres groupes} (pas de son propre groupe). Alternez les rôles : 3 fois « Reporter », 3 fois « Répondant ». Notez les réponses clés (mots-clés, pas phrases complètes) sur votre grille d'enquête. Utilisez les formules du dialogue modèle.

    \item \textbf{Compilation et analyse des données (Groupe, 15 min)} \\
    Regroupés, mettez en commun vos notes. Établissez un \textbf{tableau de fréquences} (tally chart) pour identifier les 5 activités les plus citées. Calculez les pourcentages approximatifs. Identifiez 3 expressions idiomatiques ou collocations fortes entendues chez vos camarades (ex. \eng{hang out, cup of tea, recharge batteries, kill time, guilty pleasure}).

    \item \textbf{Production de l'affiche « How Our Class Relaxes » (Groupe, 20 min)} \\
    Sur une feuille A3 (fournie), concevez l'affiche avec :
    \begin{enumerate}
        \item \textbf{Titre} : \eng{How Our Class Relaxes — Terminale A, [Année]}.
        \item \textbf{Top 5 Ranking} : Classement visuel (barres, camembert ou liste numérotée) des 5 activités favorites avec pourcentages.
        \item \textbf{Idiom Corner} : 3 phrases complètes, chacune intégrant une expression idiomatique différente entendue lors de l'enquête (ex. \eng{Playing football is our cup of tea!}).
        \item \textbf{Synthesis Paragraph} : Un paragraphe de 60--80 mots résumant les tendances de la classe, utilisant des connecteurs logiques (\eng{overall, however, surprisingly, moreover}) et du vocabulaire de la leçon (\eng{outdoor/indoor activities, socialize, unwind, leisure time}).
    \end{enumerate}

    \item \textbf{Présentation orale et correction collective (Classe, 20 min)} \\
    Chaque groupe présente son affiche en \textbf{2 minutes maximum} (tour de rôle des 4 membres). La classe note les erreurs de collocations (\eng{*make sport}), faux amis (\eng{*recreation} pour \eng{relaxation}), prononciation (\eng{/rɪˈkreɪʃn/} vs \eng{/ˌriːkriˈeɪʃn/}) sur une fiche d'évaluation pair. Le professeur synthétise les 5 erreurs majeures au tableau.
\end{enumerate}

\courssub{Proposition de production et corrigé commenté}
\begin{exoresolu}[Modèle de production attendue -- Groupe « The Chill Reporters »]
\textbf{1. Questionnaire modèle (5 questions) :}
\begin{enumerate}
    \item \eng{What do you usually do in your spare time after school?}
    \item \eng{Do you prefer playing sports, doing exercise, or going dancing?}
    \item \eng{How often do you go out with friends during the week?}
    \item \eng{Which activity helps you recharge your batteries after a hard week?}
    \item \eng{What is your guilty pleasure when you want to kill time?}
\end{enumerate}

\textbf{2. Extrait de grille de compilation (Tally Chart) :}
\begin{center}
\begin{tabular}{|l|c|c|}
\hline
\rowcolor{popblueL} \textbf{Activity} & \textbf{Tally} & \textbf{Total / 24 students} \\
\hline
\eng{Hanging out / Chilling (Call-box, Quarter)} & \eng{|||| |||| ||} & 12 (50\%) \\
\hline
\eng{Playing football (Pitch/Street)} & \eng{|||| |||} & 8 (33\%) \\
\hline
\eng{Listening to music / Dancing (Afrobeat, Coupé-Décalé)} & \eng{|||| ||} & 7 (29\%) \\
\hline
\eng{Playing video games (Phone/Console - FIFA, PES)} & \eng{|||| |} & 6 (25\%) \\
\hline
\eng{Watching movies / Series (Netflix, Local TV)} & \eng{||||} & 5 (21\%) \\
\hline
\eng{Reading / Social media scrolling} & \eng{|||} & 3 (12\%) \\
\hline
\end{tabular}
\end{center}

\textbf{3. Contenu de l'affiche « How Our Class Relaxes » :}

\eng{\textbf{HOW OUR CLASS RELAXES — Terminale A, 2024-2025}}

\eng{\textbf{TOP 5 LEISURE ACTIVITIES}}
\begin{enumerate}
    \item \eng{Hanging out \& Chilling at the "Call-box" — 50\%}
    \item \eng{Playing Football (Street/Pitch) — 33\%}
    \item \eng{Dancing to Afrobeat / Coupé-Décalé — 29\%}
    \item \eng{Gaming (FIFA, PES on Mobile) — 25\%}
    \item \eng{Watching Series / Movies — 21\%}
\end{enumerate}

\eng{\textbf{IDIOM CORNER}}
\begin{itemize}
    \item \eng{"Hanging out at the call-box is \textbf{our cup of tea} after classes!"}
    \item \eng{"A good dance session \textbf{recharges our batteries} better than sleep."}
    \item \eng{"For some, gaming is just a way to \textbf{kill time}, but for others, it's a \textbf{guilty pleasure}."}
\end{itemize}

\eng{\textbf{SYNTHESIS}}
\eng{Overall, our class is very social and active. The majority of students prefer outdoor activities like hanging out at the neighbourhood call-box or playing football on dusty pitches. Surprisingly, only a few mentioned reading as a hobby. However, music and dance remain central to our culture: dancing to Coupé-Décalé helps us unwind and recharge our batteries after stressful exam periods. Moreover, video games, especially football simulations, are a popular indoor alternative. In conclusion, leisure time in Terminale A is mostly about socializing and moving, rather than solitary relaxation.}

\tcblower

\textbf{Commentaire des choix de langue (Pourquoi c'est du niveau Terminale A) :}
\begin{itemize}
    \item \textbf{Collocations justes :} \eng{Hanging out} (go + -ing), \eng{Playing football} (play + ball sport), \eng{Recharge batteries} (idiome figé), \eng{Kill time} (collocation verbale).
    \item \textbf{Vocabulaire précis :} \eng{Call-box} (contexte camerounais : cabine téléphonique/point de rencontre), \eng{Pitch} (terrain), \eng{Coupé-Décalé} (réalité culturelle), \eng{Simulation} (jeu vidéo).
    \item \textbf{Grammaire / Cohésion :} \eng{Overall} (ouvreur), \eng{Surprisingly} (commentaire), \eng{However} (concession), \eng{Moreover} (addition), \eng{In conclusion} (clôture). Utilisation du gérondif après prépositions (\eng{after stressful exam periods}, \eng{about socializing}).
    \item \textbf{Adjectifs en -ed/-ing :} \eng{Stressful} (cause) vs \eng{Stressed} (état) — ici \eng{stressful periods} correct. \eng{Relaxing} implicite dans \eng{unwind}.
    \item \textbf{Registre :} Soutenu mais naturel (\eng{leisure activities, solitary relaxation, central to our culture}). Pas de familier excessif dans la synthèse (réservé à l'Idiom Corner).
    \item \textbf{Orthographe / Ponctuation :} Trait d'union pour \eng{call-box}, \eng{cup of tea} (pas de majuscule), guillemets anglais “ ” pour les citations idiomatiques.
\end{itemize}
\end{exoresolu}

\courssub{Bilan de l'activité}
\begin{aretenir}[Ce que vous devez retenir de cette activité d'intégration]
\begin{itemize}
    \item \textbf{Le lexique des loisirs ne s'apprend pas en liste isolée :} Il vit dans des \textbf{collocations} (\eng{play football, go dancing, do yoga}) et des \textbf{expressions figées} (\eng{cup of tea, recharge batteries}). L'enquête vous a forcé à les \textbf{entendre, prononcer, écrire et lire} en contexte réel.
    \item \textbf{L'oral précède l'écrit :} Le \textit{mingling} (déplacement, échange rapide) développe la \textbf{fluidité} et l'\textbf{écoute active} (noter des mots-clés). C'est la base de la compétence \eng{Listening \& Speaking} du Bac.
    \item \textbf{La synthèse écrite exige de la structure :} Connecteurs logiques (\eng{overall, however, moreover}), passage du brut (chiffres) à l'analyse (tendances), emploi précis du vocabulaire (\eng{socialize, unwind, outdoor/indoor}).
    \item \textbf{L'auto-correction par les pairs est cruciale :} Entendre \eng{*make sport} ou \eng{*recreation} chez un camarade et le corriger (\eng{do sport, leisure}) ancre la forme correcte mieux qu'un cours magistral.
    \item \textbf{Ancrage culturel :} Utiliser \eng{call-box}, \eng{bendskin}, \eng{Coupé-Décalé}, \eng{pitch} montre que l'anglais est une langue \textbf{vivante au Cameroun}, pas seulement un objet d'examen.
\end{itemize}
\end{aretenir}

\begin{experience}[Prolongement possible : « Letter to the Editor » -- Expression écrite type Bac]
\eng{Write a formal letter (150 words) to the Principal suggesting two new recreational facilities for the school (e.g., a modern library corner, a multi-sport court), using arguments from your survey results ("50\% hang out at call-boxes $\to$ need safe social space"). Mobilize: formal letter layout, modal verbs (\eng{should, could, would}), conditional (\eng{If we had..., students would...}), vocabulary of facilities (\eng{premises, equipment, supervision}).}
\end{experience}

\newpage

\coursec{Méthodes et exercices résolus}

\coursec{Lesson 9 — Vocabulary: Words and expressions related to recreational activities}

\courssub{Les grands champs lexicaux des loisirs}

\begin{popfigure}
\begin{tikzpicture}[mindmap, grow cyclic, every node/.style={concept, font=\small}, concept color=popblue!30,
  level 1/.append style={level distance=3.2cm, sibling angle=72, font=\small},
  level 2/.append style={level distance=2.8cm, sibling angle=45, font=\scriptsize}]
\node {Recreational Activities} [clockwise from=90]
  child { node[align=center]{Sports \\ \eng{(play / do)}} 
    child { node {\eng{football}} }
    child { node {\eng{basketball}} }
    child { node {\eng{athletics}} }
    child { node {\eng{karate}} }
    child { node {\eng{wrestling}} } }
  child { node[align=center]{Outdoor Activities \\ \eng{(go + -ing)}} 
    child { node {\eng{hiking}} }
    child { node {\eng{fishing}} }
    child { node {\eng{camping}} }
    child { node {\eng{swimming}} }
    child { node {\eng{jogging}} } }
  child { node[align=center]{Indoor Games \\ \eng{(play / do)}} 
    child { node {\eng{chess}} }
    child { node {\eng{draughts}} }
    child { node {\eng{video games}} }
    child { node {\eng{cards}} }
    child { node {\eng{yoga}} } }
  child { node {Cultural Activities} 
    child { node {\eng{reading}} }
    child { node {\eng{dancing}} }
    child { node {\eng{going to the cinema}} }
    child { node {\eng{attending concerts}} }
    child { node {\eng{visiting museums}} } }
  child { node {Social Activities} 
    child { node {\eng{chatting with friends}} }
    child { node {\eng{attending parties}} }
    child { node {\eng{hanging out}} }
    child { node {\eng{clubbing}} }
    child { node {\eng{volunteering}} } };
\end{tikzpicture}
\legende{Les cinq grandes familles d'activités récréatives et leurs verbes associés}
\end{popfigure}

\begin{methode}[Identifier et employer le vocabulaire des loisirs dans un texte]
Pour exploiter un texte (lu ou écouté) sur les \emph{recreational activities}, suis cette démarche en quatre temps :
\begin{enumerate}
\item \textbf{Repère le champ lexical} : souligne tous les mots liés aux loisirs et classe-les dans les cinq familles ci-dessus (\emph{Sports, Outdoor, Indoor, Cultural, Social}). Exemple : \eng{football, hiking, chess, dancing, chatting}.
\item \textbf{Vérifie le verbe associé (collocation)} : demande-toi systématiquement si l'activité prend \eng{play} (jeu avec ballon/adversaire/instrument), \eng{do} (discipline individuelle, art martial, exercice) ou \eng{go} (activité en \emph{-ing}, souvent de plein air). Note la collocation complète : \eng{play football}, \eng{do karate}, \eng{go hiking}.
\item \textbf{Exploite le contexte (indices textuels)} : un mot inconnu se devine grâce aux mots voisins (co-texte). Si le texte mentionne \eng{pitched a tent, lit a fire, Mount Cameroon}, l'activité est \eng{camping}.
\item \textbf{Réemploie immédiatement} : écris une phrase personnelle avec chaque nouveau mot, en lien avec ton quotidien (Yaoundé, Douala, Buea, ton quartier, ton école).
\end{enumerate}
\end{methode}

\begin{textebox}[Script d'écoute : Weekend plans at Government High School Buea]
\eng{— Hey John, what are you up to this weekend?
— I'm going hiking on Mount Cameroon with the geography club. We'll camp overnight near the summit. What about you?
— I have a chess tournament on Saturday morning at the cultural centre. In the afternoon, I'll just relax and listen to music.
— Sounds nice. My sister is attending a traditional dance rehearsal for the Ngondo festival. She's keen on cultural activities.
— Lucky her! I prefer indoor games. Let's meet up Sunday evening to watch the football match on TV?}
\end{textebox}

\begin{exoresolu}[Compréhension et lexique : Weekend plans]
Énoncé — Listen to the dialogue (or read the script above) and answer the questions.
\begin{enumerate}
\item What outdoor activity is John doing? Where?
\item What indoor game is the speaker playing competitively?
\item What cultural activity is John's sister involved in? For which event?
\item Find two expressions with \eng{go + -ing} and one with \eng{play} in the text.
\item Translate: \eng{“She's keen on cultural activities.”}
\end{enumerate}
\tcblower
Solution détaillée :
\begin{enumerate}
\item John is \eng{going hiking} on \eng{Mount Cameroon}.
\item The speaker is playing in a \eng{chess tournament}.
\item She is \eng{attending a traditional dance rehearsal} for the \eng{Ngondo festival}.
\item \eng{go hiking}, \eng{go camping} (implied by \eng{camp overnight}), \eng{play chess} (tournament), \eng{watch a match} (\eng{watch} + noun, not \eng{play}).
\item « Elle est passionnée par les activités culturelles / Elle adore les activités culturelles. » (\eng{keen on + -ing} = être passionné par / aimer beaucoup).
\end{enumerate}
\end{exoresolu}

\courssub{Maîtriser les collocations : \eng{play}, \eng{do}, \eng{go} et autres verbes}

\begin{propriete}[Règle des trois verbes de base pour les loisirs]
\begin{center}
\begin{tabularx}{\linewidth}{|X|X|X|}
\hline
\rowcolor{popblueL} \textbf{Verbe} & \textbf{Type d'activité} & \textbf{Exemples (traduction)} \\
\hline
\eng{PLAY} & Sports avec ballon, jeux d'adversaire, instruments de musique & \eng{play football} (jouer au foot), \eng{play basketball}, \eng{play chess} (jouer aux échecs), \eng{play cards}, \eng{play the guitar} (jouer de la guitare) \\
\hline
\eng{DO} & Activités individuelles, disciplines, arts martiaux, exercices & \eng{do athletics} (faire de l'athlétisme), \eng{do gymnastics}, \eng{do karate / judo / taekwondo}, \eng{do yoga}, \eng{do exercise}, \eng{do aerobics}, \eng{do sport} (général) \\
\hline
\eng{GO} & Activités en \emph{-ing} (souvent plein air, déplacement) & \eng{go swimming}, \eng{go running / jogging}, \eng{go fishing}, \eng{go hiking}, \eng{go dancing}, \eng{go shopping}, \eng{go sightseeing}, \eng{go camping} \\
\hline
\end{tabularx}
\end{center}
\textbf{Règles d'or :}
\begin{itemize}
\item \textbf{Zéro article} après le verbe : \eng{I play football.} (pas \emph{the football}).
\item \textbf{Temps verbaux} : \eng{play/do} $\rightarrow$ \eng{played/did} (prétérit) ; \eng{go} $\rightarrow$ \eng{went} (prétérit irrégulier). Ex. \eng{We went fishing last Sunday.}
\item \textbf{Autres verbes utiles} : \eng{take up a hobby} (se mettre à un loisir), \eng{give up a sport} (arrêter un sport), \eng{attend a concert / a match / a rehearsal} (assister à), \eng{watch a match / a film} (regarder, sans participer), \eng{win / lose a game / a match / a tournament}.
\end{itemize}
\end{propriete}

\begin{exoresolu}[Choix du verbe correct et mise au temps]
Énoncé — Choose the correct verb (\eng{play}, \eng{do} or \eng{go}) and put it in the correct tense.
\begin{enumerate}
\item Every morning, my grandfather \ldots\ldots{} jogging in the neighbourhood.
\item The students of Tle A \ldots\ldots{} a friendly football match against Tle C last Friday.
\item Amina \ldots\ldots{} karate twice a week at the community centre.
\item Next holidays, we \ldots\ldots{} sightseeing in Limbe.
\item My uncle \ldots\ldots{} the guitar in a local band when he was young.
\item The drama club \ldots\ldots{} a rehearsal every Wednesday. (Indice : verbe \eng{hold} ou \eng{attend} ?)
\end{enumerate}
\tcblower
Solution détaillée :
\begin{enumerate}
\item \eng{goes jogging} : \eng{jogging} en \emph{-ing} $\rightarrow$ \eng{go} ; habitude (\eng{every morning}) $\rightarrow$ présent simple, 3e pers. \eng{goes}. « Chaque matin, mon grand-père va faire du jogging. »
\item \eng{played} : \eng{football} (ballon) $\rightarrow$ \eng{play} ; \eng{last Friday} $\rightarrow$ prétérit \eng{played}. « Les élèves de Tle A ont joué un match amical… »
\item \eng{does} : \eng{karate} (art martial) $\rightarrow$ \eng{do} ; \eng{twice a week} $\rightarrow$ présent simple, 3e pers. \eng{does}. « Amina fait du karaté deux fois par semaine. »
\item \eng{will go sightseeing} (ou \eng{are going to go sightseeing}) : \eng{sightseeing} en \emph{-ing} $\rightarrow$ \eng{go} ; \eng{next holidays} $\rightarrow$ futur. « Nous irons faire du tourisme à Limbe. »
\item \eng{played} : instrument de musique $\rightarrow$ \eng{play} ; \eng{when he was young} (passé révolu) $\rightarrow$ prétérit \eng{played}. « Mon oncle jouait de la guitare… »
\item \eng{holds} (ou \eng{has}) : une répétition se \eng{hold} (organise) ou on \eng{attend} (y participe). Ici, le club l'organise : \eng{holds}. Présent simple, 3e pers. \eng{holds}. « Le club de théâtre organise une répétition… »
\end{enumerate}
\cle{Stratégie} : Deux questions dans l'ordre — (1) Quel verbe de collocation ? (2) Quel temps selon l'indice temporel ?
\end{exoresolu}

\courssub{Expressions idiomatiques et langage courant des loisirs}

\begin{definition}[Expressions figées à connaître pour le Bac]
Ces expressions sont fréquentes à l'oral (Listening) et utiles à l'écrit (Writing/Speaking). Apprends-les par blocs.
\begin{center}
\begin{tabularx}{\linewidth}{|X|X|X|}
\hline
\rowcolor{popblueL} \textbf{Expression anglaise} & \textbf{Sens français} & \textbf{Exemple d'emploi} \\
\hline
\eng{to hang out (with friends)} & traîner, passer du temps (avec des amis) & \eng{We usually hang out at the mall on Saturdays.} \\
\hline
\eng{to chill out / to chill} & se détendre, buller & \eng{After exams, I just want to chill out at home.} \\
\hline
\eng{to catch a movie / a show} & aller voir un film / un spectacle & \eng{Let's catch a movie at the cinema tonight.} \\
\hline
\eng{to hit the gym / the pool} & aller à la salle de sport / à la piscine & \eng{He hits the gym every morning before class.} \\
\hline
\eng{to kill time} & tuer le temps & \eng{I read a magazine to kill time while waiting.} \\
\hline
\eng{to have a good time / a blast} & passer un bon moment / s'éclater & \eng{We had a blast at the party last night.} \\
\hline
\eng{to be into something} & s'intéresser à, aimer (argot) & \eng{Are you into video games?} \\
\hline
\eng{to take up a hobby} & se mettre à un passe-temps & \eng{I've decided to take up photography.} \\
\hline
\eng{to unwind} & se détendre (après le stress) & \eng{Reading helps me unwind after a long day.} \\
\hline
\end{tabularx}
\end{center}
\end{definition}

\begin{exemplebox}[Dialogue naturel : Planning the weekend]
\eng{— \textbf{What are you up to} this weekend? (Qu'est-ce que tu fais ce week-end ?)
— Not much. I'll probably \textbf{chill out} at home and \textbf{catch up on} some sleep. You?
— I'm \textbf{into} hiking lately. A group of us \textbf{are hitting the trails} on Mount Cameroon Saturday.
— Sounds fun, but I'm \textbf{not really into} strenuous exercise. I'd rather \textbf{catch a film} at the cinema.
— Suit yourself. \textbf{Let me know} if you change your mind.}
\tcblower
\emph{Traduction :}
\eng{— Qu'est-ce que tu fais ce week-end ? — Pas grand-chose. Je vais probablement buller à la maison et rattraper mon sommeil. Et toi ? — Je m'y suis mis à la rando dernièrement. Un groupe part sur les sentiers du Mont Cameroun samedi. — Ça a l'air sympa, mais je ne suis pas trop fan d'exercice intense. Je préfère aller voir un film. — Comme tu veux. Tiens-moi au courant si tu changes d'avis.}
\end{exemplebox}

\courssub{Éviter les faux amis et les calques du français}

\begin{methode}[Déjouer les pièges lexicaux du thème]
\begin{enumerate}
\item \textbf{Liste personnelle} : note les faux amis du thème avec un exemple contraire pour chaque.
\item \textbf{Méfiance « transparence »} : plus un mot ressemble au français, plus tu vérifies (ex. \eng{lecture, assist, actually, sensible}).
\item \textbf{Interdiction du mot-à-mot} : « faire du sport » $\neq$ \emph{make sport} $\rightarrow$ \eng{do sport / play sports}. « Passer le temps » $\neq$ \emph{pass time} $\rightarrow$ \eng{spend time / kill time / pass the time}.
\item \textbf{Prépositions figées} : \eng{interested \textbf{in}}, \eng{keen \textbf{on}}, \eng{fond \textbf{of}}, \eng{good \textbf{at}}, \eng{bad \textbf{at}}, \eng{listen \textbf{to}}, \eng{look \textbf{at}}.
\end{enumerate}
\end{methode}

\begin{center}
\begin{tabularx}{\linewidth}{|X|X|X|}
\hline
\rowcolor{popblueL} \textbf{Faux ami anglais} & \textbf{Sens réel en anglais} & \textbf{Mot anglais pour le sens français} \\
\hline
\eng{a lecture} & un cours magistral, une conférence & \eng{a reading} (lecture personnelle), \eng{a text} \\
\hline
\eng{to assist} & aider, secourir & \eng{to attend} (assister à un événement), \eng{to watch} (regarder un match) \\
\hline
\eng{a spectator} & un spectateur & Correct pour « spectateur », mais « assister à » = \eng{attend / watch} \\
\hline
\eng{sensible} & raisonnable, sensé & \eng{sensitive} (sensible, qui ressent fort) \\
\hline
\eng{actually} & en fait, réellement & \eng{currently} / \eng{at the moment} (actuellement) \\
\hline
\eng{an experience} & une expérience (vécue), un vécu & \eng{an experiment} (expérience scientifique) ; « vivre une aventure » = \eng{have an adventure} \\
\hline
\eng{a recreation} & une récréation, un passe-temps & \eng{a break} / \eng{recess} (récréation scolaire) \\
\hline
\eng{to practise} & exercer (métier), s'entraîner (compétence) & \eng{to play} (un sport), \eng{to do} (un sport), \eng{to train} (s'entraîner sérieusement) \\
\hline
\end{tabularx}
\end{center}

\begin{exoresolu}[Correction de phrases calquées sur le français]
Énoncé — Each sentence contains a mistake caused by French interference. Correct it.
\begin{enumerate}
\item \eng{I assist to football matches every weekend.}
\item \eng{During the holidays, I like to make sport with my friends.}
\item \eng{She is actually reading a novel by a Cameroonian writer.}
\item \eng{We passed a good moment at the Christmas party.}
\item \eng{My brother is interested on video games.}
\item \eng{He practises football every evening.}
\item \eng{The match was very amused.}
\end{enumerate}
\tcblower
Solution détaillée :
\begin{enumerate}
\item \eng{I \textbf{attend} football matches every weekend} (ou \eng{I \textbf{watch} football matches}). \eng{Assist} = aider. « J'assiste à des matchs… »
\item \eng{I like to \textbf{do} sport} (ou \eng{\textbf{play} sports}). \eng{Make sport} n'existe pas. « Je fais du sport… »
\item \eng{She is \textbf{currently} reading…} (ou \eng{\textbf{at the moment}}). \eng{Actually} = « en fait ». « Elle lit actuellement… »
\item \eng{We \textbf{had a good time} at the Christmas party.} Idiome : \eng{have a good time}. « Nous avons passé un bon moment… »
\item \eng{My brother is interested \textbf{in} video games.} Préposition obligatoire \eng{in}. « Mon frère s'intéresse aux jeux vidéo. »
\item \eng{He \textbf{plays} football every evening.} \eng{Practise} ne s'emploie pas directement avec un sport-ballon. « Il joue au foot… »
\item \eng{The match was very \textbf{amusing}.} (ou \eng{entertaining}). \eng{Amused} = qui ressent de l'amusement (sujet animé). \eng{Amusing} = qui procure de l'amusement (sujet chose/événement). « Le match était très amusant. »
\end{enumerate}
\cle{Règle} : \eng{-ing} = cause du sentiment (la chose) ; \eng{-ed} = ressenti du sentiment (la personne). \eng{The film is boring} / \eng{I am bored}.
\end{exoresolu}

\courssub{Familles de mots, prononciation et orthographe}

\begin{propriete}[Construction des mots de la famille du loisir]
\begin{center}
\begin{tabularx}{\linewidth}{|X|X|X|X|}
\hline
\rowcolor{popblueL} \textbf{Verbe / Racine} & \textbf{Nom d'activité / action} & \textbf{Nom de personne} & \textbf{Adjectifs (\eng{-ing} / \eng{-ed})} \\
\hline
\eng{relax} & \eng{relaxation} & \eng{—} & \eng{relaxing} (détendant) / \eng{relaxed} (détendu) \\
\hline
\eng{recreate} & \eng{recreation} & \eng{—} & \eng{recreational} (récréatif) \\
\hline
\eng{compete} & \eng{competition} & \eng{competitor} & \eng{competitive} \\
\hline
\eng{entertain} & \eng{entertainment} & \eng{entertainer} & \eng{entertaining} / \eng{entertained} \\
\hline
\eng{enjoy} & \eng{enjoyment} & \eng{—} & \eng{enjoyable} (agréable) / \eng{enjoyed} \\
\hline
\eng{amuse} & \eng{amusement} & \eng{—} & \eng{amusing} / \eng{amused} \\
\hline
\eng{swim} & \eng{swimming} & \eng{swimmer} & \eng{—} \\
\hline
\eng{win} & \eng{winning} (rare), \eng{victory} & \eng{winner} & \eng{—} \\
\hline
\eng{perform} & \eng{performance} & \eng{performer} & \eng{—} \\
\hline
\eng{spectate} & \eng{spectator} (nom d'agent) & \eng{spectator} & \eng{—} \\
\hline
\end{tabularx}
\end{center}
\textbf{Suffixes clés :} \emph{-tion/-sion} (noms abstraits), \emph{-ment} (noms d'action/résultat), \emph{-ing} (noms d'activité/gerund), \emph{-er/-or} (agent), \emph{-ist} (spécialiste : \eng{guitarist, tourist}), \emph{-ful/-able} (adjectifs : \eng{enjoyable, restful}), \emph{-al} (adjectifs relationnels : \eng{recreational, cultural}).
\end{propriete}

\begin{aretenir}[Prononciation et orthographe : points de vigilance]
\begin{itemize}
\item \textbf{Accent tonique} : \eng{lei\textbf{sure}} (« lé-jeur » UK / « li-jur » US), \eng{recre\textbf{a}tion}, \eng{enter\textbf{tain}ment}, \eng{com\textbf{pe}tition}, \eng{re\textbf{hear}sal}.
\item \textbf{Lettres muettes} : \eng{listen} (\eng{t} muet), \eng{walk/hike/talk} (\eng{l} muet), \eng{guitar} (\eng{u} muet après \eng{g}), \eng{rhythm} (\eng{h} muet).
\item \textbf{Double consonne en \eng{-ing} / \eng{-er}} : voyelle courte accentuée + consonne finale $\rightarrow$ doublement.
\eng{swim $\rightarrow$ swimming, swimmer} ; \eng{run $\rightarrow$ running, runner} ; \eng{win $\rightarrow$ winning, winner} ; \eng{hit $\rightarrow$ hitting}.
\item \textbf{Pas de doublement} : \eng{play $\rightarrow$ playing} (voyelle + \eng{y}), \eng{read $\rightarrow$ reading} (digramme), \eng{visit $\rightarrow$ visiting} (accent sur 1ère syllabe).
\item \textbf{Orthographe \eng{rehearsal}} : \eng{re-hear-sal} (entendre \eng{hear} à l'intérieur).
\end{itemize}
\end{aretenir}

\begin{exoresolu}[Formation de mots (Word Formation)]
Énoncé — Complete each sentence with the correct form of the word in brackets.
\begin{enumerate}
\item Watching comedies is very \ldots\ldots{} (relax).
\item The \ldots\ldots{} of the singing competition received a trophy. (win)
\item The children were very \ldots\ldots{} during the long speech. (bore)
\item Our school offers many \ldots\ldots{} activities after classes. (recreation)
\item The \ldots\ldots{} of the new sports complex took place on Saturday. (open)
\item She is a talented \ldots\ldots{} ; she plays the piano beautifully. (music)
\item The \ldots\ldots{} lasted for two hours. (perform)
\end{enumerate}
\tcblower
Solution détaillée :
\begin{enumerate}
\item \eng{relaxing} : sujet = activité qui \emph{provoque} la détente $\rightarrow$ adj. en \emph{-ing}. « Regarder des comédies est très relaxant. »
\item \eng{winner} : nom de personne. \eng{win} + \eng{er} $\rightarrow$ doublement \eng{n} : \eng{winner}. « Le gagnant… a reçu un trophée. »
\item \eng{bored} : les enfants \emph{ressentent} l'ennui $\rightarrow$ adj. en \emph{-ed}. « Les enfants s'ennuyaient… »
\item \eng{recreational} : adjectif devant \eng{activities}. \eng{recreation} + \emph{-al} = \eng{recreational}. « Notre école propose des activités récréatives… »
\item \eng{opening} : nom d'action (inauguration) après \eng{the}. \eng{open} + \emph{-ing}. « L'inauguration… a eu lieu samedi. »
\item \eng{musician} : nom de personne en \emph{-ian} (musicien). Attention : pas \emph{musicer}. « C'est une musicienne talentueuse… »
\item \eng{performance} : nom d'action en \emph{-ance}. « La représentation a duré deux heures. »
\end{enumerate}
\end{exoresolu}

\courssub{Réemployer le lexique à l'oral et à l'écrit}

\begin{methode}[Produire un dialogue ou un paragraphe sur ses loisirs (Niveau Bac)]
\begin{enumerate}
\item \textbf{Structures toutes faites pour exprimer les goûts} (varie \eng{like} !) :
\begin{itemize}
\item \eng{I am fond of + -ing} / \eng{I am keen on + -ing} / \eng{I am interested in + -ing}
\item \eng{I enjoy + -ing} / \eng{I love + -ing} / \eng{I can't stand + -ing}
\item \eng{In my free time / During my leisure time, I usually\ldots}
\item \eng{What do you do for fun?} / \eng{How do you spend your weekends?}
\item \eng{I would rather + infinitif sans to} (préférence immédiate) ; \eng{I prefer + -ing} (préférence générale).
\end{itemize}
\item \textbf{Structure du paragraphe (80–100 mots)} :
\begin{enumerate}
\item Introduction : loisir préféré + fréquence.
\item Développement : quand ? avec qui ? où ? (collocations correctes).
\item Raisons (2 minimum) : bien-être, social, culture, défi.
\item Anecdote ou événement marquant (prétérit).
\item Conclusion : valeur personnelle / projection future.
\end{enumerate}
\item \textbf{Auto-correction} : relis en traquant \eng{play/do/go}, prépositions (\eng{in/on/at/to}), 3e personne du présent (\eng{she plays}), \eng{-ing/-ed}.
\end{enumerate}
\end{methode}

\begin{experience}[Jeu de rôle : “Weekend Planner” (Par groupes de 3)]
\textbf{Rôles :} A = Organisateur(trice) ; B = Sportif(ve) ; C = Casanier(ère) / Culturel(le).
\textbf{Consigne :} Planifiez le samedi idéal. Utilisez au moins 5 collocations (\eng{play/do/go}), 3 expressions d'opinion (\eng{I'm keen on…}, \eng{I'd rather…}, \eng{I'm not into…}) et 2 expressions idiomatiques (\eng{hang out, chill out, catch a movie}).
\textbf{Exemple de début :}
\eng{A: “Okay team, let's plan Saturday. What are you up to?”
B: “I’m keen on going hiking. Let’s hit the trails early.”
C: “I’m not really into strenuous exercise. I’d rather catch a movie or just chill out at home.”
A: “Let’s compromise. Morning hike, afternoon cinema?”}
\textbf{Objectif :} Fluidité, vocabulaire précis, négociation orale.
\end{experience}

\begin{savaistu}[Le sport et la fête au Cameroun : vecteurs de lien social]
Le \textbf{football} est roi : les \eng{“bendskins”} (motos-taxis) arborent les couleurs de l'équipe favorite (Canon, Coton Sport, Union de Douala). Les \eng{“buvettes”} se transforment en \eng{“viewing centres”} les soirs de match (CAN, Ligue des Champions).
La \textbf{Course de l'Espoir (Mount Cameroon Race of Hope)} à Buea est une institution : course de montagne extrême (42 km aller-retour), attirant coureurs locaux et internationaux. C'est un événement \eng{recreational} majeur.
Les \textbf{danse traditionnelles} (Bikutsi, Makossa, Assiko, Bottle Dance) ne sont pas que spectacle : elles structurent les \eng{“njangi”} (tontines), les mariages, les funérailles. \eng{Rehearsals} (répétitions) hebdomadaires sont des moments de cohésion intergénérationnelle.
\end{savaistu}

\begin{exoresolu}[Composition guidée : My favourite leisure activity]
Énoncé — Write a paragraph of about 80–100 words on: \eng{Describe your favourite leisure activity. Say when and with whom you practise it, and explain why you enjoy it.}
\tcblower
Solution détaillée — \textbf{Démarche :} (1) Choix : danse traditionnelle (contexte camerounais). (2) Temps : présent simple (habitudes), prétérit (événement passé). (3) Lexique ciblé : \eng{go dancing, attend rehearsals, keen on, have a good time, unwind, cultural heritage}. (4) Connecteurs : \eng{Besides, Whenever, For me}.

\eng{My favourite leisure activity is traditional dancing. I am a member of my school dance club in Yaoundé, and we attend rehearsals every Wednesday and Friday afternoon. I am really keen on dancing because it keeps me fit and helps me unwind after stressful classes. Besides, I have made many friends in the club, and we always have a good time together. Last term, we even won the first prize at the regional youth festival. Whenever I feel anxious before an exam, I go dancing and I immediately feel relaxed. For me, dancing is not just a hobby; it is a way of celebrating our Cameroonian cultural heritage.}

\textbf{Points forts à reproduire :} Présent simple (habitudes), prétérit (événement ponctuel \eng{won}), connecteurs logiques, variété lexicale (\eng{keen on, unwind, attend rehearsals, go dancing, cultural heritage}), collocations correctes.
\end{exoresolu}

\courssub{Entraînement autonome (Exercices non corrigés)}

\begin{exercice}[Vocabulary \& Grammar Practice]
\textbf{A. Collocations :} Complete with \eng{play}, \eng{do} or \eng{go} in the correct form.
\begin{enumerate}
\item My father \ldots\ldots{} fishing every Sunday morning. (habit)
\item The girls \ldots\ldots{} gymnastics at the sports complex. (present)
\item We \ldots\ldots{} camping in Waza Park last holiday. (past)
\item She \ldots\ldots{} the violin beautifully. (present)
\item They \ldots\ldots{} a chess tournament next week. (future)
\end{enumerate}

\textbf{B. False Friends :} Choose the correct word.
\begin{enumerate}
\item I want to \textbf{assist / attend} the conference on youth employment.
\item The lesson was very \textbf{amused / amusing}.
\item He is \textbf{actually / currently} working on his project.
\item She has a lot of \textbf{experience / experiment} in teaching.
\item Don't \textbf{make / do} sport right after eating.
\end{enumerate}

\textbf{C. Word Formation :} Complete with the right form of the word in brackets.
\begin{enumerate}
\item The \ldots\ldots{} (organize) of the festival did a great job.
\item Reading novels is an \ldots\ldots{} (enjoy) way to spend time.
\item He is a professional \ldots\ldots{} (swim).
\item The \ldots\ldots{} (compete) was tough but fair.
\item We need more \ldots\ldots{} (recreation) facilities in our neighbourhood.
\end{enumerate}

\textbf{D. Translation :} Translate into English.
\begin{enumerate}
\item Pendant mon temps libre, j'aime bien traîner avec mes amis au marché.
\item Ma sœur s'est mise à la photographie l'année dernière.
\item Nous avons passé un moment formidable à la fête de la jeunesse.
\item Il est passionné de jeux vidéo, mais il n'est pas très bon aux échecs.
\item Chaque matin, elle fait du yoga pour se détendre.
\end{enumerate}
\end{exercice}

\begin{corrige}[Exercices d'entraînement]
\textbf{A. Collocations}
\begin{enumerate}
\item \eng{goes fishing} (\eng{go} + \emph{-ing}, présent 3e pers.)
\item \eng{do gymnastics} (\eng{do} + discipline, présent)
\item \eng{went camping} (\eng{go} + \emph{-ing}, prétérit irrégulier \eng{went})
\item \eng{plays the violin} (\eng{play} + instrument, présent 3e pers.)
\item \eng{will play} / \eng{are going to play} (\eng{play} + tournoi/jeu adversaire, futur)
\end{enumerate}

\textbf{B. False Friends}
\begin{enumerate}
\item \eng{attend} (\eng{assist} = aider)
\item \eng{amusing} (cause du sentiment, chose)
\item \eng{currently} (\eng{actually} = en fait)
\item \eng{experience} (vécu ; \eng{experiment} = expérience scientifique)
\item \eng{do} (\eng{make sport} n'existe pas)
\end{enumerate}

\textbf{C. Word Formation}
\begin{enumerate}
\item \eng{organizers} (ou \eng{organisation} si singulier, mais \eng{The organizers} pluriel plus naturel pour "did a great job") — \eng{organizer} (nom agent \emph{-er}).
\item \eng{enjoyable} (adj. \emph{-able}, "qui procure du plaisir").
\item \eng{swimmer} (nom agent \emph{-er}, doublement \eng{m}).
\item \eng{competition} (nom en \emph{-tion}) ou \eng{competitors} (participants). Contexte "was tough" $\rightarrow$ \eng{competition} (l'événement).
\item \eng{recreational} (adj. en \emph{-al} pour qualifier \eng{facilities}).
\end{enumerate}

\textbf{D. Translation (propositions)}
\begin{enumerate}
\item \eng{In my free time, I like hanging out with my friends at the market.}
\item \eng{My sister took up photography last year.}
\item \eng{We had a blast / a great time at the Youth Day party.}
\item \eng{He is keen on video games, but he is not very good at chess.}
\item \eng{Every morning, she does yoga to unwind / relax.}
\end{enumerate}
\end{corrige}

\begin{attention}[Les 5 erreurs fatales au Bac — Check-list finale]
\begin{enumerate}
\item \textbf{Collocation incorrecte} : \emph{make sport, practise football, go to swim}. \rightarrow \eng{do sport / play sports}, \eng{play football}, \eng{go swimming}.
\item \textbf{Faux amis} : \eng{I assist to the match}, \eng{She is actually here}, \eng{It was very amused}. \rightarrow \eng{I attend/watch the match}, \eng{She is currently here}, \eng{It was very amusing}.
\item \textbf{Confusion \eng{-ing} / \eng{-ed}} : \eng{I am boring}, \eng{The match was bored}. \rightarrow \eng{I am bored} (ressenti), \eng{The match was boring} (cause).
\item \textbf{Prépositions fautives} : \eng{interested on, keen in, fond in, good in, listen music}. \rightarrow \eng{interested \textbf{in}}, \eng{keen \textbf{on}}, \eng{fond \textbf{of}}, \eng{good \textbf{at}}, \eng{listen \textbf{to} music}.
\item \textbf{Orthographe \eng{-ing} / \eng{-er}} : \emph{swiming, wining, runing, beginer}. \rightarrow \eng{swimming, winning, running, beginner} (doublement consonne après voyelle brève accentuée).
\end{enumerate}
\end{attention}

\newpage

\coursec{Exercices}

\courssub{Je vérifie mes connaissances}

\begin{exercice}[QCM : la bonne collocation]
Pour chaque activité, choisis la bonne collocation parmi \textbf{a}, \textbf{b} ou \textbf{c}. Entoure la lettre de la bonne réponse.

\begin{enumerate}
\item .............. football \quad a) do \quad b) play \quad c) go
\item .............. swimming \quad a) play \quad b) do \quad c) go
\item .............. karate \quad a) do \quad b) play \quad c) make
\item .............. chess \quad a) do \quad b) go \quad c) play
\item .............. shopping \quad a) play \quad b) go \quad c) do
\item .............. gymnastics \quad a) play \quad b) do \quad c) go
\item .............. the guitar \quad a) do \quad b) play \quad c) go
\item .............. hiking \quad a) play \quad b) do \quad c) go
\item .............. video games \quad a) play \quad b) make \quad c) go
\item .............. yoga \quad a) play \quad b) go \quad c) do
\end{enumerate}
\end{exercice}

\begin{exercice}[Vrai ou faux ?]
Lis le texte suivant, puis indique si chaque affirmation est \textbf{True} ou \textbf{False}. Justifie chaque réponse par une courte citation du texte.

\begin{textebox}[Leisure time in Yaoundé]
After school, many young people in Yaoundé spend their free time playing football in the neighbourhood or watching matches on television. Others prefer indoor activities: they listen to music, chat with friends on their phones or play video games. At weekends, some families go hiking around Mount Fébé, while others attend church choir practice. Surprisingly, very few teenagers read novels for pleasure; most say they have no time for books.
\end{textebox}

\begin{enumerate}
\item All young people in Yaoundé play football after school.
\item Some leisure activities mentioned take place indoors.
\item Families never go hiking at weekends.
\item Most teenagers say they read a lot of novels.
\end{enumerate}
\end{exercice}

\begin{exercice}[Vocabulaire : trouve le mot]
Complète chaque définition par le mot anglais du champ lexical des loisirs qui convient. Choisis parmi : \eng{hobby, pastime, leisure, entertainment, recreation, amusement}.

\begin{enumerate}
\item An activity you do regularly for pleasure in your free time : \eng{a ..............}
\item Free time, when you are not working or studying : \eng{..............}
\item Something that makes time pass pleasantly, like a film or a show : \eng{..............}
\item An enjoyable way of spending time, often a game or a sport : \eng{a ..............}
\end{enumerate}
\end{exercice}

\begin{exercice}[Familles de mots : complète]
Complète chaque phrase avec la forme correcte du mot entre parenthèses (nom, verbe, adjectif).

\begin{enumerate}
\item Gardening is my grandmother's favourite .............. . (\eng{relax})
\item After the exams, the students felt completely .............. . (\eng{relax})
\item We spent a very .............. evening by the swimming pool. (\eng{relax})
\item Reading is a good form of .............. for me. (\eng{entertain})
\item The clown was very .............. ; all the children laughed. (\eng{amuse})
\item My father finds fishing extremely .............. . (\eng{bore})
\end{enumerate}
\end{exercice}

\courssub{Je m'entraîne}

\begin{exercice}[Traduction]
Traduis les phrases suivantes en anglais, en utilisant le vocabulaire et les collocations étudiés dans la leçon.

\begin{enumerate}
\item Le week-end, je fais de la natation et je joue aux échecs avec mon frère.
\item Beaucoup de jeunes Camerounais jouent au football après les cours.
\item Ma sœur n'aime pas les jeux vidéo ; elle préfère lire des romans.
\item Pour me détendre, j'écoute de la musique ou je me promène avec mes amis.
\item Pendant les vacances, nous faisons de la randonnée dans les montagnes de l'Ouest.
\end{enumerate}
\end{exercice}

\begin{exercice}[Corrige les faux amis]
Chaque phrase contient une erreur typique de francophone (faux ami ou calque du français). Souligne l'erreur et réécris la phrase correctement.

\begin{enumerate}
\item I assisted to the football match last Saturday.
\item I make sport every day to stay healthy.
\item In my free time, I like to repose myself.
\item She passed her weekend at her grandmother's village.
\item We amused ourselves very much at the party. (\emph{attention au sens exact !})
\item He plays with the computer all evening.
\end{enumerate}
\end{exercice}

\begin{exercice}[Texte à trous]
Complète le texte suivant avec les mots de la banque ci-dessous. Chaque mot ne peut être utilisé qu'une seule fois.

\begin{center}
\eng{hobbies — relaxing — leisure — play — go — entertainment — pastime — bored}
\end{center}

\begin{textebox}[Free time in Douala]
In Douala, young people have many different .............. . After school, boys often .............. football in the streets, while girls sometimes .............. dancing or listen to makossa music. For many families, television is the main form of .............. in the evening. At weekends, some people .............. swimming at the beach in Limbe. Reading is not a very popular .............. among teenagers, and many say they feel .............. when they have no phone. However, doctors say that .............. activities like walking and gardening are excellent for our health, especially during our .............. time.
\end{textebox}
\end{exercice}

\begin{exercice}[Appariement : expressions idiomatiques]
Associe chaque expression anglaise (1--8) à son équivalent français (a--h).

\begin{center}
\begin{tabularx}{\linewidth}{|X|X|}
\hline
\rowcolor{popblueL} \textbf{English expression} & \textbf{Équivalent français} \\
\hline
1. \eng{to let off steam} & a) perdre son temps \\
\hline
2. \eng{to kill time} & b) se défouler, se détendre \\
\hline
3. \eng{it's not my cup of tea} & c) se détendre, se relaxer \\
\hline
4. \eng{to chill out} & d) ce n'est pas mon truc \\
\hline
5. \eng{to have a whale of a time} & e) s'amuser comme un fou \\
\hline
6. \eng{to be a couch potato} & f) être un passionné de sport \\
\hline
7. \eng{to be a sports fanatic} & g) être un vrai paresseux devant la télé \\
\hline
8. \eng{to take up a hobby} & h) se mettre à une activité de loisir \\
\hline
\end{tabularx}
\end{center}
\end{exercice}

\courssub{Je me prépare au Bac}

\begin{exercice}[Compréhension et langue — type Bac]
Lis attentivement le texte suivant, puis réponds aux questions.

\begin{textebox}[How Cameroonian teenagers spend their free time]
In many towns of Cameroon, free time is a precious thing for secondary school students. After long hours in the classroom and homework in the evening, weekends are the only real moment for recreation. In Yaoundé and Douala, football remains king: boys organise matches on dusty pitches, and even girls are joining school teams more and more. Basketball and handball are also gaining popularity in some lycées.

Not all leisure activities are sporty, however. Many teenagers prefer quieter pastimes. Listening to music — makossa, afrobeats or gospel — is extremely common, and dancing competitions are often organised during school cultural weeks. Some students join drama clubs or choirs, where they learn to speak and sing in front of an audience. A smaller group enjoys reading novels or writing poems, although teachers complain that books are losing the battle against mobile phones.

Mobile phones have indeed changed the way young people relax. Social media, video games and online chats now occupy a large part of their evenings. Parents worry that their children spend too much time on screens and not enough on outdoor activities. “When I was young, we played outside until nightfall,” says Mr Ndi, a father of three in Bamenda. “Today, my son plays football on his phone!”

Yet specialists remind us that recreation is not a waste of time. Playing, singing, reading or simply resting helps the brain to recover and improves school results. The secret, they say, is balance: a little sport, a little culture, a little rest — and not too much screen time.
\end{textebox}

\textbf{A. Comprehension questions} — Answer in your own words.

\begin{enumerate}
\item According to the text, when do Cameroonian students really have free time?
\item Name two sports mentioned in the text, apart from football.
\item What quiet pastimes do some teenagers prefer? Give two examples.
\item What do parents worry about?
\item In your opinion, why is recreation important for students? (2--3 sentences)
\end{enumerate}

\textbf{B. Language work}

\begin{enumerate}
\item Find in the text words or expressions meaning : \textbf{a)} \eng{free time activities} (paragraph 1) ; \textbf{b)} \eng{becoming more liked} (paragraph 1) ; \textbf{c)} \eng{to rest and get energy back} (paragraph 4).
\item Complete with the correct collocation (\eng{play}, \eng{do} or \eng{go}) : \eng{.............. basketball} ; \eng{.............. dancing} ; \eng{.............. video games}.
\item Rewrite using your own words : \eng{“books are losing the battle against mobile phones”} (paragraph 2).
\end{enumerate}
\end{exercice}

\begin{exercice}[Traduction — type Bac]
Traduis en anglais le passage suivant.

\begin{textebox}[Texte à traduire]
Pendant leur temps libre, les jeunes de ma ville pratiquent beaucoup d'activités. Les garçons jouent au football le samedi après-midi, tandis que les filles préfèrent souvent danser ou écouter de la musique. Certains élèves font de la natation ou de la randonnée avec leurs parents. Moi, je ne suis pas un fanatique du sport : je préfère lire des romans ou jouer aux échecs avec mon grand-père. Ces loisirs m'aident à me détendre après une longue semaine de cours.
\end{textebox}

\emph{Conseil : attention aux collocations avec} \eng{play}, \eng{do} \emph{et} \eng{go}, \emph{aux faux amis et au temps des verbes.}
\end{exercice}

\begin{exercice}[Composition — type Bac]
\textbf{Topic :} \eng{Write an article of about 150 words for your school magazine about the recreational activities of students in your school. Say which activities are the most popular, which ones you personally prefer, and why recreation is important for students.}

\medskip
\textbf{Grille d'évaluation (10 points) :}

\begin{center}
\begin{tabularx}{\linewidth}{|X|l|}
\hline
\rowcolor{popblueL} \textbf{Critère} & \textbf{Points} \\
\hline
Lexique des loisirs riche et varié (champs lexicaux, expressions idiomatiques) & 4 pts \\
\hline
Collocations correctes (\eng{play / do / go} + activité) & 3 pts \\
\hline
Correction grammaticale (temps, articles, prépositions) & 2 pts \\
\hline
Cohérence et organisation de l'article (introduction, développement, conclusion) & 1 pt \\
\hline
\end{tabularx}
\end{center}

\emph{Ton article doit contenir au moins trois collocations différentes et une expression idiomatique vue dans la leçon.}
\end{exercice}

\newpage

\coursec{Corrigés des exercices}

\begin{corrige}[Exercice 1 — QCM : la bonne collocation]

\textbf{Réponses :}

\begin{center}
\begin{tabular}{|l|l|l|}
\hline
\rowcolor{popblueL} \textbf{N°} & \textbf{Réponse} & \textbf{Justification} \\
\hline
1 & b) \eng{play football} & sport avec ballon $\rightarrow$ \eng{play} \\
\hline
2 & c) \eng{go swimming} & activité en \eng{-ing} $\rightarrow$ \eng{go} \\
\hline
3 & a) \eng{do karate} & art martial / sport individuel $\rightarrow$ \eng{do} \\
\hline
4 & c) \eng{play chess} & jeu de société $\rightarrow$ \eng{play} \\
\hline
5 & b) \eng{go shopping} & activité en \eng{-ing} $\rightarrow$ \eng{go} \\
\hline
6 & b) \eng{do gymnastics} & sport individuel / discipline $\rightarrow$ \eng{do} \\
\hline
7 & b) \eng{play the guitar} & instrument de musique $\rightarrow$ \eng{play + the} \\
\hline
8 & c) \eng{go hiking} & activité en \eng{-ing} $\rightarrow$ \eng{go} \\
\hline
9 & a) \eng{play video games} & jeu $\rightarrow$ \eng{play} \\
\hline
10 & c) \eng{do yoga} & activité individuelle de détente $\rightarrow$ \eng{do} \\
\hline
\end{tabular}
\end{center}

\textbf{Rappel de la règle :} \eng{play} + sports d'équipe ou avec ballon et jeux (\eng{play football, play cards}) ; \eng{go} + activités en \eng{-ing} (\eng{go swimming, go hiking}) ; \eng{do} + sports individuels, arts martiaux et activités de remise en forme (\eng{do karate, do yoga, do aerobics}). Attention à l'article : \eng{play} + \textbf{the} + instrument (\eng{play the guitar}), mais \eng{play} + \textbf{Ø} + sport ou jeu (\eng{play football, play chess}).
\end{corrige}

\begin{corrige}[Exercice 2 — Vrai ou faux ?]

\begin{enumerate}
\item \textbf{False.} Le texte dit \eng{“many young people ... spend their free time playing football”}, et non \emph{all} : d'autres préfèrent les activités d'intérieur.
\item \textbf{True.} \eng{“Others prefer indoor activities: they listen to music, chat with friends on their phones or play video games.”}
\item \textbf{False.} \eng{“At weekends, some families go hiking around Mount Fébé”} — des familles font bien de la randonnée.
\item \textbf{False.} \eng{“very few teenagers read novels for pleasure; most say they have no time for books.”}
\end{enumerate}

\textbf{Point de méthode :} au Bac, la justification par une courte citation entre guillemets est exigée ; recopier la phrase entière du texte n'est pas nécessaire, un extrait suffit. Attention aux quantifieurs (\eng{many, some, few, all}) : c'est souvent là que se joue la réponse.
\end{corrige}

\begin{corrige}[Exercice 3 — Vocabulaire : trouve le mot]

\begin{enumerate}
\item \eng{a hobby} — une activité pratiquée régulièrement pour le plaisir.
\item \eng{leisure} — le temps libre (invariable, sans article : \eng{leisure time}).
\item \eng{entertainment} — ce qui divertit (film, spectacle) ; on accepte aussi \eng{amusement}.
\item \eng{a pastime} — une façon agréable d'occuper son temps (jeu, sport).
\end{enumerate}

\textbf{Nuances :} \eng{hobby} est concret et comptable (\eng{my hobbies are...}) ; \eng{leisure} désigne le temps libre lui-même ; \eng{recreation} est plus soutenu et général ; \eng{amusement} insiste sur le rire et le plaisir (\eng{amusement park} = parc d'attractions).
\end{corrige}

\begin{corrige}[Exercice 4 — Familles de mots : complète]

\begin{enumerate}
\item \eng{Gardening is my grandmother's favourite \textbf{relaxation}.} (nom après l'adjectif possessif \eng{'s})
\item \eng{After the exams, the students felt completely \textbf{relaxed}.} (adjectif en \eng{-ed} : ce que l'on ressent)
\item \eng{We spent a very \textbf{relaxing} evening by the swimming pool.} (adjectif en \eng{-ing} : qualité de la soirée)
\item \eng{Reading is a good form of \textbf{entertainment} for me.} (nom après \eng{form of})
\item \eng{The clown was very \textbf{amusing}; all the children laughed.} (adjectif en \eng{-ing} : le clown fait rire)
\item \eng{My father finds fishing extremely \textbf{boring}.} (la pêche \emph{est} ennuyeuse ; \eng{bored} décrirait mon père)
\end{enumerate}

\textbf{Point de vigilance :} la paire \eng{-ed / -ing} est un classique du Bac : \eng{I am bored} (je m'ennuie) mais \eng{the film is boring} (le film est ennuyeux). Ne jamais écrire \eng{I am boring} pour dire « je m'ennuie » — cela signifie « je suis ennuyeux » !
\end{corrige}

\begin{corrige}[Exercice 5 — Traduction]

\begin{enumerate}
\item \eng{At the weekend, I go swimming and play chess with my brother.}
\item \eng{Many young Cameroonians play football after school.}
\item \eng{My sister does not like video games; she prefers reading novels.} (ou \eng{she prefers to read novels})
\item \eng{To relax, I listen to music or go for a walk with my friends.}
\item \eng{During the holidays, we go hiking in the mountains of the West Region.}
\end{enumerate}

\textbf{Points de vigilance :}
\begin{itemize}
\item « faire de la natation » = \eng{go swimming} (et non \eng{do swimming} ni \eng{make swimming}).
\item « jouer aux échecs / au football » = \eng{play chess / football} sans article.
\item « se promener » = \eng{go for a walk} (faux ami : \eng{to walk} seul = marcher).
\item « après les cours » = \eng{after school} ; « pendant les vacances » = \eng{during the holidays}.
\end{itemize}
\end{corrige}

\begin{corrige}[Exercice 6 — Corrige les faux amis]

\begin{enumerate}
\item Erreur : \eng{assisted to}. Correction : \eng{I \textbf{watched} the football match last Saturday.} (ou \eng{attended}). \eng{Assist to} n'existe pas ; \eng{attend} = assister à (soutenu), mais pour un match on dit plutôt \eng{watch} ou \eng{go to}.
\item Erreur : \eng{make sport}. Correction : \eng{I \textbf{do sport} every day to stay healthy.} (ou \eng{play sports} en anglais américain). \eng{Make} ne s'emploie jamais avec \eng{sport}.
\item Erreur : \eng{repose myself}. Correction : \eng{In my free time, I like to \textbf{relax}.} (ou \eng{to rest}). \eng{Repose} est littéraire et ne se construit pas ainsi ; de plus, \eng{relax} ne prend pas de pronom réfléchi en anglais.
\item Erreur : \eng{passed her weekend}. Correction : \eng{She \textbf{spent} her weekend at her grandmother's village.} Faux ami : « passer du temps » = \eng{spend time} ; \eng{pass the time} signifie « tuer le temps ».
\item Erreur : \eng{amused ourselves}. Correction : \eng{We \textbf{had a lot of fun} at the party.} (ou \eng{enjoyed ourselves very much}). \eng{Amuse oneself} signifie « s'occuper, se divertir » et non « bien s'amuser » au sens de la fête ; \eng{have fun} est l'équivalent naturel.
\item Erreur : \eng{plays with the computer}. Correction : \eng{He \textbf{plays on the computer} all evening.} (ou \eng{plays computer games}). \eng{Play with} suggère un jouet ; pour un appareil, on dit \eng{play on}.
\end{enumerate}
\end{corrige}

\begin{corrige}[Exercice 7 — Texte à trous]

Texte complété :

\eng{In Douala, young people have many different \textbf{hobbies}. After school, boys often \textbf{play} football in the streets, while girls sometimes \textbf{go} dancing or listen to makossa music. For many families, television is the main form of \textbf{entertainment} in the evening. At weekends, some people \textbf{go} swimming at the beach in Limbe. Reading is not a very popular \textbf{pastime} among teenagers, and many say they feel \textbf{bored} when they have no phone. However, doctors say that \textbf{relaxing} activities like walking and gardening are excellent for our health, especially during our \textbf{leisure} time.}

\textbf{Justifications :} \eng{play football} (sport avec ballon) ; \eng{go dancing / go swimming} (activités en \eng{-ing}) ; \eng{feel bored} (sentiment, adjectif en \eng{-ed}) ; \eng{relaxing activities} (qualité, adjectif en \eng{-ing}) ; \eng{leisure time} (collocation fixe = temps libre).
\end{corrige}

\begin{corrige}[Exercice 8 — Appariement : expressions idiomatiques]

\begin{center}
\begin{tabular}{|l|l|}
\hline
\rowcolor{popblueL} \textbf{Expression} & \textbf{Réponse} \\
\hline
1. \eng{to let off steam} & b) se défouler, se détendre \\
\hline
2. \eng{to kill time} & a) perdre son temps (tuer le temps) \\
\hline
3. \eng{it's not my cup of tea} & d) ce n'est pas mon truc \\
\hline
4. \eng{to chill out} & c) se détendre, se relaxer \\
\hline
5. \eng{to have a whale of a time} & e) s'amuser comme un fou \\
\hline
6. \eng{to be a couch potato} & g) être un vrai paresseux devant la télé \\
\hline
7. \eng{to be a sports fanatic} & f) être un passionné de sport \\
\hline
8. \eng{to take up a hobby} & h) se mettre à une activité de loisir \\
\hline
\end{tabular}
\end{center}

\textbf{Exemple d'emploi :} \eng{After the exams, students need to let off steam.} « Après les examens, les élèves ont besoin de se défouler. »
\end{corrige}

\begin{corrige}[Exercice 9 — Compréhension et langue — type Bac]

\textbf{Barème indicatif :} A. 5 pts (1 pt par question) ; B. 5 pts (B1 : 1,5 pt ; B2 : 1,5 pt ; B3 : 2 pts).

\textbf{A. Comprehension — réponses modèles :}

\begin{enumerate}
\item \eng{They really have free time at weekends, because weekdays are taken up by long hours in the classroom and homework in the evening.}
\item \eng{Basketball and handball are mentioned, apart from football.}
\item \eng{Some teenagers prefer listening to music (makossa, afrobeats or gospel) and reading novels or writing poems.} (on accepte aussi : \eng{joining drama clubs or choirs})
\item \eng{Parents worry that their children spend too much time on screens and not enough time on outdoor activities.}
\item Réponse personnelle attendue, par exemple : \eng{In my opinion, recreation is important because it helps the brain to recover after hard work. Sport keeps students healthy, and hobbies like music or reading reduce stress. That is why balanced free time can even improve school results.}
\end{enumerate}

\textbf{B. Language work :}

\begin{enumerate}
\item \textbf{a)} \eng{recreation} (ou \eng{leisure activities}) ; \textbf{b)} \eng{gaining popularity} ; \textbf{c)} \eng{to recover} (dans \eng{helps the brain to recover}).
\item \eng{\textbf{play} basketball} ; \eng{\textbf{go} dancing} ; \eng{\textbf{play} video games}.
\item Reformulations possibles : \eng{“teenagers read fewer and fewer books because they prefer their mobile phones”} ; \eng{“mobile phones are replacing books in teenagers' free time”}.
\end{enumerate}

\textbf{Points de vigilance :} répondre avec des phrases complètes et \emph{avec ses propres mots} (recopier le texte est pénalisé) ; vérifier l'accord sujet-verbe (\eng{parents worry}, \eng{football remains}) ; ne pas confondre \eng{recover} (se remettre, récupérer) et \eng{discover} (découvrir).
\end{corrige}

\begin{corrige}[Exercice 10 — Traduction — type Bac]

\textbf{Barème indicatif :} 10 pts — 2 pts par phrase (lexique 1 pt, grammaire/collocations 1 pt).

\textbf{Traduction modèle :}

\eng{In their free time, the young people of my town do many activities. Boys play football on Saturday afternoons, while girls often prefer dancing or listening to music. Some students go swimming or hiking with their parents. As for me, I am not a sports fanatic: I prefer reading novels or playing chess with my grandfather. These pastimes help me (to) relax after a long week of classes.}

\textbf{Points de vigilance :}
\begin{itemize}
\item « temps libre » = \eng{free time} ou \eng{leisure time} (jamais \eng{free times}).
\item « jouent au football » = \eng{play football} (sans article) ; « jouer aux échecs » = \eng{play chess}.
\item « font de la natation ou de la randonnée » = \eng{go swimming or (go) hiking}.
\item « fanatique du sport » = \eng{a sports fanatic} (expression idiomatique de la leçon).
\item « m'aident à me détendre » = \eng{help me (to) relax} — pas de réflexif en anglais (\eng{relax} sans \eng{myself}).
\item « le samedi après-midi » (habitude) = \eng{on Saturday afternoons} ; présent simple partout car il s'agit d'habitudes.
\end{itemize}
\end{corrige}

\begin{corrige}[Exercice 11 — Composition — type Bac]

\textbf{Barème :} voir grille (lexique 4 pts ; collocations 3 pts ; grammaire 2 pts ; organisation 1 pt).

\textbf{Article modèle (environ 150 mots) :}

\begin{textebox}[How we relax at G.B.H.S. Biyem-Assi]
What do students of our school do when the last bell rings? A quick survey shows that football is still king: every evening, boys play football on the school pitch, and more and more girls are joining them. Basketball and volleyball are also gaining popularity.

Not everyone is a sports fanatic, however. Many students prefer quieter pastimes: they listen to afrobeats, chat with friends or play video games at home. Some go swimming at weekends, while others have taken up new hobbies such as photography or drama. Reading novels, unfortunately, is not many students' cup of tea.

Personally, I love playing chess with my grandfather and going hiking with my family during the holidays. After a long week of classes, these activities help me to let off steam.

Recreation is not a waste of time: it rests the brain, keeps the body fit and makes us better students. So, dear readers, find your hobby — and enjoy it!
\end{textebox}

\textbf{Points de vigilance pour l'élève :}
\begin{itemize}
\item Employer au moins trois collocations correctes : \eng{play football / chess / video games}, \eng{go swimming / hiking}, \eng{do sport}.
\item Insérer une expression idiomatique : \eng{let off steam}, \eng{not my cup of tea}, \eng{a sports fanatic}, \eng{take up a hobby}.
\item Utiliser le présent simple (habitudes) et varier les connecteurs (\eng{however, while, personally, unfortunately}).
\item Structurer l'article : titre, introduction accrocheuse (question), développement, opinion personnelle, conclusion.
\item Éviter les calques : « faire du sport » = \eng{do sport}, « passer le temps » = \eng{spend time}, « s'amuser » = \eng{have fun / enjoy oneself}.
\end{itemize}
\end{corrige}

\newpage

\coursec{Fiche bilan}

\coursec{Lesson 9 — Vocabulary: Words and expressions related to recreational activities}

\begin{popfigure}
\centering
\begin{tikzpicture}[
    mindmap,
    grow cyclic,
    every node/.style={concept, execute at begin node=\hskip0pt},
    concept color=popblue!60,
    level 1/.append style={level distance=4.5cm, sibling angle=360/7},
    level 2/.append style={level distance=3cm, sibling angle=360/4},
    text=white,
    font=\small\bfseries
]
\node[concept, scale=1.2, align=center]{Recreational\\ Activities}
    child[concept color=popgreen] { node[concept, align=center]{Sports \&\\ Games} }
    child[concept color=poporange] { node[concept, align=center]{Arts \&\\ Culture} }
    child[concept color=poppurple] { node[concept, align=center]{Outdoor \&\\ Nature} }
    child[concept color=poppink] { node[concept, align=center]{Social \&\\ Digital} }
    child[concept color=popdark] { node[concept, align=center]{Verbs:\\ Play / Do / Go} }
    child[concept color=popgold] { node[concept, align=center]{Idioms \&\\ Expressions} }
    child[concept color=popmuted] { node[concept, align=center]{Word Families\\ \& Pronunciation} };
\end{tikzpicture}
\legende{Carte mentale : Vocabulaire des activités récréatives (Tle A)}
\end{popfigure}

\courssub{Champs lexicaux : les grandes familles de loisirs}

\begin{definition}[Lexique clé — English / Français]
\begin{center}
\begin{tabularx}{\linewidth}{|>{\bfseries}X|X|X|}\hline
\rowcolor{popblueL} \textbf{Domaine} & \textbf{English} & \textbf{Français} \\ \hline
\cellcolor{popgreenL}\textbf{Sports \& Games} & athletics, badminton, basketball, boxing, chess, cycling, football, gymnastics, handball, jogging, martial arts, rugby, swimming, table tennis, tennis, volleyball, weightlifting, yoga & athlétisme, badminton, basket, boxe, échecs, cyclisme, football, gymnastique, handball, jogging, arts martiaux, rugby, natation, tennis de table, tennis, volley-ball, haltérophilie, yoga \\ \hline
\cellcolor{poporangeL}\textbf{Arts \& Culture} & acting, ballet, cinema, concert, drawing, exhibition, festival, painting, photography, playing an instrument, pottery, reading, sculpture, theatre & jeu d'acteur, ballet, cinéma, concert, dessin, exposition, festival, peinture, photographie, jouer d'un instrument, poterie, lecture, sculpture, théâtre \\ \hline
\cellcolor{poppurpleL}\textbf{Outdoor \& Nature} & birdwatching, camping, fishing, gardening, hiking, trekking, hunting, picnic, safari, sightseeing, sunbathing, swimming (outdoor) & ornithologie, camping, pêche, jardinage, randonnée, trek, chasse, pique-nique, safari, visite touristique, bain de soleil, baignade (extérieur) \\ \hline
\cellcolor{poppinkL}\textbf{Social \& Digital} & board games, bowling, chatting (online), clubbing, dancing, gaming, karaoke, night out, party, scrolling (social media), streaming, video calls & jeux de société, bowling, discuter (en ligne), boîte de nuit, danser, jeux vidéo, karaoké, sortie nocturne, fête, défiler (réseaux), streaming, appels vidéo \\ \hline
\end{tabularx}
\end{center}
\medskip
\textbf{Note culturelle :} Au Cameroun, on ajoute : \eng{bendskin} (danse/musique urbaine), \eng{buvette} (bar populaire), \eng{playground} (terrain de quartier), \eng{canoe racing} / \eng{pirogue racing} (courses de pirogues), \eng{mountain racing} (course de côte, ex. Mont Cameroun).
\end{definition}

\courssub{Collocations : play, do, go et autres verbes associés}

\begin{propriete}[Règles de collocations — \cle{Play / Do / Go} (à connaître par cœur)]
\begin{center}
\begin{tabularx}{\linewidth}{|>{\bfseries}X|>{\bfseries}X|X|}\hline
\rowcolor{popblueL} \textbf{Verbe} & \textbf{Règle générale} & \textbf{Exemples types (EN $\rightarrow$ FR)} \\ \hline
\cellcolor{popgreenL}\textbf{PLAY} & Jeux de balle, sports d'équipe, jeux de compétition/règles, instruments de musique. & \eng{play football, play chess, play the piano, play cards, play volleyball} (jouer au foot, aux échecs, du piano, aux cartes, au volley) \\ \hline
\cellcolor{poporangeL}\textbf{DO} & Activités individuelles, exercices corporels, arts martiaux, loisirs non-compétitifs, tâches/jeux d'esprit. & \eng{do yoga, do gymnastics, do karate, do exercise, do a puzzle, do crosswords, do judo} (faire du yoga, de la gym, du karaté, de l'exercice, un puzzle, des mots croisés, du judo) \\ \hline
\cellcolor{poppurpleL}\textbf{GO} & Activités en \cle{-ing} impliquant un déplacement vers un lieu de pratique (souvent outdoor). & \eng{go swimming, go hiking, go fishing, go camping, go jogging, go dancing, go shopping, go skiing} (aller nager, randonner, pêcher, camper, faire du jogging, danser, faire les courses, skier) \\ \hline
\end{tabularx}
\end{center}
\medskip
\textbf{Autres verbes utiles :} \eng{attend} (assister à un événement), \eng{watch} / \eng{see} (regarder un match/film), \eng{practise} (s'entraîner à — verbe UK), \eng{enjoy} (apprécier), \eng{be keen on} / \eng{be fond of} / \eng{be into} (+ \cle{-ing}), \eng{hang out} (traîner avec des amis), \eng{chill out} / \eng{unwind} (se détendre, décompresser).
\end{propriete}

\courssub{Expressions idiomatiques et langage courant des loisirs}

\begin{definition}[Expressions idiomatiques \& langage courant]
\begin{center}
\begin{tabularx}{\linewidth}{|>{\bfseries}X|X|X|}\hline
\rowcolor{popblueL} \textbf{Expression} & \textbf{Sens} & \textbf{Exemple / Contexte} \\ \hline
\eng{to kill time} & Occuper le temps, s'occuper en attendant & \eng{We played cards to kill time while waiting for the bus.} (On a joué aux cartes pour passer le temps en attendant le bus.) \\ \hline
\eng{to recharge one's batteries} & Se ressourcer, se reposer profondément & \eng{I need a weekend in the village to recharge my batteries.} \\ \hline
\eng{to let one's hair down} & Se détendre, faire la fête sans contrainte & \eng{After the Bac, the students let their hair down.} \\ \hline
\eng{a night out} & Une sortie (soirée) entre amis & \eng{Let's have a night out in Douala this Saturday.} \\ \hline
\eng{to be a couch potato} & Être un fainéant qui passe son temps devant la télé/écrans & \eng{Don't be a couch potato; go outside and play!} \\ \hline
\eng{to hang out} & Passer du temps ensemble (sans but précis) & \eng{Where do you usually hang out after school?} \\ \hline
\eng{to catch a movie / a show} & Aller voir un film / un spectacle (familier) & \eng{Do you want to catch a movie at the cinema tonight?} \\ \hline
\eng{to take up a hobby} & Se mettre à un loisir, commencer une activité & \eng{She took up photography last year.} \\ \hline
\eng{to wind down} & Se relaxer, ralentir le rythme (après le travail/études) & \eng{I read a book to wind down before bed.} \\ \hline
\end{tabularx}
\end{center}
\end{definition}

\courssub{Faux amis et pièges des francophones}

\begin{attention}[Faux amis \& Pièges fréquents (Francophones)]
\begin{itemize}
    \item \textbf{Fun vs Funny :} \cle{Fun} = amusant, divertissant (nom ou adjectif) \eng{It was great fun. / A fun party.} \cle{Funny} = drôle (qui fait rire) ou bizarre \eng{A funny joke. / A funny smell.}
    \item \textbf{Entertainment :} Orthographe \cle{Entertainment} (pas de 'a' après 't' : *entertainement).
    \item \textbf{Practice (n) / Practise (v) :} UK \eng{practice} (nom), \eng{practise} (verbe). US : \eng{practice} pour les deux. Au Cameroun (norme UK) : \eng{I practise the piano every day. Piano practice is at 5 pm.}
    \item \textbf{Play vs Play at :} \eng{Play football, play the guitar} (pas \eng{play at football}). \eng{Play at} = faire semblant \eng{He's only playing at being brave.}
    \item \textbf{Go to the cinema / Go to the movies :} \eng{Go to the cinema} (UK), \eng{Go to the movies} (US). Jamais \eng{go at the cinema}.
    \item \textbf{Do sport vs Play a sport :} \eng{Do sport} (général, UK) / \eng{Play a sport} (spécifique). \eng{I do sport every day. I play tennis.}
    \item \textbf{Hobby + -ing :} \eng{My hobby is reading} (gérondif naturel). \eng{My hobby is to read} (possible mais moins idiomatique).
    \item \textbf{Leisure :} Prononciation « léjœur » (UK) / « li-jour » (US). Pas « lei-zyur ». Orthographe : \cle{leisure} (ie).
    \item \textbf{Competition :} Orthographe \cle{competition} (pas *compétition). Adjectif : \cle{competitive} (pas *compétitif).
\end{itemize}
\end{attention}

\courssub{Familles de mots, prononciation et orthographe}

\begin{propriete}[Familles de mots, Prononciation, Orthographe]
\begin{center}
\begin{tabularx}{\linewidth}{|>{\bfseries}l|>{\bfseries}l|>{\bfseries}l|X|}\hline
\rowcolor{popblueL} \textbf{Verbe / Base} & \textbf{Nom (Action/Personne)} & \textbf{Adjectif} & \textbf{Prononciation (approx. FR) / Note} \\ \hline
recreate & recreation, recreational & recreational & « ré-kri-é-ï-chone » / stress sur \textbf{é} (3e syll.) \\ \hline
entertain & entertainment, entertainer & entertaining & « èn-tèr-tèin-mente » / pas de 'a' après 't' \\ \hline
compete & competition, competitor & competitive & « com-pé-ti-chone » / adj \eng{competitive} \\ \hline
relax & relaxation & relaxed / relaxing & \cle{Relaxed} = soulagé, détendu (état) ; \cle{Relaxing} = qui détend (cause) \\ \hline
amuse & amusement & amused / amusing & \cle{Amused} = content, diverti ; \cle{Amusing} = drôle, divertissant \\ \hline
exercise & exercise & exercising & « èk-sèr-saïze » (verbe \& nom, son /z/ final) \\ \hline
photograph & photography, photographer & photographic & Stress shift : \textbf{PHO}-to-graph $\rightarrow$ pho-\textbf{TOG}-ra-phy $\rightarrow$ pho-\textbf{TOG}-ra-pher $\rightarrow$ pho-to-\textbf{GRAPH}-ic \\ \hline
\end{tabularx}
\end{center}
\medskip
\textbf{Rappel \cle{-ed / -ing} (adjectifs participiaux) :} \eng{bored/boring, interested/interesting, excited/exciting, tired/tiring, fascinated/fascinating}. Règle : \cle{-ed} = sentiment ressenti par une personne ; \cle{-ing} = caractéristique de la chose/personne qui provoque le sentiment.
\end{propriete}

\courssub{Emploi en contexte : dialogues et court texte}

\begin{textebox}[Dialogue : Planning a weekend in Yaoundé]
\textbf{Context:} Two friends, \eng{Alice} and \eng{Bruno}, discuss their plans for Saturday.

\eng{Alice:} Hey Bruno, what are you up to this weekend? Any plans?
\eng{Bruno:} Not really. I'll probably just \textbf{hang out} at home, maybe \textbf{catch a movie} on Netflix. I'm exhausted after the mock exams.
\eng{Alice:} Come on! Don't be a \textbf{couch potato}. Let's \textbf{go hiking} on Mount Fébé. It's great to \textbf{recharge your batteries}.
\eng{Bruno:} Hiking? I haven't done that in ages. I usually \textbf{do yoga} on Saturdays to \textbf{wind down}.
\eng{Alice:} You can do yoga anytime. This is a \textbf{night out}... well, a day out! We could \textbf{have a picnic} at the top. \textbf{Kill two birds with one stone}: exercise and fun.
\eng{Bruno:} \textbf{Kill time} on the mountain? Okay, you talked me into it. What time?
\eng{Alice:} 7 am sharp. Wear good trainers. And bring your camera if you're \textbf{into photography}.
\eng{Bruno:} Deal. It'll be \textbf{fun}. Not just \textbf{funny} selfies!
\eng{Alice:} Exactly. See you Saturday!
\end{textebox}

\begin{textebox}[Texte court : Blog post — My ideal weekend]
\textbf{Title:} \eng{Weekends are for Unwinding}

\eng{Weekends are sacred for me. After a hectic week of classes at the lycée, I need to \textbf{unwind} properly. My perfect Saturday starts with a lie-in, followed by a hearty breakfast of \eng{beans and plantains}. Then, I \textbf{go jogging} around the \eng{playground} in my neighbourhood — it's my way to \textbf{stay fit} and clear my head.}

\eng{In the afternoon, it's time for culture. I'm \textbf{keen on photography}, so I often \textbf{go sightseeing} in the city centre or visit an \textbf{exhibition} at the National Museum. Capturing the light on the \eng{Reunification Monument} is \textbf{relaxing}. Later, I meet my crew at a \eng{buvette} for a \textbf{night out}. We listen to \textbf{bendskin}, \textbf{chill out}, and \textbf{catch up} on gossip.}

\eng{Sunday is slower. I \textbf{do some reading} (African literature mostly) or \textbf{play chess} with my little brother. Sometimes I \textbf{practise the guitar} — I'm learning \eng{Makossa} rhythms. Before bed, I \textbf{stream a series} to \textbf{wind down}. That's how I \textbf{recharge my batteries} for Monday. What about you?}
\end{textebox}

\begin{experience}[Activité orale : Role-play — "Weekend Planner"]
\textbf{Consigne :} En binôme. L'élève A propose une activité pour samedi (choisie dans les 4 champs lexicaux). L'élève B accepte ou refuse poliment en donnant une raison (utiliser \eng{I'd rather...}, \eng{I'm not really into...}, \eng{I have to...}) et propose une alternative. Utiliser au moins 5 collocations \cle{PLAY/DO/GO}, 2 idiomes et 2 adjectifs \cle{-ed/-ing}. Durée : 3 min. Exemple de début :
\eng{A: "Hey, do you want to go hiking on Mount Cameroon this Saturday?"}
\eng{B: "I'd love to, but I'm exhausted. I'd rather go swimming at the pool. It's more relaxing."}
\end{experience}

\courssub{Méthodes express}

\begin{methode}[Apprendre le vocabulaire des loisirs efficacement]
\begin{enumerate}
    \item \textbf{Classer par verbe opérateur :} Faire 3 colonnes \cle{PLAY / DO / GO} et y ranger chaque nouvelle activité apprise.
    \item \textbf{Contextualiser personnellement :} Écrire 1 phrase vraie par mot (ex. \eng{I go hiking in the mountains near Buea.}).
    \item \textbf{Collocations figées :} Apprendre les blocs \eng{do yoga}, \eng{play the guitar}, \eng{go shopping}, \eng{catch a movie} comme des mots uniques.
    \item \textbf{Prononciation active :} Utiliser la synthèse vocale (Google Translate, DeepL, Forvo) pour répéter les mots difficiles : \eng{photography} (« fé-tog-ra-fi »), \eng{recreation} (« ré-kri-é-ï-chone »), \eng{leisure} (« léjœur »).
    \item \textbf{Réviser espacé :} Revoir les listes J+1, J+7, J+30 (appli Anki / Quizlet / cartes physiques).
\end{enumerate}
\end{methode}

\begin{methode}[Réemployer à l'oral (Speaking) et à l'écrit (Writing)]
\begin{enumerate}
    \item \textbf{Speaking (Part 1 / Interview) :} Préparer des réponses types : \eng{In my free time, I usually...}, \eng{I'm keen on...}, \eng{I unwind by + -ing...}, \eng{I'm into...}. Utiliser des \cle{linking words} : \eng{besides, actually, to be honest, as a matter of fact}.
    \item \textbf{Writing (Article / Blog / Email informel) :} Structure : 1. \textbf{Hook} (\eng{Weekends are perfect for...}), 2. \textbf{Description} activité + verbes variés (\cle{PLAY/DO/GO}), 3. \textbf{Sentiments} (\eng{It helps me recharge / I find it relaxing}), 4. \textbf{Recommandation} (\eng{You should try... / I highly recommend...}).
    \item \textbf{Éviter la répétition de "like" :} Varier : \eng{enjoy, love, be fond of, be keen on, be into, fancy, appreciate, be a fan of}.
    \item \textbf{Position de l'adverbe de fréquence :} \eng{I often go swimming} (avant verbe principal) / \eng{I am always keen on...} (après \cle{be}). Jamais \eng{*I go often swimming}.
\end{enumerate}
\end{methode}

\begin{aretenir}[Checklist avant le Bac]
\begin{itemize}
    \item ☐ Je distingue \cle{PLAY} (ballon/équipe/instrument), \cle{DO} (individuel/exercice), \cle{GO} (+ \cle{-ing} déplacement).
    \item ☐ Je sais orthographier \cle{Entertainment}, \cle{Leisure}, \cle{Recreation}, \cle{Competition}, \cle{Photography} sans faute.
    \item ☐ Je n'écris jamais \eng{play at football} ni \eng{go at the cinema}.
    \item ☐ Je différencie \cle{Fun} (nom/adj : amusant, divertissant) et \cle{Funny} (adj : qui fait rire / bizarre).
    \item ☐ Je maîtrise 5 idiomes minimum : \eng{kill time, recharge batteries, let hair down, night out, hang out, catch a movie, take up a hobby, wind down}.
    \item ☐ Je connais les paires \cle{-ed / -ing} : \eng{relaxed/relaxing}, \eng{amused/amusing}, \eng{bored/boring}, \eng{interested/interesting}, \cle{excited/exciting}, \eng{tired/tiring}.
    \item ☐ Je place correctement l'adverbe de fréquence : \eng{I often go swimming} (pas \eng{*I go often swimming}).
    \item ☐ Je sais décrire une activité : \textbf{Verbe + Lieu + Fréquence + Pourquoi} (ex. \eng{I play basketball at the college court twice a week to stay fit.}).
    \item ☐ Je connais le vocabulaire spécifique camerounais : \eng{bendskin}, \eng{buvette}, \eng{playground}, \eng{canoe racing / pirogue racing}, \eng{mountain racing}.
    \item ☐ Je sais transformer \eng{photograph} $\rightarrow$ \eng{photography} $\rightarrow$ \eng{photographer} $\rightarrow$ \eng{photographic} (stress shift).
    \item ☐ Je peux écrire un paragraphe de 80--100 mots sur mes loisirs en utilisant 10 mots du cours.
    \item ☐ Je comprends un dialogue audio (Listening) sur la planification d'un week-end (vocabulaire clé : \eng{book tickets, queue, crowd, venue, schedule, entrance fee, refund}).
\end{itemize}
\end{aretenir}

\findecours

\end{document}
